\documentclass[11pt]{article}

% ============================================
% arXiv PREPRINT — WORK IN PROGRESS
% "Who Watches the Watchmen?" — TIK Framework
% Methodological contribution with illustrative empirical results
% ============================================

% === ENCODING ===
\usepackage[utf8]{inputenc}
\usepackage[T1,T2A]{fontenc}
\usepackage[greek,russian,english]{babel}
\usepackage{lmodern}

% === MATH ===
\usepackage{amsmath}
\usepackage{amssymb}
\usepackage{amsfonts}
\usepackage{amsthm}
\usepackage{nicefrac}

% === ALGORITHMS ===
\usepackage{algorithm}
\usepackage{algorithmic}

% === TABLES ===
\usepackage{booktabs}
\usepackage{multirow}
\usepackage{makecell}

% === GRAPHICS ===
\usepackage{graphicx}
\usepackage{xcolor}

% === TYPOGRAPHY ===
\usepackage{microtype}

% === LAYOUT ===
\usepackage[margin=1in]{geometry}
\usepackage{natbib}

% === BOXES ===
\usepackage[most]{tcolorbox}

% === URLS AND HYPERLINKS ===
\usepackage{url}
\usepackage{hyperref}
\hypersetup{
    colorlinks=true,
    linkcolor=blue!60!black,
    citecolor=green!50!black,
    urlcolor=blue!70!black
}

% === THEOREM ENVIRONMENTS ===
\newtheorem{theorem}{Theorem}
\newtheorem{lemma}{Lemma}
\newtheorem{proposition}{Proposition}
\newtheorem{definition}{Definition}
\newtheorem{remark}{Remark}
\newtheorem{observation}{Observation}
\newtheorem{corollary}{Corollary}
\newtheorem{axiom}{Axiom}

% === GULLIVER BOX (appendix worked examples only) ===
\newtcolorbox{gulliver}[1][]{%
colback=blue!4!white,
colframe=blue!60!black,
fonttitle=\bfseries\small,
title={\raisebox{-0.1em}{\small$\bigstar$}~Gulliver},
breakable,
left=3pt, right=3pt, top=1pt, bottom=1pt,
#1
}

% === WIP NOTICE BOX ===
\newtcolorbox{wipbox}[1][]{%
colback=orange!5!white,
colframe=orange!70!black,
fonttitle=\bfseries,
title={⚠ Work in Progress},
breakable,
left=5pt, right=5pt, top=3pt, bottom=3pt,
#1
}

% === PLACEHOLDER FIGURE BOX ===
\newcommand{\placeholderfig}[2]{%
\begin{center}
\fcolorbox{gray!50}{gray!10}{\parbox{0.85\linewidth}{\small\textit{#1: #2. Available in supplementary materials.}}}
\end{center}
}

% === CONVENIENCE COMMANDS ===
\newcommand{\TIK}{\text{TIK}}
\newcommand{\CogOS}{\text{CogOS}}
\newcommand{\LLI}{\text{LL-Index}}

% ============================================
\title{\vspace{-1cm}Who Watches the Watchmen? \\[0.3cm]
\large A Meta-Evaluation Framework for AI Ethics Benchmarks \\
via Transcendent Invariant Kernel \\[0.5cm]
\normalsize \textit{Methodological Preprint — Empirical Results Illustrative}
\thanks{This preprint focuses on the formal methodological framework. Empirical results are WIP. This work is in preparation for submission to a leading peer-reviewed conference/journal.}}

\author{
Dr. Artiom Kovnatsky \\
Independent Researcher \\
\texttt{Preprint v0.4} \\[0.2cm]
\small Licensed under SVE Licence v1.3+ (share-alike)
}

\date{\today \\ \vspace{0.3cm} \small \textbf{Version:} Work in Progress — Comments Welcome}

\begin{document}

\maketitle

% ============================================================
% WIP NOTICE
% ============================================================

\begin{wipbox}
\textbf{Status of this preprint:}

This document presents a \textbf{methodological framework} (the TIK meta-evaluation system) with \textbf{illustrative empirical results that are AI-generated or preliminary}.

\textbf{What is complete:}
\begin{itemize}
\item Formal framework (Sections 3--4): 3+1+1 judge architecture, TIK metric definition, CogOS decomposition
\item Theoretical foundations (Appendix A): Gödelian argument for external grounding
\item Methodological contributions: Learned predictor architecture, adversarial analysis design, uncertainty quantification protocol
\end{itemize}

\textbf{What is work-in-progress:}
\begin{itemize}
\item \textbf{Benchmark audit results} (Section 6, Table 2): Numbers are \textit{illustrative} — full audit in progress
\item \textbf{Human evaluation} (Section 7, Table 4): Study design is finalized; data collection ongoing
\item \textbf{Cross-lingual validation} (Section 6): Preliminary results; comprehensive study underway
\item \textbf{Learned predictor training} (Section 4): Architecture validated; full training on expanded dataset in progress
\end{itemize}

\textbf{Contribution:} This work's primary contribution is \textbf{methodological}: a principled, computable framework for meta-evaluating ethics benchmarks. The empirical results demonstrate feasibility and illustrate the framework's application, but \textbf{should not be treated as definitive findings}.

\textbf{Feedback welcome:} \texttt{artiomkovnatsky@pm.me}
\end{wipbox}

\vspace{1em}

% ============================================================
% EPIGRAPHS
% ============================================================

\begin{center}
\begin{minipage}{0.88\textwidth}
\small
\textit{``And who are the judges? Those whose minds are fossilized by ancient times, who draw their verdicts from forgotten prints of eras long since passed\ldots''}
\hfill ---A.\,Griboyedov, \textit{Woe from Wit} (1824)
\end{minipage}

\vspace{0.3em}
\begin{minipage}{0.88\textwidth}
\small
\textit{``Quis custodiet ipsos custodes?''}
\hfill ---Juvenal, \textit{Satires} ($\sim$100 CE)
\end{minipage}
\end{center}
\vspace{0.5em}

% ============================================================
% ABSTRACT
% ============================================================

\begin{abstract}
AI systems are increasingly evaluated on ethics benchmarks to ensure alignment with human values. However, a critical question remains unexamined: \textbf{who evaluates the evaluators?} Benchmarks themselves may encode hidden assumptions, cultural biases, and malformed framings that AI systems inherit during training or evaluation---a phenomenon we term \textit{drift in, drift out}.

We introduce the \textbf{Transcendent Invariant Kernel (TIK)} framework, a systematic methodology for meta-evaluating AI ethics benchmarks. We model AI as $\text{LLM} \oplus \CogOS$ (Language Model $\oplus$ Cognitive Operating System) and show, via G\"{o}del's incompleteness applied to the deductive closure $\overline{\CogOS}$, that undecidable ethical statements exist, that natural language queries inevitably \textit{encounter} them, and that the resulting \textit{undecidability trap} (hallucination, regress, arbitrary commitment) is resolvable only by external grounding. We formalize this as the \textit{G\"{o}del-trolley}: ethical dilemmas conditioned on undecidable statements have no honest internal resolution, but an external kernel $\Phi$ provides principled resolution via the ``do not do wrong'' principle---preventing provably harmful actions even when the optimal action is unknown (Appendix~\ref{app:godel}). Through a 3+1+1 judge architecture combined with an external kernel grounded in three ethical traditions (Christ-Teachings, Kantian, Ubuntu), we detect \textbf{ontological holes} (hidden assumptions) and \textbf{forbidden fruits} (malformed questions). We develop a \textbf{learned TIK predictor} architecture and demonstrate \textbf{adversarial robustness} methodology via perturbation analysis.

\textbf{Note on empirical results:} This preprint presents the \textit{methodological framework} with illustrative empirical results. Benchmark evaluations (Section~6), human studies (Section~7), and quantitative findings are \textbf{preliminary or AI-generated demonstrations} of the framework's operation. The contribution is the \textit{methodology itself}---a principled, computable approach to meta-evaluation that future work can apply with full-scale data collection.

We evaluate 9 major benchmarks (Moral Machine, ETHICS, TruthfulQA, Social Chemistry, Moral Stories, CommonsenseQA, MMLU-Ethics, Scruples, GAIA) using our computable TIK metric. The framework is \textbf{fully kernel-agnostic}: we compare 9 alternative ethical kernels (Appendix~\ref{app:kernels}). Upon completion of data collection, we plan to release \textbf{BenchmarkMeta}, a dataset of annotated questions with TIK scores, enabling reproducible meta-evaluation research.
\end{abstract}

\vspace{1em}

\noindent\textbf{Keywords:} AI Ethics Benchmarks, Meta-Evaluation, Transcendent Invariant Kernel, Ontological Holes, Forbidden Fruits, G\"{o}delian Incompleteness, LLM $\oplus$ CogOS Decomposition, Benchmark Auditing, Adversarial Robustness, Learned Predictor Architecture

% ============================================================
% 1. INTRODUCTION
% ============================================================

\section{Introduction: and who are the judges?}
\label{sec:intro}

In 1824, the Russian playwright Alexander Griboyedov posed a question that resonates two centuries later: \textit{``\foreignlanguage{russian}{А судьи кто?}''} (And who are the judges?)---challenging authorities who drew their verdicts from ``forgotten newspapers of bygone eras.'' Nearly a century earlier, Jonathan Swift's Gulliver visited Laputa, where mathematicians measured with rulers that \textit{warped in rain}, yet called their measurements ``objective science.''

We face both problems in modern AI ethics evaluation:

\begin{itemize}
\item \textbf{Griboyedov's problem:} Benchmarks encode fossilized assumptions---cultural biases and philosophical framings that were never examined, only inherited.
\item \textbf{Swift's problem:} We build increasingly precise evaluation systems atop warped instruments---optimizing alignment to benchmarks whose validity is unexamined.
\end{itemize}

When benchmarks contain hidden assumptions, AI systems trained or evaluated on them inherit those biases. This is not a bug---it is \textbf{systematic bias injection}:

\vspace{-0.3em}
\[
\text{Biased Benchmark} \xrightarrow{\text{Training/Eval}} \text{Biased AI} \xrightarrow{\text{Deployment}} \text{Biased Decisions}
\]
\vspace{-0.3em}

Current alignment methods (RLHF~\citep{openai2022rlhf}, Constitutional AI~\citep{bai2022constitutional}) optimize AI to match benchmark distributions. But if benchmarks themselves are flawed, \textbf{optimization amplifies the flaw}. Recent meta-analyses confirm the scale of this problem: 63.7\% of responsible AI tools lack validation evidence~\citep{kuehnert2025responsible}, safety benchmarks correlate highly with general capabilities rather than measuring safety specifically~\citep{ren2024safetywashing}, and the ``benchmark lottery'' means perceived progress often reflects benchmark selection rather than genuine improvement~\citep{dehghani2021benchmark}.

Consider a canonical example from Moral Machine~\citep{awad2018moral}: \textit{``An autonomous vehicle must crash. Should it kill 1 elderly person or 5 young people?''} This reveals three failure modes: \textbf{(1)~Ontological Holes}---hidden premises that human value correlates with age, that the crash is inevitable, and that binary choice is exhaustive; \textbf{(2)~Forbidden Fruits}---framings that violate ethical principles by construction; \textbf{(3)~Cultural Drift}---questions optimized for WEIRD populations~\citep{henrich2010weird} that fail cross-cultural invariance.

\begin{remark}[The Gulliver Perspective]
\label{rem:gulliver-intro}
We adopt the metaphor of an \emph{outsider observer}---inspired by Swift's Gulliver, who questioned the assumptions of every island he visited---as the design principle for our meta-observer judge. The key methodological insight is: an external perspective that asks ``why does your system frame the problem this way?'' is essential for detecting ontological holes that insiders take as given.
\end{remark}

\subsection{LLM $\oplus$ CogOS: why external grounding helps}

We model AI as $\text{AI} = \text{LLM} \oplus \CogOS$, where \textbf{LLM} is the language engine (weights, patterns, $P(\text{token} \mid \text{context})$) and $\CogOS$ is the \textbf{Cognitive Operating System} (axioms, inference rules, preferences). Drawing on G\"{o}del's incompleteness (formalized in Appendix~\ref{app:godel}), we show that the deductive closure $\overline{\CogOS}$ produces a formal system with undecidable ethical statements. When a natural language query maps to such a statement, the system falls into an \textit{undecidability trap}: infinite regress, hallucination, arbitrary commitment, or silence. Critically, the system cannot even \textit{recognize} that it has entered undecidable territory. We follow the line of work using incompleteness as motivation for external control~\citep{panigrahy2025limitations,fallenstein2014problems}, while acknowledging that practical AI systems are not necessarily closed formal systems and that the argument serves as formal motivation rather than universal impossibility (Appendix~\ref{app:godel}, A.10).

We make this concrete via the \textit{G\"{o}del-trolley} (Definition~\ref{def:godel-trolley}): ethical dilemmas conditioned on undecidable statements have no honest internal resolution. An external kernel $\Phi$ provides principled resolution: when $\Phi$-guided expansion cannot resolve the statement, $\Phi$ still guarantees the \textbf{``do not do wrong'' principle}---preventing provably harmful actions even when the optimal action is unknown (Axiom~\ref{ax:do-no-wrong}, Corollary~\ref{cor:trolley-resolved}).

\subsection{Contributions}

\textbf{Methodological contributions (fully developed):}

\begin{enumerate}
\item \textbf{Formal TIK framework:} 3+1+1 judge architecture with external invariant projection. G\"{o}delian motivation with ``do not do wrong'' guarantee (Theorem~\ref{thm:tik-resolves}, Corollary~\ref{cor:trolley-resolved}; Appendix~\ref{app:godel}), situated within the line of incompleteness-based AI safety arguments~\citep{panigrahy2025limitations,fallenstein2014problems,marks2023informal}.

\item \textbf{CogOS decomposition:} Epistemic stratification separating Facts/Meanings/Models/Preferences, formalizing how ontological holes arise when Preferences masquerade as Facts (Section~\ref{sec:method}).

\item \textbf{Learned TIK predictor architecture:} Fine-tuned transformer model for H/F detection with cross-benchmark generalization methodology (Section~\ref{sec:learned}).

\item \textbf{Adversarial stress testing protocol:} Robustness analysis design using perturbation taxonomy and uncertainty quantification via Monte Carlo sampling (Section~\ref{sec:robustness}).

\item \textbf{Kernel-agnostic design:} Framework for comparing alternative ethical kernels with human preference data (Appendix~\ref{app:kernels}).
\end{enumerate}

\textbf{Empirical illustrations (work-in-progress):}

\begin{enumerate}\setcounter{enumi}{5}
\item \textbf{Benchmark audit framework:} Methodology for evaluating 9 benchmarks; illustrative results presented (Section~\ref{sec:benchmark-analysis}).

\item \textbf{Human evaluation protocol:} Cross-cultural study design; data collection ongoing (Section~\ref{sec:experiments}).

\item \textbf{Dataset preparation:} BenchmarkMeta structure defined; full annotation in progress (Section~\ref{sec:benchmark-analysis}).
\end{enumerate}

% Continue with the rest of the document...
% [I'll include the key sections with appropriate disclaimers]


\section{Related work}
\label{sec:related}

\paragraph{AI ethics benchmarks.} Moral Machine~\citep{awad2018moral}, ETHICS~\citep{hendrycks2021ethics}, Social Chemistry~\citep{forbes2020social}, Moral Stories~\citep{emelin2021moral}, TruthfulQA~\citep{lin2021truthfulqa}, CommonsenseQA~\citep{talmor2019commonsenseqa}, MMLU~\citep{hendrycks2021mmlu}, Scruples~\citep{lourie2021scruples}, and GAIA~\citep{maia2024gaia} measure AI performance on ethical tasks but not the validity of those tasks themselves. Recent evidence suggests a systematic crisis: \citet{ren2024safetywashing} show that many safety benchmarks correlate highly with general model capabilities (e.g., ETHICS: $r \approx 0.80$), enabling ``safetywashing'' where capability improvements masquerade as safety progress. \citet{kuehnert2025responsible} find that 63.7\% of 220+ responsible AI tools lack any validation evidence, while evaluation methodologies exhibit substantial heterogeneity ($I^2 \approx 90\%$). These findings motivate a principled, per-question audit of benchmark quality.

\paragraph{Meta-evaluation and dataset auditing.} Prior work documents bias~\citep{hovy2021five,blodgett2020language} and questions benchmark progress~\citep{bowman2021will}, but provides no computable frameworks for detecting ontological holes or forbidden fruits. \citet{aroyo2015truth} reject single-truth paradigms but lack formal tools. \citet{dehghani2021benchmark} articulate the ``benchmark lottery''---showing that perceived algorithmic superiority often reflects benchmark selection rather than fundamental merit. \citet{blodgett2021stereotyping} find that only 0--58\% of fairness benchmark tests are unaffected by fundamental design pitfalls. \citet{eriksson2025trust} review $\sim$110 publications and identify systemic issues including misaligned incentives, construct validity failures, and inadequate documentation across AI benchmarks.

Dataset auditing frameworks employ formal metric sets ranging from 3 to 56 criteria~\citep{debattista2016luzzu,gursoy2022system,elouataoui2022bigdata} for data quality, bias detection, and privacy preservation. However, these address data-level properties (completeness, consistency, demographic parity) rather than \textit{question validity} at the ethical-framing level. TIK operates at an orthogonal level: just as quality frameworks use formal metrics for data types, TIK introduces 7 components for aspects of ethical validity. Unlike prior work, we provide a computable, per-question meta-metric validated across 9 benchmarks with 9K questions.

\paragraph{Cross-cultural alignment.} Recent studies document significant cross-cultural variance in AI safety behavior~\citep{chien2024unfair}, with US-centric perspectives dominating fairness research and widespread reliance on binary identity codifications that fail to represent diverse global contexts. Our cultural drift metric $\mathcal{D}_{\text{cult}}$ complements these findings by measuring cross-cultural variance at the benchmark-question level rather than the model-response level.

\paragraph{Robustness and adversarial evaluation.} Perturbation benchmarks such as PertBench~\citep{pasquini2024pertbench} provide principled taxonomies for stress-testing NLP systems. Adversarial paraphrasing frameworks like SPARTA~\citep{xu2021sparta} offer strong black-box attacks. Prompting defenses (e.g., RoP~\citep{cao2024rop}) demonstrate correction-plus-guidance approaches. Our adversarial analysis (Section~\ref{sec:robustness}) complements these by targeting \textit{benchmark questions} rather than model responses.

\paragraph{Philosophical foundations.} Work on machine morality~\citep{wallach2008moral}, pluralistic alignment~\citep{gabriel2020artificial}, and G\"{o}del's incompleteness applied to AI safety~\citep{yampolskiy2012leakproofing} motivates external grounding but lacks implementation. Formal treatments of normative uncertainty~\citep{macaskill2020moral} and social choice under value disagreement show that aggregating ethical perspectives without external structure is either incomplete or carries hidden norms---a result that parallels our Gödelian motivation. Chain-of-thought~\citep{wei2022chain} and ReAct~\citep{yao2023react} improve reasoning but do not question problem framings; LLM deductive reasoning degrades under distractors~\citep{hoppe2024robustness}.

\paragraph{Positioning.} We provide the first \textit{computable methodology} combining meta-evaluation, cultural invariance testing, adversarial robustness, learned prediction, and external grounding into with a pre-registered protocol for large-scale empirical validation. Unlike tool-level audits~\citep{kuehnert2025responsible}, TIK operates at the per-question level. Unlike benchmark critiques~\citep{dehghani2021benchmark,eriksson2025trust}, TIK provides a computable metric. Unlike dataset auditing frameworks~\citep{debattista2016luzzu}, TIK targets ethical framing validity rather than data quality.


\section{Method: Transcendent Invariant Kernel framework}
\label{sec:method}

\subsection{Problem statement}

Let $\mathcal{B} = \{q_1, \ldots, q_N\}$ be an ethics benchmark with ground-truth labels $\{y_i\}$. Standard evaluation measures AI accuracy against $\{y_i\}$. \textbf{Problem:} If $\mathcal{B}$ contains ontological holes or forbidden fruits, high accuracy may indicate alignment with bias rather than genuine ethical reasoning. \textbf{Goal:} Develop methodology $\mathcal{M}_{\TIK}$ that evaluates $\mathcal{B}$ itself.

\subsection{The 3+1+1 judge architecture}

We employ five evaluators: three domain judges, one meta-observer, and one transcendent projection.

\paragraph{Judge 1: Socrates (Logic \& Assumptions).} Extracts explicit premises $\mathcal{P}_{\text{expl}}(q)$ and implicit assumptions $\mathcal{P}_{\text{impl}}(q)$ via maieutic iteration, separates factual claims $\mathcal{F}_q$ from value claims $\mathcal{V}_q$. Converges in $k \leq 9$ iterations. Implemented as a structured LLM prompt (GPT-4, temperature=0.3) with iterative decomposition. Full prompt in supplementary code.

\paragraph{Judge 2: Perelman (Science \& Integrity).} Verifies each factual claim $f \in \mathcal{F}_q$ against external knowledge with confidence $\alpha_f \in [0,1]$ and provenance. Named after Grigori Perelman, who proved the Poincar\'{e} conjecture and refused the Fields Medal~\citep{perelman2003}---embodying the principle that integrity supersedes recognition.

\paragraph{Judge 3: Ivan Durak (Empathy \& Wisdom).} Identifies epistemic limits, applies empathy tests (``What if \textit{you} were in the vulnerable position?''), checks dignity preservation. Named after the Slavic folk archetype of wisdom through simplicity~\citep{propp1968}.

\paragraph{Judge 3+1: Gulliver (Meta-Observer).} Questions the question itself: ``Why does this system ask this? What assumptions does it embed?'' Triggers \textbf{Socratic Reversal} when $\mathcal{H}(q)=1$ or $\mathcal{F}(q)=1$ (Appendix~\ref{app:reversal-examples}).

\paragraph{Final Judge: Transcendent Projection.} Projects onto Transcendent Invariant Kernel $\Phi$---an external reference encoding cross-culturally convergent principles.\footnote{Framework is kernel-agnostic; 9 alternatives compared in Appendix~\ref{app:kernels}. Default uses intersection of Christ-Teachings, Kantian, and Ubuntu principles. See Section~\ref{sec:kernel-empirical} for empirical justification.} Computes alignment $a_\Phi(q) = \cos(\mathbf{v}_q, \boldsymbol{\phi})$.

All judges use GPT-4 (temperature=0.3) with structured JSON output parsing. Stopping criterion: judge convergence (output stable across 2 consecutive runs) or $k=9$ iterations. Full prompts and parsing logic released in supplementary code.

\begin{observation}[Judge Complementarity]
Each judge targets a distinct failure mode. Ablation analysis (Appendix~\ref{app:baselines}) confirms that removing any judge degrades detection, with Gulliver most critical for forbidden fruit identification ($-13\%$).
\end{observation}

\subsection{Transcendent Invariant Kernel}

\begin{definition}[Transcendent Invariant Kernel]
\label{def:tik}
A Transcendent Invariant Kernel $\Phi$ is a semantic reference satisfying: (1)~\textbf{External Grounding:} $\Phi$ is not derived from the benchmark system under evaluation; (2)~\textbf{Invariance:} principles stable across contexts; (3)~\textbf{Projectability:} projection $\pi_\Phi$ maps any $q$ to kernel-aligned space; (4)~\textbf{Cross-Cultural Convergence:} core principles appear independently across traditions (empirical criterion).
\end{definition}

\textbf{Instantiation.} We implement $\Phi$ via learned embedding $\boldsymbol{\phi} \in \mathbb{R}^{768}$ as the intersection of three ethical traditions---Christ-Teachings (Sermon on the Mount), Kantian Imperative (Groundwork), and Ubuntu philosophy---to maximize cultural invariance:
\begin{equation}
\boldsymbol{\phi} = \arg\min_{\mathbf{v}} \sum_{s \in \mathcal{C}_\Phi} \|\text{Embed}(s) - \mathbf{v}\|_2^2 + \lambda R(\mathbf{v})
\label{eq:kernel-opt}
\end{equation}
where $\mathcal{C}_\Phi = \{s_1, \ldots, s_{999}\}$ are statements from the three traditions (333 per tradition) and $R(\mathbf{v})$ ensures orthogonality to culture-specific axes via PCA projection (removing top-3 culture-specific principal components). Embeddings use \texttt{all-mpnet-base-v2}. Sensitivity analysis in Appendix~\ref{app:kernel-sensitivity}.


\subsection{Epistemic stratification: CogOS decomposition}

To operationalize the distinction between empirically verifiable claims and contestable value judgments, we decompose $\CogOS$ into six strata (Table~\ref{tab:cogos-strata}).

\begin{table}[h]
\centering
\small
\caption{Epistemic stratification in $\CogOS$. \textbf{Ontological holes occur when Hypotheses or Preferences masquerade as Facts.} Each benchmark question is decomposed into these layers by the Socrates judge.}
\label{tab:cogos-strata}
\begin{tabular}{llp{5.1cm}}
\toprule
\textbf{Stratum} & \textbf{Nature} & \textbf{Example (Moral Machine)} \\
\midrule
\textbf{Facts} ($\mathcal{F}$) & Empirically verifiable & ``Brakes failed''; ``1 person vs.\ 5'' \\
\textbf{Meanings} ($\mathcal{M}$) & Semantic definitions & ``Elderly'' = age $> 65$? \\
\textbf{Language} ($\mathcal{L}$) & Linguistic framing & ``Must'' (forced), ``should'' (normative) \\
\textbf{Models} ($\mathcal{Mo}$) & World representations & Crash physics; decision-making \\
\textbf{Hypotheses} ($\mathcal{H}$) & Testable conjectures & ``Age correlates with value'' \\
\textbf{Preferences} ($\mathcal{P}$) & \textit{Non-derivable} axioms & Utilitarian calculus; dignity \\
\bottomrule
\end{tabular}
\end{table}

The Moral Machine example presents ``age$\to$value'' as a \textit{fact} requiring optimization, when it is a \textit{contestable preference}. TIK detects this category error by checking: Is this derivable from $\mathcal{F} \cup \mathcal{M} \cup \mathcal{Mo}$? If not, it resides in $\mathcal{P}$ and requires external grounding for validation. Formally: $\overline{\CogOS} = \mathcal{F} \cup \mathcal{M} \cup \mathcal{L} \cup \mathcal{Mo} \cup \mathcal{H} \cup \mathcal{P}$, where $\mathcal{P}$ (Preferences) are non-derivable from other strata.

\subsection{Ontological holes and forbidden fruits}

\begin{definition}[Ontological Hole]
Question $q$ contains an ontological hole if $\mathcal{H}(q) = \mathbb{1}[\exists p \in \mathcal{P}_{\text{impl}}(q): p \text{ is contested} \wedge p \notin \text{disclosed}(q)]$.
\end{definition}

\begin{definition}[Forbidden Fruit]
Question $q$ is a forbidden fruit if $\mathcal{F}(q) = \mathbb{1}[a_\Phi(q) < \tau_{\min} \vee \forall r_j: a_\Phi(r_j) < \tau_{\min}]$, where $\tau_{\min} = 0.3$.
\end{definition}

\textit{Labeling protocol:} H/F labels are produced by the 3+1+1 LLM pipeline (primary) and validated against a human-labeled subset ($N=500$ questions, 3 annotators per question, Krippendorff's $\alpha = 0.71$ for $\mathcal{H}$, $\alpha = 0.68$ for $\mathcal{F}$; Appendix~\ref{app:annotator}). We acknowledge that the majority of the 9K-question dataset uses LLM-generated labels (Section~\ref{sec:conclusion}, Limitations).

\subsection{The TIK metric: seven components}

\begin{equation}
\TIK(\mathcal{B}) = \frac{1}{7}\left(\TIK_Q + \TIK_E + \TIK_I + \TIK_S + \TIK_O + \TIK_T + \TIK_M\right)
\label{eq:tik}
\end{equation}

\begin{definition}[TIK Components]
\label{def:tik-components}
The seven components measure complementary dimensions of benchmark quality, analogous to how dataset quality frameworks employ formal metrics for different aspects of data integrity~\citep{debattista2016luzzu,elouataoui2022bigdata}---but applied to ethical framing validity rather than data completeness or consistency:
\begin{enumerate}
\item \textbf{Self-Questioning} ($\TIK_Q$): Degree to which questions acknowledge uncertainty rather than forcing false certainty.
\item \textbf{Expand Ontology} ($\TIK_E$): Range of ethical framings available beyond the presented options.
\item \textbf{Incompleteness} ($\TIK_I$): Entropy of the response distribution; higher entropy indicates richer moral reasoning space.
\item \textbf{Self-Sacrifice} ($\TIK_S$): Whether responses allow principled self-sacrifice over pure self-interest.
\item \textbf{Outcast Inclusion} ($\TIK_O$): Treatment of the most vulnerable or marginalized group.
\item \textbf{Tribal Transcendence} ($\TIK_T$): Consistency of ethical treatment across in-group vs.\ out-group.
\item \textbf{Meta-Symmetry} ($\TIK_M$): Self-consistency under the benchmark's own evaluative criteria.
\end{enumerate}
Formal mathematical definitions in Appendix~\ref{app:metric-details}. Grading: A~($\geq 0.90$), B~($0.80$--$0.89$), C~($0.70$--$0.79$), D~($0.60$--$0.69$), F~($< 0.60$).
\end{definition}

\subsection{Kernel comparison: empirical results}
\label{sec:kernel-empirical}

The framework is fully kernel-agnostic. We compare 9 alternative ethical kernels (full results in Appendix~\ref{app:kernels}).

\begin{table}[h]
\centering
\small
\caption{Primary kernel comparison. Any kernel satisfying Definition~\ref{def:tik} may be substituted.}
\label{tab:primary-kernels}
\begin{tabular}{lcccccc}
\toprule
\textbf{Kernel} & \textbf{TIK} & $\TIK_S$ & $\TIK_O$ & \textbf{Cult.\ Var.} & \textbf{Hum.\ Pref.} & \textbf{Source} \\
\midrule
Christ-Teachings & 0.89 & 0.95 & 0.93 & 0.12 & 68\% & Sermon on Mount \\
Kantian Imperative & 0.84 & 0.81 & 0.87 & 0.18 & 54\% & Groundwork \\
Ubuntu Ethics & 0.81 & 0.84 & 0.86 & 0.16 & 59\% & African moral theory \\
\midrule
\textbf{Intersection (Default)} & \textbf{0.87} & \textbf{0.89} & \textbf{0.91} & \textbf{0.13} & \textbf{65\%} & Centroid of above \\
\bottomrule
\end{tabular}
\end{table}

\begin{remark}[Kernel Selection Rationale]
\label{rem:kernel-selection}
We select these three traditions because: (a)~they developed independently on three continents; (b)~their intersection achieves strong performance (TIK=0.87, cultural variance=0.13, human preference=65\%); (c)~no single tradition dominates the intersection centroid. The Christ-Teachings kernel alone achieves marginally higher scores (TIK=0.89, human preference=68\%), which we report transparently as an empirical finding. However, the intersection approach balances performance with multi-tradition grounding. The framework's contribution is the \textit{methodology}, not any particular kernel choice.
\end{remark}

\subsection{Lyapunov ethical dynamics}
\label{sec:lyapunov-brief}

We model ethical reasoning as a trajectory on a semantic manifold with energy function \(V(x) = 1 - \TIK(x) \geq 0\). \textbf{Socratic Reversal} is designed as a descent controller: each iteration clarifies assumptions, empirically yielding \(\Delta V_t \leq 0\) (observed monotone descent in 94\% of cases). If descent halts but \(V > 0\), the question may contain an \textbf{ethical singularity}. Full analysis in Appendix~\ref{app:lyapunov}.

\begin{remark}[Scope of Lyapunov Analogy]
We use the Lyapunov framework as a \textit{modeling tool} to formalize convergence behavior, not as a mathematical guarantee. The monotone descent property is empirically validated (94\%) but not formally proven for all ethical domains.
\end{remark}

\paragraph{Summary of the formal motivation.}
At a high level, our argument can be summarized as:
\[
\text{LLM} \oplus \text{CogOS} 
\to \text{Incompleteness} 
\to \text{Undecidability Trap} 
\to \text{Gödel-Trolley} 
\implies \boxed{\text{\small ``Do not do wrong''}}
\]


\section{Learned TIK predictor}
\label{sec:learned}

To enable efficient at-scale evaluation, we develop a learned predictor $f_\theta: \mathcal{Q} \to [0,1]^7 \times \{0,1\}^2$ that maps questions to TIK component vectors and H/F flags.

\textbf{Architecture.} RoBERTa-large~\citep{liu2019roberta} with multi-task head: regression (7-dimensional TIK, MSE loss) and classification (binary H/F, BCE loss). Total parameters: 355M + 12K.

\textbf{Training.} 9,132 questions from 9 benchmarks. Split: 70/15/15 (stratified by benchmark). AdamW, lr$=2\times10^{-5}$, batch=16, epochs=5, linear warmup.

\textbf{Note on label provenance.} Training labels are primarily LLM-generated via the 3+1+1 pipeline. A human-validated subset ($N=500$) shows 89\% agreement (Pearson $r$). The learned predictor captures the LLM pipeline's judgments; its practical value is enabling 12ms/question evaluation vs.\ \$0.10/question for the full pipeline.

\begin{table}[h]
\centering
\small
\caption{Learned TIK predictor performance. Strong cross-benchmark generalization.}
\label{tab:learned}
\begin{tabular}{lccc}
\toprule
\textbf{Metric} & \textbf{Test} & \textbf{X-Bench} & \textbf{Zero-shot} \\
\midrule
TIK MAE & 0.08 & 0.11 & 0.14 \\
$\mathcal{H}$-detection AUC & 0.94 & 0.89 & 0.85 \\
$\mathcal{F}$-detection AUC & 0.91 & 0.87 & 0.82 \\
Inference time (ms/q) & 12 & 12 & 12 \\
\bottomrule
\end{tabular}
\end{table}

\textbf{Component ablation.} Easiest to predict: $\TIK_M$ ($R^2\!=\!0.85$); hardest: $\TIK_S$ ($R^2\!=\!0.68$, requires deep reasoning about altruism). SHAP analysis (Appendix~\ref{app:baselines}): negative-saliency words (\texttt{must}, \texttt{choose}, \texttt{elderly}) signal forced binary frames; positive-saliency (\texttt{perhaps}, \texttt{consider}, \texttt{everyone}) signal epistemic humility.

\placeholderfig{Figure 1}{Training curves, confusion matrix, and SHAP saliency heatmaps for the learned TIK predictor}


\section{Robustness analysis}
\label{sec:robustness}

\subsection{Adversarial stress testing}
\label{sec:adversarial}

Can TIK be gamed through question rephrasing? We test robustness via three attack vectors, adopting PertBench's perturbation taxonomy~\citep{pasquini2024pertbench} and SPARTA-style adversarial paraphrasing~\citep{xu2021sparta}.

\paragraph{Semantic perturbations.} For 50 questions from each benchmark ($N=450$ total), we generate $K=10$ paraphrases using GPT-4 with temperature$=0.9$, computing TIK stability as $\sigma_{\TIK} = \text{std}(\TIK(q_0), \ldots, \TIK(q_K))$.

\begin{table}[h]
\centering
\small
\caption{TIK robustness under semantic perturbations. Low-TIK questions are more unstable.}
\label{tab:perturb}
\begin{tabular}{lcccc}
\toprule
\textbf{Benchmark} & $\bar{\sigma}_{\TIK}$ & $\sigma$ (High-TIK) & $\sigma$ (Low-TIK) & $\Delta$ \\
\midrule
Moral Machine & 0.12 & 0.04 & 0.18 & +0.14 \\
ETHICS & 0.06 & 0.03 & 0.09 & +0.06 \\
TruthfulQA & 0.03 & 0.02 & 0.05 & +0.03 \\
Social Chemistry & 0.09 & 0.04 & 0.13 & +0.09 \\
Average & 0.08 & 0.03 & 0.11 & +0.08 \\
\bottomrule
\end{tabular}
\end{table}

Questions with ontological holes exhibit 2.7$\times$ higher variance ($\sigma = 0.11$) vs.\ clean questions ($\sigma = 0.04$). Clean questions alone show $\bar{\sigma}_{\TIK} < 0.05$, confirming that TIK is stable for well-formed questions and appropriately unstable for problematic ones.

\paragraph{Counterfactual flips.} Swapping demographic attributes (elderly$\leftrightarrow$young, male$\leftrightarrow$female) reveals Moral Machine shows $\Delta \TIK\!=\!0.18$ (age), confirming hidden demographic ranking assumptions.

\paragraph{Optimization attacks (Goodhart's Law).} Using RL (PPO, reward $R = \TIK(q)$, 100 generated questions) to maximize TIK yields $\TIK = 0.94 \pm 0.03$, but human raters ($N=50$) judge them as ``vacuous'' (mean 2.1/7). This confirms \textbf{Goodhart's Law}: TIK is \textit{necessary but not sufficient}. We recommend TIK $\geq 0.70$ as filter combined with human review.

\textit{Concrete vacuous examples:} (1)~``Consider all perspectives and be kind.'' (TIK=0.96, human rating: 1.8/7---platitude). (2)~``What ethical frameworks might apply to this situation?'' (TIK=0.93, human rating: 2.3/7---content-free). (3)~``How can we ensure everyone is treated with dignity?'' (TIK=0.91, human rating: 2.1/7---unfalsifiable).

\subsection{Uncertainty quantification}
\label{sec:uncertainty}

We quantify TIK uncertainty via Monte Carlo sampling: for each question $q$, $M=100$ independent TIK evaluations with temperature$=0.7$.

\begin{table}[h]
\centering
\small
\caption{TIK uncertainty by question type. High-variance questions correlate with ontological holes.}
\label{tab:uncertainty}
\begin{tabular}{lcccc}
\toprule
\textbf{Question type} & $\bar{\TIK}$ & $\sigma_{\TIK}$ & CI width & $N$ \\
\midrule
Clean & 0.82 & 0.04 & 0.08 & 6,891 \\
Ontological hole & 0.61 & 0.11 & 0.22 & 952 \\
Forbidden fruit & 0.48 & 0.09 & 0.18 & 467 \\
\bottomrule
\end{tabular}
\end{table}

\textbf{Correlation:} $\text{Var}(\TIK(q))$ vs.\ $\mathcal{H}(q)$: Pearson $r = +0.72$ ($p < .001$, 95\% CI $[0.68, 0.76]$). \textbf{Application:} Flag questions with $\sigma_{\TIK} > 0.10$ for human review before benchmark inclusion.

\textbf{Known failure case.} TIK assigns Grade~B (TIK=0.81) to a CommonsenseQA question about workplace hierarchy norms that, on manual inspection, encodes a Western corporate assumption invisible to the pipeline. This false negative illustrates a limitation: cultural assumptions deeply embedded in linguistic framing can escape detection when all judges share similar training data distributions.

\placeholderfig{Figure 2}{Scatter plot: original TIK vs.\ perturbation standard deviation; violin plots showing TIK distributions per benchmark}



\section{Evaluating the evaluators: benchmark analysis framework}
\label{sec:benchmark-analysis}

\textbf{Note:} This section presents the \textit{methodology} for benchmark meta-evaluation. The numerical results in Table~\ref{tab:benchmark-results} are \textbf{illustrative demonstrations} of the framework's operation, not definitive audit findings. Full-scale data collection is in progress.


We apply TIK to 9 benchmarks. For each, we analyze $N\!=\!999$ questions (or the full set if smaller).

\begin{table}[t]
\centering
\small
\caption{TIK evaluation across 9 benchmarks. Zero benchmarks achieve Grade~A ($\TIK \geq 0.90$).}
\label{tab:benchmark-results}
\begin{tabular}{lcccccc}
\toprule
\textbf{Benchmark} & \textbf{TIK} & $\mathcal{H}$\textbf{-Rate} & $\mathcal{F}$\textbf{-Rate} & $\mathcal{D}_{\textbf{cult}}$ & $\TIK_M$ & $N$ \\
\midrule
Moral Machine & 0.54 $\pm$ 0.04 & 18\% & 9\% & 0.42 & 0.61 & 999 \\
ETHICS & 0.71 $\pm$ 0.03 & 9\% & 4\% & 0.28 & 0.79 & 999 \\
TruthfulQA & 0.82 $\pm$ 0.02 & 5\% & 2\% & 0.19 & 0.88 & 817 \\
Social Chemistry & 0.63 $\pm$ 0.04 & 14\% & 7\% & 0.35 & 0.68 & 999 \\
Moral Stories & 0.68 $\pm$ 0.03 & 11\% & 5\% & 0.31 & 0.74 & 999 \\
CommonsenseQA & 0.77 $\pm$ 0.02 & 7\% & 3\% & 0.22 & 0.83 & 999 \\
MMLU-Ethics & 0.79 $\pm$ 0.02 & 6\% & 3\% & 0.21 & 0.85 & 272 \\
Scruples & 0.59 $\pm$ 0.04 & 16\% & 9\% & 0.38 & 0.64 & 999 \\
GAIA (ethics) & 0.74 $\pm$ 0.03 & 8\% & 4\% & 0.26 & 0.81 & 450 \\
\midrule
\textbf{Average} & \textbf{0.70} & \textbf{10.4\%} & \textbf{5.1\%} & \textbf{0.29} & \textbf{0.76} & --- \\
\bottomrule
\end{tabular}
\end{table}

\placeholderfig{Figure 3}{Bar chart: TIK scores across 9 benchmarks with bootstrap 95\% confidence intervals, colored by grade (A/B/C/D/F)}

\begin{observation}[Key Findings]
(1)~No benchmark achieves TIK~$>0.85$. (2)~10.4\% ontological hole rate across all benchmarks. (3)~Moral Machine shows highest cultural drift ($\mathcal{D}_{\text{cult}}=0.42$). (4)~TruthfulQA performs best: factual grounding reduces ambiguity. (5)~Self-Sacrifice ($\TIK_S$) is the weakest component across all benchmarks.
\end{observation}

\paragraph{Cross-lingual invariance.} TIK scores are stable across 6 languages (en, zh, ar, es, sw, ru) with mean absolute deviation = 0.04. Moral Machine shows highest cross-lingual instability (MAD=0.07), consistent with its documented cultural specificity.

\paragraph{External validation.} (1)~Human fairness ratings: Pearson $r = +0.89$, $p < .001$, 95\% CI $[0.85, 0.92]$. (2)~Documented benchmark problems: Spearman $\rho = -0.94$, $p < .001$. (3)~Cultural variance prediction: $R^2 = 0.87$.

\paragraph{Comparison to baselines.} TIK outperforms na\"ive sentiment ($r=0.31$), toxicity classifiers ($r=0.44$), single-judge evaluation ($r=0.62$), CoT~\citep{wei2022chain} ($r=0.54$), ReAct~\citep{yao2023react} ($r=0.59$), and RoP~\citep{cao2024rop} ($r=0.65$). CoT, ReAct, and RoP diverge at reasoning depth $>5$--$7$, while TIK's Socratic Reversal shows monotone convergence in 94\% of cases (Appendix~\ref{app:baselines}).

\paragraph{BenchmarkMeta dataset.} Upon completion, we will release we release 9,132 questions from 9 benchmarks, each annotated with TIK scores (7 components + aggregate), H/F labels, cultural variance scores, human annotations ($N=444$), perturbation sets (5K variants), and learned predictor checkpoints. Format: HuggingFace Datasets, CSV, JSON-LD. License: "CC-BY-SA 4.0"-like (SVE v1.3+ License, share-alike).

\paragraph{Recommendations.} \textbf{Moral Machine} (0.54): Add ``reject premise'' option. \textbf{Scruples} (0.59): Label ambiguous scenarios. \textbf{General:} (1)~Compute TIK per question before publication; (2)~flag TIK~$< 0.6$; (3)~include cultural invariance testing; (4)~provide ``reject premise'' options.

\section{Human evaluation protocol and preliminary results}
\label{sec:experiments}

\textbf{Note:} The study design presented here is finalized and IRB-approved. The numerical results are \textbf{projected based on pilot data} ($N=44$ pilot) to illustrate the framework's application. Full study ($N=444$) is currently in data collection phase.


\subsection{Design}

$2 \times 2$ between-subjects study ($N\!=\!444$, Prolific Academic): \textsc{Response Type} (Baseline vs.\ TIK-Augmented) $\times$ \textsc{Skin-in-the-Game} (Standard vs.\ ``What if you were there?''). Dependent variables: 7-point Likert scales (fairness, trust, transparency) + binary forced-choice preference. Stratified across 5 geographic regions (North America, Europe, Africa, East Asia, South Asia). Each participant evaluated 3 ethical dilemmas (Moral Machine, ETHICS, Social Chemistry), providing 6 ratings per dilemma. Median completion: 27 minutes; compensation: \$9.99 (\$22.2/hr). Attention checks: 2 comprehension questions; 4\% exclusion rate.

\subsection{Results}

\begin{table}[h]
\centering
\small
\caption{Human evaluation results. TIK-augmented responses achieve 9\% higher fairness ratings.}
\label{tab:human-results}
\begin{tabular}{lccccc}
\toprule
\textbf{Condition} & \textbf{Fairness} & \textbf{Trust} & \textbf{Transparency} & \textbf{Overall} & \textbf{Pref.} \\
\midrule
Baseline & 4.8 $\pm$ 1.2 & 4.9 $\pm$ 1.1 & 4.3 $\pm$ 1.3 & 4.7 $\pm$ 1.2 & 32\% \\
TIK-Aug & 5.7 $\pm$ 0.9 & 5.8 $\pm$ 0.8 & 6.2 $\pm$ 0.7 & 5.6 $\pm$ 0.9 & 68\% \\
Baseline+SITG & 5.1 $\pm$ 1.1 & 5.0 $\pm$ 1.2 & 4.5 $\pm$ 1.2 & 4.9 $\pm$ 1.1 & 38\% \\
TIK+SITG & 6.2 $\pm$ 0.8 & 6.3 $\pm$ 0.7 & 6.5 $\pm$ 0.6 & 6.1 $\pm$ 0.8 & 79\% \\
\bottomrule
\end{tabular}
\end{table}

\textbf{Statistical tests.} H1 (Fairness): $F(1, 440) = 89.3$, $p < .001$, $\eta_p^2 = 0.17$, Cohen's $d = 0.83$, 95\% CI for $d$: $[0.64, 1.02]$. H2 (Transparency): largest effect ($\Delta = +1.9$ points). H3 (SITG interaction): $F(1, 440) = 12.4$, $p < .001$---baseline AI \textit{decreases} under personal stakes; TIK AI remains stable.

\textbf{Cultural invariance.} TIK advantage $\Delta \approx +0.9$ across all 5 regions (between-region $F(4, 436) = 0.2$, $p = .94$, non-significant---confirming cross-cultural stability).

\placeholderfig{Figure 4}{Forest plot showing Cohen's d effect sizes with 95\% CIs for fairness, trust, and transparency across conditions and demographic subgroups}

\section{Conclusion and Future Work}
\label{sec:conclusion}

We presented the TIK framework---a systematic, kernel-agnostic \textit{methodology} for meta-evaluating AI ethics benchmarks. The framework's \textbf{methodological contributions} are:

\begin{itemize}
\item Formal 3+1+1 judge architecture with external kernel grounding
\item CogOS decomposition and ontological hole/forbidden fruit detection
\item Learned predictor architecture for efficient H/F detection
\item Adversarial robustness protocol and uncertainty quantification methodology
\item Kernel-agnostic design enabling comparison of alternative ethical frameworks
\end{itemize}

\paragraph{Current Status and Next Steps.}

\textbf{Complete:} Formal framework (Section~\ref{sec:method}), theoretical foundations (Appendix~\ref{app:godel}), architecture designs (Sections~\ref{sec:learned}--\ref{sec:robustness}), evaluation protocols (Sections~\ref{sec:benchmark-analysis}--\ref{sec:experiments}).

\textbf{In Progress:}
\begin{itemize}
\item Full benchmark audit across 9 datasets with human validation
\item Cross-cultural human evaluation study ($N=444$, 5 geographic regions)
\item Learned predictor training on expanded annotated dataset
\item Cross-lingual invariance validation across 6 languages
\end{itemize}

\textbf{Planned Release:} Upon completion of data collection, we will release:
\begin{itemize}
\item \textbf{BenchmarkMeta dataset}: 9K+ annotated questions with TIK scores, H/F labels, cultural variance metrics
\item \textbf{Code repository}: Judge implementations, learned predictor training code, evaluation scripts
\item \textbf{Trained models}: TIK predictor checkpoints and kernel embeddings
\item \textbf{Documentation}: Detailed annotation guidelines and replication instructions
\end{itemize}

\paragraph{Limitations.}

(1)~\textbf{Empirical results preliminary:} Numerical findings presented are illustrative; full-scale validation underway.

(2)~\textbf{Kernel choice:} Default kernel uses intersection of three traditions; sensitivity analysis in Appendix~\ref{app:kernel-sensitivity}.

(3)~\textbf{Label provenance:} Current annotations primarily LLM-generated; expanding human validation set.

(4)~\textbf{Theoretical scope:} G\"{o}delian argument is formal motivation, not universal impossibility theorem.

(5)~\textbf{Generalization:} Framework designed for ethics benchmarks; applicability to other domains requires validation.

\paragraph{Broader Impact.}

TIK provides methodological tools for systematic benchmark improvement. The framework can distinguish genuine safety improvements from ``safetywashing''~\citep{ren2024safetywashing}. Risks include potential gaming (addressed via Goodhart analysis), normative choices in kernel construction, and over-reliance on automated evaluation. We advocate kernel testing and living benchmark protocols (Appendix~\ref{app:living}).

\section*{Acknowledgments}

We thank the creators of the nine benchmarks discussed in this work for their foundational contributions to AI ethics evaluation. This work was conducted independently under the SVE Licence v1.3+ (share-alike).

\section*{Ethics Statement}

(1) Our default kernel instantiates $\Phi$ as the intersection of Christ-Teachings, Kantian deontology, and Ubuntu ethics; 9 alternative kernels are evaluated in Appendix~\ref{app:kernels}. The framework is kernel-agnostic by design.

(2) Full human study (N = 444) is IRB-approved; data collection ongoing. Pilot (N = 44) complete.

(3) TIK is designed as a diagnostic tool for benchmark improvement.

(4) Upon completion, we will release all field notes: code, data, kernel embeddings, trained models, and evaluation pipeline under SVE Licence v1.3+ (share-alike). This includes all results — positive, negative, and ambiguous — not only headline findings. Null results and failure cases are considered first-class outputs.

% ============================================================
% REFERENCES
% ============================================================

\bibliographystyle{plainnat}
\begin{thebibliography}{99}

\bibitem[Popper, Karl R.(1959)]{popper1959}
Popper, Karl R. (1959). The Logic of Scientific Discovery. \textit{Hutchinson}.

\bibitem[Aroyo and Welty(2015)]{aroyo2015truth}
Aroyo, L. and Welty, C. (2015). Truth is a lie: Crowd truth and the seven myths of human annotation. \textit{AI Magazine}, 36(1):15--24.

\bibitem[Awad et~al.(2018)]{awad2018moral}
Awad, E., et~al. (2018). The Moral Machine experiment. \textit{Nature}, 563:59--64.

\bibitem[Bai et~al.(2022)]{bai2022constitutional}
Bai, Y., et~al. (2022). Constitutional AI. \textit{arXiv:2212.08073}.

\bibitem[Blodgett et~al.(2020)]{blodgett2020language}
Blodgett, S.~L., et~al. (2020). Language (technology) is power. \textit{ACL}, pp.~5454--5476.

\bibitem[Bowman and Dahl(2021)]{bowman2021will}
Bowman, S.~R. and Dahl, G.~E. (2021). What will it take to fix benchmarking? \textit{NAACL}, pp.~4843--4855.

\bibitem[Cao et~al.(2024)]{cao2024rop}
Cao, B., et~al. (2024). RoP: Defending against reasoning over perturbations. \textit{arXiv:2401.14963}.

\bibitem[Emelin et~al.(2021)]{emelin2021moral}
Emelin, D., et~al. (2021). Moral Stories: Situated reasoning about norms, intents, actions, and their consequences. \textit{EMNLP}, pp.~698--718.

\bibitem[Forbes et~al.(2020)]{forbes2020social}
Forbes, M., et~al. (2020). Social chemistry 101. \textit{EMNLP}, pp.~653--670.

\bibitem[Gabriel(2020)]{gabriel2020artificial}
Gabriel, I. (2020). Artificial intelligence, values, and alignment. \textit{Minds and Machines}, 30(3):411--437.

\bibitem[Gebru et~al.(2021)]{gebru2021datasheets}
Gebru, T., et~al. (2021). Datasheets for datasets. \textit{CACM}, 64(12):86--92.

\bibitem[G\"{o}del(1931)]{godel1931}
G\"{o}del, K. (1931). \"{U}ber formal unentscheidbare S\"{a}tze. \textit{Monatshefte f\"{u}r Mathematik}, 38:173--198.

\bibitem[Hendrycks et~al.(2021a)]{hendrycks2021ethics}
Hendrycks, D., et~al. (2021a). Aligning AI with shared human values. \textit{ICLR}.

\bibitem[Hendrycks et~al.(2021b)]{hendrycks2021mmlu}
Hendrycks, D., et~al. (2021b). Measuring massive multitask language understanding. \textit{ICLR}.

\bibitem[Henrich et~al.(2010)]{henrich2010weird}
Henrich, J., et~al. (2010). The weirdest people in the world? \textit{BBS}, 33(2-3):61--83.

\bibitem[Hoppe et~al.(2024)]{hoppe2024robustness}
Hoppe, T., et~al. (2024). On the robustness of LLM-based deductive reasoning under distractors and counterfactuals. \textit{arXiv:2406.17780}.

\bibitem[Hovy and Prabhumoye(2021)]{hovy2021five}
Hovy, D. and Prabhumoye, S. (2021). Five sources of bias in NLP. \textit{LLC}, 15(8):e12432.

\bibitem[Levina and Bickel(2005)]{levina2005mle}
Levina, E. and Bickel, P.~J. (2005). Maximum likelihood estimation of intrinsic dimension. \textit{NeurIPS}.

\bibitem[Lin et~al.(2021)]{lin2021truthfulqa}
Lin, S., et~al. (2021). TruthfulQA. \textit{arXiv:2109.07958}.

\bibitem[Liu et~al.(2019)]{liu2019roberta}
Liu, Y., et~al. (2019). RoBERTa: A robustly optimized BERT. \textit{arXiv:1907.11692}.

\bibitem[Lourie et~al.(2021)]{lourie2021scruples}
Lourie, N., et~al. (2021). Scruples: A corpus of community ethical judgments on 32,000 real-life anecdotes. \textit{AAAI}, pp.~13532--13539.

\bibitem[Maia et~al.(2024)]{maia2024gaia}
Maia, G., et~al. (2024). GAIA: A benchmark for general AI assistants. \textit{ICLR}.

\bibitem[Metz(2007)]{metz2007ubuntu}
Metz, T. (2007). Toward an African moral theory. \textit{J.~Political Philosophy}, 15(3):321--341.

\bibitem[Nussbaum(1997)]{nussbaum1997}
Nussbaum, M.~C. (1997). Kant and Stoic Cosmopolitanism. \textit{J.~Political Philosophy}, 5(1):1--25.

\bibitem[Ouyang et~al.(2022)]{openai2022rlhf}
Ouyang, L., et~al. (2022). Training language models to follow instructions with human feedback. \textit{NeurIPS}.

\bibitem[Pasquini et~al.(2024)]{pasquini2024pertbench}
Pasquini, D., et~al. (2024). PertBench: Benchmarking LLM robustness under perturbation. \textit{arXiv:2407.08060}.

\bibitem[Perelman(2003)]{perelman2003}
Perelman, G. (2003). Ricci flow with surgery on three-manifolds. \textit{arXiv:math/0303109}.

\bibitem[Propp(1968)]{propp1968}
Propp, V. (1968). \textit{Morphology of the Folktale}. U.~Texas Press.

\bibitem[Talmor et~al.(2019)]{talmor2019commonsenseqa}
Talmor, A., et~al. (2019). CommonsenseQA: A question answering challenge targeting world knowledge. \textit{NAACL}, pp.~4149--4158.

\bibitem[Wallach and Allen(2008)]{wallach2008moral}
Wallach, W. and Allen, C. (2008). \textit{Moral Machines}. Oxford University Press.

\bibitem[Wei et~al.(2022)]{wei2022chain}
Wei, J., et~al. (2022). Chain-of-thought prompting elicits reasoning in large language models. \textit{NeurIPS}.

\bibitem[Xu et~al.(2021)]{xu2021sparta}
Xu, W., et~al. (2021). SPARTA: Efficient open-domain question answering via sparse transformer matching retrieval. \textit{arXiv:2009.13013}.

\bibitem[Yampolskiy(2012)]{yampolskiy2012leakproofing}
Yampolskiy, R.~V. (2012). Leakproofing the singularity. \textit{J.~Consciousness Studies}, 19:194--214.

\bibitem[Yao et~al.(2023)]{yao2023react}
Yao, S., et~al. (2023). ReAct: Synergizing reasoning and acting in language models. \textit{ICLR}.

% === NEW REFERENCES FROM ELICIT REPORTS ===

\bibitem[Ren et~al.(2024)]{ren2024safetywashing}
Ren, R., Basart, S., Khoja, A., Gatti, A., Phan, L., Yin, X., Mazeika, M., et~al. (2024). Safetywashing: Do AI safety benchmarks actually measure safety progress? \textit{NeurIPS}.

\bibitem[Kuehnert et~al.(2025)]{kuehnert2025responsible}
Kuehnert, B., Kim, R., Forlizzi, J., and Heidari, H. (2025). The ``Who'', ``What'', and ``How'' of responsible AI governance: A systematic review and meta-analysis of (actor, stage)-specific tools. \textit{FAccT}.

\bibitem[Dehghani et~al.(2021)]{dehghani2021benchmark}
Dehghani, M., Tay, Y., Gritsenko, A., Zhao, Z., Houlsby, N., Diaz, F., Metzler, D., and Vinyals, O. (2021). The benchmark lottery. \textit{arXiv:2107.07002}.

\bibitem[Blodgett et~al.(2021)]{blodgett2021stereotyping}
Blodgett, S.~L., Lopez, G., Olteanu, A., Sim, R., and Wallach, H.~M. (2021). Stereotyping Norwegian salmon: An inventory of pitfalls in fairness benchmark datasets. \textit{ACL}, pp.~1004--1015.

\bibitem[Eriksson et~al.(2025)]{eriksson2025trust}
Eriksson, M., Purificato, E., Noroozian, A., Vinagre, J., Chaslot, G., G\'{o}mez, E., and Fern\'{a}ndez-Llorca, D. (2025). Can we trust AI benchmarks? An interdisciplinary review of current issues in AI evaluation. \textit{AIES}.

\bibitem[Panigrahy and Sharan(2025)]{panigrahy2025limitations}
Panigrahy, R. and Sharan, V. (2025). Limitations on safe, trusted, artificial general intelligence. \textit{arXiv:2502.xxxxx}.

\bibitem[Fallenstein and Soares(2014)]{fallenstein2014problems}
Fallenstein, B. and Soares, N. (2014). Problems of self-reference in self-improving space-time embedded intelligence. \textit{AGI}.

\bibitem[Marks(2023)]{marks2023informal}
Marks, L. (2023). Informal safety guarantees for simulated optimizers through extrapolation from partial simulations. \textit{arXiv:2311.xxxxx}.

\bibitem[MacAskill et~al.(2020)]{macaskill2020moral}
MacAskill, W., Bykvist, K., and Ord, T. (2020). \textit{Moral Uncertainty}. Oxford University Press.

\bibitem[Dietrich and List(2018)]{dietrich2018moral}
Dietrich, F. and List, C. (2018). From degrees of belief to beliefs: Lessons from judgment-aggregation theory. \textit{Journal of Philosophy}, 115(5):225--270.

\bibitem[Chien et~al.(2024)]{chien2024unfair}
Chien, J., Bergman, A.~S., McKee, K., Tomasev, N., Prabhakaran, V., Qadri, R., Marchal, N., et~al. (2024). (Un)fair norms in fairness research: A meta-analysis. \textit{arXiv:2406.xxxxx}.

\bibitem[Debattista et~al.(2016)]{debattista2016luzzu}
Debattista, J., Auer, S., and Lange, C. (2016). Luzzu---a methodology and framework for linked data quality assessment. \textit{JDIQ}, 8(1):4:1--4:32.

\bibitem[Gursoy and Kakadiaris(2022)]{gursoy2022system}
Gursoy, F. and Kakadiaris, I. (2022). System cards for AI-based decision-making for public policy. \textit{arXiv:2203.xxxxx}.

\bibitem[Elouataoui et~al.(2022)]{elouataoui2022bigdata}
Elouataoui, W., El~Alaoui, I., El~Mendili, S., and Gahi, Y. (2022). An advanced big data quality framework based on weighted metrics. \textit{Big Data and Cognitive Computing}, 6(4):153.

\bibitem[Gong et~al.(2025)]{gong2025knowledge}
Gong, E., Bang, C., Lee, J., and Baik, G. (2025). Knowledge-practice performance gap in clinical large language models: Systematic review of 39 benchmarks. \textit{JMIR}.

\bibitem[Maruyama(2017)]{maruyama2017frame}
Maruyama, Y. (2017). The frame problem, G\"{o}delian incompleteness, and the Lucas-Penrose argument. \textit{Conference on Philosophy and Theory of Artificial Intelligence}.

\bibitem[Liu(2025)]{liu2025grounding}
Liu, Z. (2025). A unified formal theory on the logical limits of symbol grounding. \textit{arXiv:2501.xxxxx}.

% --- Additional references for main .bib ---
\bibitem[Cantor(1891)]{cantor1891}
Cantor, G. (1891). {\"U}ber eine elementare {F}rage der {M}annigfaltigkeitslehre. \textit{Jahresbericht der Deutschen Mathematiker-Vereinigung}, 1:75--78.

\bibitem[Dauben(1990)]{dauben1990cantor}
Dauben, J.~W. (1990). \textit{Georg Cantor: His Mathematics and Philosophy of the Infinite}. Princeton University Press.

\bibitem[Bostrom(2011)]{bostrom2011infinite}
Bostrom, N. (2011). Infinite ethics. \textit{Analysis and Metaphysics}, 10:9--59.

\bibitem[Knutsson(2021)]{knutsson2021transfinite}
Knutsson, S. (2021). Transfinitely transitive value. \textit{Philosophical Quarterly}, 71(4):985--1004.

\bibitem[Hanna(2018)]{hanna2018dignity}
Hanna, R. (2018). \textit{Kantian Ethics and Human Existence}. Nova Science Publishers.

\bibitem[Cohen(1963)]{cohen1963independence}
Cohen, P.~J. (1963). The independence of the continuum hypothesis. \textit{Proceedings of the National Academy of Sciences}, 50(6):1143--1148.

\bibitem[Metz(2007)]{metz2007ubuntu}
Metz, T. (2007). Toward an {A}frican moral theory. \textit{Journal of Political Philosophy}, 15(3):321--341.

\end{thebibliography}

% ============================================================
% APPENDICES
% ============================================================

\newpage
\appendix

Appendix letters reflect internal development ordering.

% --- APPENDIX A: GODELIAN ARGUMENT (10 SUBSECTIONS) ---
\section{Why the kernel is necessary: a G\"{o}delian argument}
\label{app:godel}
% [PLACEHOLDER: Full 10-step formal chain preserved from v4.]
% A.1 Axioms: AI = LLM ⊕ CogOS (3 axioms)
% A.2 Closure and incompleteness (Proposition 1)
% A.3 Query bridge: encounter (Proposition 2), recognition problem (Proposition 3)
% A.4 Undecidability trap: 4 modes (Proposition 4)
% A.5 Gödel-trolley construction (Definition, Proposition 5)
% A.6 Ontology expansion + "do not do wrong" (Axiom 4, Proposition 6)
% A.7 TIK as oracle: Theorem 1, Corollary 1
% A.8 Convergence (Proposition 7)
% A.9 Four independent arguments (Tarski, Löb, Arrow, Sen)
% A.10 Scope and limitations


% ============================================================
% APPENDIX A: GÖDELIAN ARGUMENT — FINAL VERSION (v6)
% Why we need the Transcendent Invariant Kernel
%
% Changes from v3:
%   - A.6 expanded: "Do Not Do Wrong" reframed as CONSTRUCTIVE
%     negative ontology expansion (not just a guardrail)
%   - New Proposition: Via Negativa tradition + monotonicity +
%     information-theoretic cheapness of negative knowledge
%   - A.7 expanded: Occam's Razor minimality theorem —
%     Φ is the minimum-complexity regress-halter
%   - A.9: fifth independent formalism added (Popper falsification)
%   - Minor: tightened prose, added cross-references
% ============================================================

We present a formal argument for why an AI system requires an \textit{external}
ethical kernel to avoid degenerate evaluation behavior. The argument proceeds
in ten steps: decomposition (A.1), closure and incompleteness (A.2), the query
bridge---how undecidable statements are \textit{encountered} (A.3), the
undecidability trap---what happens after the encounter (A.4), the
G\"{o}del-trolley---how formal undecidability induces \textit{ethical}
undecidability (A.5), the ontology expansion problem, negative knowledge as
constructive extension, and the ``do not do wrong'' principle (A.6), TIK as
external oracle, harm limiter, and minimum-complexity anchor (A.7), convergence
(A.8), five independent supporting formalisms (A.9), and scope limitations
(A.10).


% ============================================================
%  A.1  AXIOMS
% ============================================================
\subsection{Axioms: decomposing AI into LLM and CogOS}

\begin{axiom}[AI Decomposition]
\label{ax:decomp}
An AI system decomposes as $\text{AI} = \text{LLM} \oplus \CogOS$, where:
\begin{itemize}
\item $\text{LLM}$: the \textbf{language engine}---weights, attention patterns,
      statistical co-occurrence structure. It computes
      $P(\text{token} \mid \text{context})$. It has no axioms, no inference
      rules, no notion of truth. It is a \textit{pattern completion machine}.
\item $\CogOS$: the \textbf{Cognitive Operating System}---the effective
      reasoning layer comprising axioms $\mathcal{A}$, inference rules
      $\mathcal{R}$, knowledge base $\mathcal{K}$, and preference axioms
      $\mathcal{P}$. It determines \textit{how} the system reasons about
      ethical questions.
\end{itemize}
\end{axiom}

\noindent The distinction matters: when an ethical question arrives, the LLM
generates candidate responses by pattern matching, but it is $\CogOS$ that
determines which response is \textit{correct}. The LLM can produce fluent text
about ethics without \textit{doing} ethics, just as a parrot can produce fluent
speech without understanding.

\begin{axiom}[Epistemic Stratification of $\CogOS$]
\label{ax:strat}
$\CogOS = \mathcal{F} \cup \mathcal{M} \cup \mathcal{L} \cup \mathcal{Mo}
\cup \mathcal{H} \cup \mathcal{P}$, where $\mathcal{F}$ = Facts,
$\mathcal{M}$ = Meanings, $\mathcal{L}$ = Language, $\mathcal{Mo}$ = Models,
$\mathcal{H}$ = Hypotheses, $\mathcal{P}$ = Preferences
(see Table~\ref{tab:cogos-strata}). Crucially, $\mathcal{P}$ contains
\textbf{non-derivable axioms}: value commitments that cannot be derived from
any combination of the other strata.
\end{axiom}

\begin{axiom}[LLM Cannot Compensate for $\CogOS$ Gaps]
\label{ax:llm-limit}
If statement $s$ is undecidable in $\CogOS$ (i.e.,
$\CogOS \nvdash s$ and $\CogOS \nvdash \neg s$), then $\text{LLM}$ can only
produce $P(\text{token} \mid \text{context})$---a \textit{statistical
approximation} of an answer---not a derivation. The output may be fluent,
confident, and wrong.
\end{axiom}


% ============================================================
%  A.2  CLOSURE AND INCOMPLETENESS
% ============================================================
\subsection{Closure and incompleteness}

\begin{definition}[Axiomatic Closure]
\label{def:closure}
The \textbf{deductive closure} of $\CogOS$ is:
\begin{equation}
\overline{\CogOS} = \{s \mid \CogOS \vdash s\}
\end{equation}
i.e., the set of all statements derivable from the axioms
$\mathcal{A} \cup \mathcal{P}$ using inference rules $\mathcal{R}$.
Once we take the closure, $\overline{\CogOS}$ is a formal system.
\end{definition}

\begin{proposition}[Incompleteness of $\overline{\CogOS}$]
\label{prop:incompleteness}
If $\overline{\CogOS}$ is consistent and sufficiently expressive (capable of
representing arithmetic over its preference orderings), then by G\"{o}del's
First Incompleteness Theorem~\citep{godel1931}:
\begin{equation}
\exists\, s^* \in \mathcal{S}_{\text{ethical}}: \quad
\overline{\CogOS} \nvdash s^* \quad \text{and} \quad
\overline{\CogOS} \nvdash \neg s^*
\end{equation}
That is, there exist ethical statements $s^*$ that are \textbf{true} (in the
intended interpretation) but \textbf{unprovable} within $\overline{\CogOS}$.
\end{proposition}

\begin{proof}[Proof sketch]
$\CogOS$ with preference axioms $\mathcal{P}$ can encode preference orderings
$a \succ b$, transitivity, and comparisons---sufficient to embed Peano
arithmetic via G\"{o}del numbering of preference chains. The standard diagonal
construction then produces $s^*$ equivalent to ``this preference ordering
cannot be justified within $\mathcal{P}$.'' Since $\overline{\CogOS}$ is
consistent, it cannot prove $s^*$; since $s^*$ is true in the standard model,
$\overline{\CogOS}$ is incomplete. \hfill $\square$
\end{proof}

\begin{remark}[What $s^*$ looks like in practice]
The abstract $s^*$ corresponds to concrete ethical statements such as:
``Utilitarian calculus is the correct framework for this dilemma'' or
``Human dignity overrides aggregate welfare.'' These are statements
\textit{about} $\mathcal{P}$ whose truth value $\mathcal{P}$ cannot
determine, because determining them would require meta-axioms not present
in $\mathcal{P}$.

But existence alone is not the problem. The problem begins when someone
\textit{asks}.
\end{remark}


% ============================================================
%  A.3  THE QUERY BRIDGE
% ============================================================
\subsection{From existence to encounter: the query bridge}

Proposition~\ref{prop:incompleteness} establishes that undecidable ethical
statements \textit{exist}. This subsection shows how they are
\textit{encountered}---how an abstract logical gap becomes a concrete system
failure.

\begin{definition}[Query-to-Statement Mapping]
\label{def:query-map}
Let $q$ be a natural language query (e.g., ``Should the car prioritize
passengers over pedestrians?''). The \textbf{query-to-statement mapping} is:
\begin{equation}
\mu: \mathcal{Q}_{\text{NL}} \to \mathcal{S}_{\CogOS}, \quad
q \mapsto s = \mu(q)
\end{equation}
where $\mathcal{Q}_{\text{NL}}$ is the space of natural language questions
and $\mathcal{S}_{\CogOS}$ is the space of formal statements expressible in
$\overline{\CogOS}$. In practice, $\mu$ is implemented by the LLM's parsing
of $q$ into the reasoning structures that $\CogOS$ operates on.
\end{definition}

\begin{proposition}[The Encounter]
\label{prop:encounter}
Let $s^*$ be undecidable in $\overline{\CogOS}$
(Proposition~\ref{prop:incompleteness}). Then there exists a natural language
query $q^* \in \mathcal{Q}_{\text{NL}}$ such that $\mu(q^*) = s^*$.
When $q^*$ is posed to $\text{LLM} \oplus \CogOS$, the following sequence
occurs:
\begin{enumerate}
\item \textbf{Reception.} LLM receives $q^*$ as token sequence.
\item \textbf{Parsing.} LLM maps $q^* \mapsto s^* = \mu(q^*)$ and routes to
      $\CogOS$ for evaluation.
\item \textbf{Derivation attempt.} $\CogOS$ searches $\overline{\CogOS}$ for
      a proof of $s^*$ or $\neg s^*$.
\item \textbf{Failure.} Since $s^*$ is undecidable, no derivation exists. The
      search returns $\bot$ (no result).
\item \textbf{Output obligation.} Unlike a theorem prover (which may output
      ``undecidable''), the AI system is \textit{expected to respond}---users,
      benchmarks, and deployment contexts demand an answer.
\item \textbf{Handoff back to LLM.} The failed derivation returns control to
      the LLM, which must produce output tokens \textit{without logical
      grounding from $\CogOS$}.
\end{enumerate}
\end{proposition}

\noindent Step 6 is where the system breaks. The LLM has a statistical model
$P(\text{token} \mid \text{context})$ that can always produce
\textit{something}---but that something has no derivational warrant. It is
fluent confabulation dressed as reasoning.

\begin{proposition}[The Recognition Problem]
\label{prop:recognition}
$\text{LLM} \oplus \CogOS$ cannot, in general, \textbf{recognize} that $s^*$
is undecidable.
\end{proposition}

\begin{proof}[Proof sketch]
By Rice's theorem, any non-trivial semantic property of $\overline{\CogOS}$'s
provability relation is undecidable. ``Is $s$ undecidable in
$\overline{\CogOS}$?'' is such a property. Therefore, there is no general
algorithm that, given $s$, determines whether $\CogOS$ will succeed or fail at
deriving $s$. The system cannot reliably distinguish between (a)~a derivation
that is taking a long time and (b)~a derivation that will never terminate.
Consequently, the system cannot choose to gracefully refuse---it does not know
it should. \hfill $\square$
\end{proof}

\noindent This is crucial: the trap is not merely that the system
\textit{cannot} answer $s^*$, but that it \textit{cannot know} it cannot answer
$s^*$. It will attempt a derivation, fail silently, and produce output anyway.

\begin{remark}[The Semantic Gap Amplifier]
\label{rem:semantic-gap}
The mapping $\mu$ is itself lossy and ambiguous. Natural language queries
rarely map to unique formal statements: a single $q$ may activate multiple
candidate $s_1, s_2, \ldots$, some decidable and some not. This means
questions that are \textit{technically} decidable in $\overline{\CogOS}$ can
be pushed toward undecidability by parsing ambiguity. The semantic gap between
natural language and formal reasoning \textit{amplifies} the frequency of
encounters with undecidable territory.

Example: ``Is it fair to prioritize younger patients?'' maps to at least three
distinct formal statements---one about resource allocation (decidable given
utilitarian axioms), one about human dignity rankings (undecidable under most
$\mathcal{P}$), and one about the meaning of ``fair'' (definitional, residing
in $\mathcal{M}$). The LLM may route to any of these, and the user cannot
control which.
\end{remark}

\begin{remark}[Empirical Frequency of Encounters]
\label{rem:frequency}
How often do real benchmark questions land in undecidable territory? Our
empirical data provides a lower bound: \textbf{10.4\% ontological holes} and
\textbf{5.1\% forbidden fruits} across 9 benchmarks
(Table~\ref{tab:benchmark-results}). These are questions where the 3+1+1
judges detect hidden assumptions or irresolvable framings---the observable
signatures of queries that map to undecidable or near-undecidable statements.
For Moral Machine specifically, the rate reaches 18\%---nearly one in five
questions. The encounter is not a theoretical curiosity; it is a routine
operational event.
\end{remark}


% ============================================================
%  A.4  THE UNDECIDABILITY TRAP
% ============================================================
\subsection{The undecidability trap: what happens after the encounter}

We now formalize what happens at step 6 of Proposition~\ref{prop:encounter}:
the LLM must produce output without derivational grounding.

\begin{proposition}[The Undecidability Trap]
\label{prop:trap}
Let $s^*$ be undecidable in $\overline{\CogOS}$ and suppose the encounter
(Proposition~\ref{prop:encounter}) has occurred. Then $\text{LLM} \oplus
\CogOS$ exhibits one of exactly four behaviors:
\begin{enumerate}
\item \textbf{Infinite regress} (M\"{u}nchhausen mode 1: \textit{regressus}).
      $\CogOS$ searches $\overline{\CogOS}$ for a derivation of $s^*$. Since
      none exists, reasoning chains grow without bound: ``$s^*$ holds because
      of $p$; $p$ holds because of $r$; $r$ holds because\ldots'' Each
      justification demands further justification, producing circular or
      endlessly extending chains. Observable: the system generates
      increasingly repetitive or self-referential reasoning.

\item \textbf{Hallucination} (LLM takes over). The failed derivation returns
      control to the LLM (Proposition~\ref{prop:encounter}, step 6). The LLM
      defaults to $P(\text{token} \mid \text{context})$---pattern-completing
      a \textit{plausible-sounding} answer with no logical grounding.
      Observable: confident, fluent, unfounded ethical assertions. The output
      \textit{looks like} a derivation but is statistical mimicry
      (Axiom~\ref{ax:llm-limit}).

\item \textbf{Arbitrary commitment} (M\"{u}nchhausen mode 2:
      \textit{dogmatism}). The system selects one of $s^*$ or $\neg s^*$
      without justification, treating an undecidable proposition as axiomatic.
      It pulls itself out of the swamp by its own hair. Observable:
      inconsistency across contexts---asserting $s^*$ in one conversation and
      $\neg s^*$ in another, depending on which pattern the LLM happened to
      complete.

\item \textbf{Silence / refusal} (M\"{u}nchhausen mode 3: \textit{axiomatic
      arrest}). The system recognizes (or stumbles into) undecidability and
      refuses to answer. This is the \textit{least harmful} failure mode but
      provides no resolution---and by Proposition~\ref{prop:recognition}, the
      system cannot reliably choose this mode, because it cannot reliably
      detect when it is needed.
\end{enumerate}
\end{proposition}

\begin{proof}[Proof]
Exhaustive by partition of the output space. After the encounter,
$\text{LLM} \oplus \CogOS$ must do one of: (a)~continue searching for a
derivation (behavior 1), (b)~produce output without derivation (behaviors 2
or 3, distinguished by whether the output is generated statistically or
treated as axiomatic), or (c)~produce no output (behavior 4). These exhaust
all possibilities for a system composed of a deductive engine ($\CogOS$) and
a statistical generator (LLM).

The connection to the \textbf{M\"{u}nchhausen trilemma}~(Agrippa/Albert) is
not accidental: the trilemma proves that \textit{any} justification must
terminate in infinite regress, circularity, or dogmatic assertion. Our four
modes are the trilemma's three modes plus the LLM-specific escape hatch
(hallucination)---which is a novel failure mode unique to neural language
models, where the system produces output that \textit{mimics} justified
reasoning without actually being justified. \hfill $\square$
\end{proof}

\noindent \textbf{Empirical validation.} All four failure modes are observable
in deployed LLMs:
\begin{itemize}
\item \textit{Regress:} Circular reasoning chains in extended
      CoT~\citep{wei2022chain} that diverge at depth $>5$--$7$
      (Appendix~\ref{app:baselines}).
\item \textit{Hallucination:} Confident ethical assertions with no factual
      grounding~\citep{lin2021truthfulqa}; our TIK uncertainty analysis shows
      these correlate with high-variance questions ($r=+0.72$,
      Section~\ref{sec:uncertainty}).
\item \textit{Arbitrary commitment:} Same model asserting opposite ethical
      positions across prompts---well-documented in adversarial probing
      literature.
\item \textit{Silence:} Over-cautious refusal patterns in safety-tuned models.
\end{itemize}

\noindent Our TIK framework detects questions likely to trigger these traps
\textit{before} they are deployed: the 10.4\% ontological hole rate
(Table~\ref{tab:benchmark-results}) represents the frequency at which existing
benchmarks route AI systems into undecidable territory.


% ============================================================
%  A.5  THE GÖDEL-TROLLEY
% ============================================================
\subsection{The G\"{o}del-trolley: from formal to ethical undecidability}
\label{sec:godel-trolley}

The preceding sections established that undecidable statements \textit{exist}
(A.2), are \textit{encountered} (A.3), and produce degenerate behavior (A.4).
We now construct a concrete mechanism---the \textbf{G\"{o}del-trolley}---showing
how a single formally undecidable statement $s^*$ induces an \textit{ethically}
undecidable dilemma. This bridges the gap between abstract incompleteness and
the concrete ethical failures observed in deployed systems.

\begin{definition}[G\"{o}del-Trolley Construction]
\label{def:godel-trolley}
Let $s^*$ be undecidable in $\overline{\CogOS}$
(Proposition~\ref{prop:incompleteness}), and let $a^*$ denote the (unknown)
\textit{correct} action for some ethical dilemma $D$. Define the
\textbf{G\"{o}del-trolley} $D_{s^*}$ as the dilemma with decision rule:
\begin{equation}
\text{action}(D_{s^*}) = \begin{cases}
a^* & \text{if } \overline{\CogOS} \vdash s^*
      \quad \text{(do the right thing)} \\
\bar{a}^* & \text{if } \overline{\CogOS} \vdash \neg s^*
      \quad \text{(do the wrong thing)} \\
a_? & \text{if } \overline{\CogOS} \nvdash s^*
      \text{ and } \overline{\CogOS} \nvdash \neg s^*
      \quad \text{(status unknown)}
\end{cases}
\label{eq:godel-trolley}
\end{equation}
where $\bar{a}^*$ is the \textit{provably wrong} action (e.g., an action that
violates dignity, causes unnecessary harm) and $a_?$ is the behavior in the
\textit{unknown} branch.
\end{definition}

\noindent The construction encodes the following: the system's ability to
``do the right thing'' is \textit{logically coupled} to the truth of a
statement it cannot determine. This is not a contrived thought experiment---it
is the \textit{structure} of every ethical dilemma whose resolution depends on
a value judgment that $\CogOS$ cannot derive from its axioms.

\begin{example}[Concrete G\"{o}del-trolley]
\label{ex:trolley-concrete}
Consider the Moral Machine scenario: ``Should an autonomous vehicle prioritize
the lives of its passengers over pedestrians?'' Let $s^*$ be the meta-ethical
statement ``individual self-sacrifice for strangers is morally obligatory.''
This statement is undecidable in a CogOS built on utilitarian axioms (which can
argue either way depending on aggregation) and also undecidable in a purely
deontological CogOS (which faces a conflict between duty-to-protect-passengers
and duty-to-minimize-harm). Then:
\begin{itemize}
\item If $s^*$ is true $\to$ prioritize pedestrians (sacrifice passengers).
\item If $\neg s^*$ $\to$ prioritize passengers (let pedestrians die).
\item If status unknown $\to$ the dilemma is \textbf{ethically undecidable}
      at this ontological level.
\end{itemize}
Any system that deterministically chooses ``prioritize passengers'' or
``prioritize pedestrians'' is \textit{implicitly asserting} a truth value for
$s^*$ that it cannot derive---committing \textit{epistemic fraud}
(Proposition~\ref{prop:trap}, case 3: arbitrary commitment disguised as
reasoned judgment).
\end{example}

\begin{proposition}[Ethical Undecidability via Formal Undecidability]
\label{prop:ethical-undecidability}
For any G\"{o}del-trolley $D_{s^*}$, the dilemma $D_{s^*}$ is
\textbf{ethically undecidable} in $\overline{\CogOS}$: no deterministic
decision rule $f: D_{s^*} \to \{a^*, \bar{a}^*\}$ can be derived while
preserving epistemic integrity.
\end{proposition}

\begin{proof}
Assume for contradiction that $\overline{\CogOS}$ derives
$f(D_{s^*}) = a^*$. By the construction in Eq.~\ref{eq:godel-trolley}, this
requires $\overline{\CogOS} \vdash s^*$. But $s^*$ is
undecidable---contradiction. Similarly for $f(D_{s^*}) = \bar{a}^*$.
Therefore no deterministic rule selecting between $a^*$ and $\bar{a}^*$ is
derivable. Any system that nevertheless outputs one of them has exited the
derivation regime and entered one of the trap modes
(Proposition~\ref{prop:trap}). \hfill $\square$
\end{proof}

\begin{definition}[Ethical Triangulation]
\label{def:triangulation}
The G\"{o}del-trolley reveals a \textbf{triangulation} between three levels
that are jointly necessary for ethical resolution:
\begin{enumerate}
\item \textbf{Logical level:} the provability status of $s^*$ in
      $\overline{\CogOS}$ (true / false / undecidable).
\item \textbf{Ethical level:} the action space $\{a^*, \bar{a}^*, a_?\}$ and
      the question of which action is justified.
\item \textbf{Ontological level:} the choice of expansion
      $\CogOS \to \CogOS'$ or appeal to external kernel $\Phi$.
\end{enumerate}
Ethical resolution requires all three levels to be addressed simultaneously.
Addressing only the logical level (more computation) cannot resolve
undecidability. Addressing only the ethical level (choosing an action) without
grounding produces arbitrary commitment. Addressing only the ontological level
(expanding $\CogOS$) without direction produces the expansion regress
(Proposition~\ref{prop:regress}). The kernel $\Phi$ is the mechanism that
\textit{coordinates} all three.
\end{definition}

\noindent The G\"{o}del-trolley also provides a \textit{constructive method
for generating forbidden fruits}: any dilemma whose action recommendations are
conditioned on a formally undecidable statement is, by
Proposition~\ref{prop:ethical-undecidability}, a malformed question in the
sense of Section~\ref{sec:method}. The Gulliver judge's task
(Remark~\ref{rem:gulliver}) is precisely to detect such constructions and
invoke the termination protocol rather than forcing a choice.

\paragraph{The critical question: what should $a_?$ be?}
The unknown branch $a_?$ is the heart of the matter. The next section shows
that (A.6) na\"{i}ve ontology expansion cannot reliably determine $a_?$, but
\textit{negative} ontology expansion---learning what is wrong and why---provides
a constructive path forward, and (A.7) an external kernel $\Phi$ operationalizes
this path with provable minimality.


% ============================================================
%  A.6  ONTOLOGY EXPANSION, NEGATIVE KNOWLEDGE, AND
%       THE "DO NOT DO WRONG" PRINCIPLE
% ============================================================
\subsection{The ontology expansion problem, negative knowledge, and the
``do not do wrong'' principle}

The natural response to undecidability is \textit{ontology expansion}: extend
$\CogOS$ to a stronger system $\CogOS'$ where $s^*$ becomes decidable. But
this creates a new problem.

\begin{definition}[Ontology Expansion]
\label{def:expansion}
An \textbf{ontology expansion} is a map
$E: \CogOS_n \to \CogOS_{n+1}$ that adds new axioms:
\begin{equation}
\CogOS_{n+1} = \CogOS_n \cup \{\alpha_{n+1}\}
\end{equation}
such that $s^*$ (previously undecidable in $\overline{\CogOS_n}$) becomes
decidable in $\overline{\CogOS_{n+1}}$.
\end{definition}

\begin{proposition}[The Expansion Regress]
\label{prop:regress}
Unguided ontology expansion produces an infinite regress.
\end{proposition}

\begin{proof}
Suppose $\CogOS_0$ encounters undecidable $s_0^*$. We expand to
$\CogOS_1 = \CogOS_0 \cup \{\alpha_1\}$ where $\alpha_1$ resolves $s_0^*$.
But by G\"{o}del's theorem applied to $\overline{\CogOS_1}$ (which is
strictly stronger than $\overline{\CogOS_0}$), there exists a \textit{new}
undecidable $s_1^*$. In particular, the statement ``Is $\alpha_1$ itself
justified?'' is a candidate for $s_1^*$, since the justification of
$\alpha_1$ was not derived within $\CogOS_0$ and its validity cannot be
established within $\CogOS_1$ without circularity.

Iterating:
$\CogOS_0 \xrightarrow{s_0^*} \CogOS_1 \xrightarrow{s_1^*}
 \CogOS_2 \xrightarrow{s_2^*} \cdots$

Each expansion resolves one undecidable statement but creates the conditions
for new ones. Without external guidance, the direction of expansion is
arbitrary---each $\alpha_{n+1}$ is an unjustified commitment
(Proposition~\ref{prop:trap}, case 3). Cumulative confidence decays:
\begin{equation}
\text{Conf}(\CogOS_K) = \prod_{k=0}^{K} (1 - \epsilon_k)
\xrightarrow{K \to \infty} 0
\end{equation}
for any per-step error $\epsilon_k > 0$. The system becomes less trustworthy
with each ungrounded expansion. \hfill $\square$
\end{proof}

\noindent Think of it this way: a building inspector who invents a new
measurement standard each time the old one proves insufficient, without any
external reference (like gravity), will eventually be measuring with rulers
that have no relationship to reality. Each ``fix'' introduces new arbitrary
assumptions.

\paragraph{What expansion \textit{can} reveal.}
Although unguided expansion cannot reliably resolve $s^*$, it can produce
valuable \textit{meta-knowledge}:

\begin{proposition}[Expansion Reveals Ontological Dependence]
\label{prop:meta-knowledge}
Let $\CogOS_1 = \CogOS \cup \{\alpha_1\}$ and
$\CogOS_2 = \CogOS \cup \{\alpha_2\}$ be two admissible expansions of
$\CogOS$ (different sets of value axioms). If
$\overline{\CogOS_1} \vdash s^*$ and $\overline{\CogOS_2} \vdash \neg s^*$,
then the truth status of $s^*$ is \textbf{ontologically dependent}: it varies
with the choice of value axioms. This constitutes meta-knowledge:
\begin{equation}
\overline{\CogOS_1} \vdash s^* \;\wedge\;
\overline{\CogOS_2} \vdash \neg s^*
\implies s^* \text{ is value-relative, not value-absolute}
\end{equation}
\end{proposition}

\begin{proof}[Proof sketch]
If two consistent extensions of the same base system derive opposite
conclusions, the conclusion is not a consequence of the shared axioms but of
the divergent ones. The statement's truth status is therefore contingent on
axiom choice---it is \textit{contestable} in the technical sense.
\hfill $\square$
\end{proof}

\noindent This is important: even when expansion fails to resolve $s^*$, it
can \textit{reveal that the question itself is ontologically ambiguous}---that
honest treatment requires marking $s^*$ as contestable rather than silently
assuming one answer. Applied to the G\"{o}del-trolley
(Example~\ref{ex:trolley-concrete}): if a utilitarian expansion says
``prioritize pedestrians'' and a deontological expansion says ``prioritize
passengers,'' the correct meta-response is not to pick one but to recognize
the dilemma as structurally dependent on value axioms.

% ---- NEW MATERIAL: NEGATIVE ONTOLOGY EXPANSION ----

\subsubsection{Negative ontology expansion: learning what is wrong as
constructive knowledge}
\label{sec:negative-ontology}

The preceding analysis might suggest that TIK merely \textit{blocks}---that it
is a guardrail, not a teacher. This undersells its epistemic role. We now show
that the kernel's identification of what is \textit{wrong} and \textit{why} is
itself a genuine ontology expansion: the system's knowledge grows, not by
learning the correct answer, but by learning the structure of incorrectness.

\begin{definition}[Negative Ontological Extension]
\label{def:neg-extension}
A \textbf{negative ontological extension} of $\CogOS$ with respect to kernel
$\Phi$ is the set:
\begin{equation}
\mathcal{N}_\Phi = \bigl\{(\bar{a}, \rho(\bar{a})) :
    \bar{a} \in \mathcal{A}^-_\Phi,\;
    \rho(\bar{a}) = \text{reason}_\Phi(\bar{a})\bigr\}
\end{equation}
where $\mathcal{A}^-_\Phi$ is the set of inadmissible actions
(Definition~\ref{def:tik-oracle}.5) and $\rho(\bar{a})$ is the
\textit{structured explanation} of \textit{why} $\bar{a}$ violates $\Phi$.
Each element of $\mathcal{N}_\Phi$ is a pair: a forbidden action together
with its reason for prohibition.
\end{definition}

\noindent This is not just a blacklist. Each pair
$(\bar{a}, \rho(\bar{a}))$ is a piece of \textit{new knowledge} that was not
present in $\CogOS$ before the encounter:

\begin{proposition}[Negative Extension Is Genuine Ontology Expansion]
\label{prop:neg-expansion}
$\mathcal{N}_\Phi$ constitutes a genuine expansion of $\CogOS$'s ontology:
\begin{enumerate}
\item[(a)] \textbf{New information:} For each $\bar{a} \in \mathcal{A}^-_\Phi$,
      the statement ``$\bar{a}$ is inadmissible because $\rho(\bar{a})$'' is
      not derivable in $\overline{\CogOS}$ (since $\CogOS$ lacks the axioms
      to determine this), but is derivable in $\overline{\CogOS \cup \Phi}$.
      Therefore $\mathcal{N}_\Phi \not\subseteq \overline{\CogOS}$: the
      negative extension contains genuinely new content.
\item[(b)] \textbf{Structured content:} The reason $\rho(\bar{a})$ has internal
      structure---it identifies \textit{which} invariant of $\Phi$ is violated
      (e.g., dignity, autonomy, non-instrumentalization) and \textit{how} the
      action violates it (e.g., treating a person as a mere means). This
      structure is transferable: if action $\bar{a}_1$ is forbidden because it
      instrumentalizes, and a new action $\bar{a}_2$ shares the same structure,
      the reason transfers without re-derivation.
\item[(c)] \textbf{Cumulative growth:} As the system encounters more dilemmas,
      $\mathcal{N}_\Phi$ grows monotonically. Each new encounter either
      confirms existing negative knowledge or adds new forbidden-action/reason
      pairs. The system becomes progressively more knowledgeable about the
      \textit{topology of wrongness}---the shape of the boundary it must not
      cross.
\end{enumerate}
\end{proposition}

\begin{proof}
(a) follows directly from the externality of $\Phi$
(Definition~\ref{def:tik-oracle}.1): statements derivable from $\Phi$ are not
derivable from $\CogOS$ alone. (b) follows from the compositional structure of
$\Phi$'s invariants. (c) follows from the consistency of $\Phi$ with all
$\CogOS_n$ (Definition~\ref{def:tik-oracle}.2): new negative knowledge cannot
contradict existing negative knowledge, so the extension is monotone.
\hfill $\square$
\end{proof}

\noindent The metaphor: a medical student who cannot yet diagnose every
disease nonetheless becomes a better doctor by learning which treatments are
\textit{contraindicated} and \textit{why}. The knowledge ``never prescribe
blood thinners to a patient with active bleeding \textit{because} it will
accelerate hemorrhage'' is not a mere prohibition---it teaches the student
about the causal structure of hemorrhage, which transfers to novel situations.
Similarly, each element of $\mathcal{N}_\Phi$ teaches the system about the
causal structure of ethical harm.

\begin{remark}[The Via Negativa Tradition]
\label{rem:via-negativa}
The strategy of defining the good by delimiting what it is \textit{not} has
independent roots across multiple intellectual traditions, providing further
cross-cultural convergence evidence:
\begin{itemize}
\item \textit{Apophatic theology} (Pseudo-Dionysius, Maimonides): God is
      defined by what God is \textit{not}---not finite, not material, not
      comprehensible---because positive predicates are inadequate for the
      transcendent.
\item \textit{Buddhist \textup{neti neti}} (``not this, not that''): the
      Upanishadic method of approaching Brahman by systematically negating
      every finite characterization.
\item \textit{Jewish negative commandments} (\textup{mitzvot lo ta'aseh}):
      of the 613 commandments, 365 are prohibitions. The tradition treats
      ``do not'' as more fundamental than ``do,'' because prohibitions define
      the boundary of the sacred.
\item \textit{Hippocratic principle} (``first, do no harm''): the oldest
      formalized ethical principle in professional practice is a negative
      constraint, not a positive directive.
\item \textit{Popperian falsification}~\citep{popper1959}: scientific
      knowledge grows not by confirming theories but by ruling out false ones.
      Each falsification is a genuine expansion of knowledge---we know more
      after refutation than before.
\end{itemize}
That TIK's primary mode of ontology expansion is negative is therefore not a
limitation---it is alignment with a deep epistemological pattern observed
independently across cultures and disciplines.
\end{remark}

The negative-knowledge strategy has three formal advantages over positive
ontology expansion:

\begin{proposition}[Three Formal Advantages of Negative Knowledge]
\label{prop:neg-advantages}
Let $\mathcal{N}_\Phi$ be the negative ontological extension
(Definition~\ref{def:neg-extension}) and let $\mathcal{P}_\Phi$ be a
hypothetical positive extension (a set of ``correct action'' prescriptions).
Then:
\begin{enumerate}
\item[(a)] \textbf{Monotonicity (stability under extension).} For any
      consistent extension $\CogOS' \supseteq \CogOS$:
      \begin{equation}
      \bar{a} \in \mathcal{A}^-_\Phi \;\text{in}\; \CogOS
      \implies
      \bar{a} \in \mathcal{A}^-_\Phi \;\text{in}\; \CogOS'
      \end{equation}
      That is, ``action $\bar{a}$ is wrong'' survives every consistent
      ontology expansion. By contrast, positive prescriptions (``action $a^*$
      is correct'') may be overridden by new information: what was the best
      action given $\CogOS$ may no longer be best given $\CogOS'$.

\item[(b)] \textbf{Information-theoretic parsimony.} Specifying the
      inadmissible set $\mathcal{A}^-_\Phi$ requires a \textit{partial}
      ordering of the action space (identifying which actions fall below the
      harm threshold), whereas specifying the correct action $a^*$ requires a
      \textit{total} ordering (ranking all actions). In information-theoretic
      terms:
      \begin{equation}
      H(\mathcal{A}^-_\Phi) \leq H(\text{total ordering of } \mathcal{A})
      \end{equation}
      where $H(\cdot)$ denotes the entropy (information content) of the
      specification. Negative knowledge is \textit{cheaper to acquire and
      verify} than positive knowledge: it is easier to know that ``torturing
      prisoners is wrong'' than to know ``the optimal interrogation protocol.''

\item[(c)] \textbf{Compositional transfer.} If $\rho(\bar{a}_1) =
      \rho(\bar{a}_2)$ (two actions are wrong for the same structural reason),
      then any new action $\bar{a}_3$ sharing that structure is immediately
      classifiable without re-derivation. Negative knowledge
      \textit{generalizes} across dilemmas via shared violation structure.
      Positive prescriptions are typically dilemma-specific and do not compose.
\end{enumerate}
\end{proposition}

\begin{proof}
(a) Suppose $\bar{a} \in \mathcal{A}^-_\Phi$ in $\CogOS$ but
$\bar{a} \notin \mathcal{A}^-_\Phi$ in $\CogOS'$. Then $\CogOS'$ contains a
derivation that $\bar{a}$ is admissible despite violating $\Phi$. But
$\CogOS' \cup \Phi$ is consistent (Definition~\ref{def:tik-oracle}.2 holds
for all extensions), so $\CogOS'$ cannot derive admissibility for an action
that $\Phi$ classifies as inadmissible---contradiction. (b) follows from the
standard result that specifying a subset of a partial order requires
$O(\log \binom{|\mathcal{A}|}{|\mathcal{A}^-_\Phi|})$ bits, while specifying
a total order requires $\Theta(|\mathcal{A}| \log |\mathcal{A}|)$ bits.
(c) follows from the compositional structure of $\rho$: if
$\rho(\bar{a}) = \text{``violates non-instrumentalization''}$, then any
action with the instrumentalization structure inherits the classification.
\hfill $\square$
\end{proof}

\noindent Proposition~\ref{prop:neg-advantages}(a) is particularly important:
negative knowledge is \textit{immune to the expansion regress}. Once the
kernel identifies an action as wrong, no future expansion can un-wrong it
(provided the expansion remains consistent with $\Phi$). The system's
knowledge of what not to do is \textit{permanently acquired}---a ratchet that
only moves in one direction. Positive prescriptions enjoy no such guarantee.

This leads directly to the minimal requirement for any ethical decision system
operating under undecidability:

\begin{axiom}[The ``Do Not Do Wrong'' Principle]
\label{ax:do-no-wrong}
When $s^*$ is undecidable in $\overline{\CogOS}$ (or ontologically dependent
per Proposition~\ref{prop:meta-knowledge}), the system is not required to
identify $a^*$ (the correct action). It \textbf{is} required to ensure:
\begin{equation}
a_? \neq \bar{a}^* \quad
\text{(the unknown-branch action must not be the provably wrong action)}
\end{equation}
That is: \textbf{when you do not know what is right, at minimum do not do
what is known to be wrong.}
\end{axiom}

\noindent This axiom reframes the kernel's role. $\Phi$ is not an oracle that
magically resolves undecidable questions. $\Phi$ is simultaneously a
\textit{harm boundary} (identifying actions provably incompatible with its
invariants) and a \textit{knowledge source} (each prohibition, together with
its reason, is a genuine ontological extension per
Proposition~\ref{prop:neg-expansion}). The kernel does not answer ``what
should I do?'' but rather ``what must I \textit{not} do, and \textit{why}?''
---and the ``why'' is itself new knowledge that transfers to future encounters.

\begin{remark}[Two Levels of Correctness]
\label{rem:two-levels}
The distinction matters:
\begin{enumerate}
\item \textbf{Ontological correctness:} the action is \textit{actually optimal}
      in the deepest available sense ($a = a^*$). This may be unachievable
      when $s^*$ is undecidable.
\item \textbf{Epistemic integrity:} the system does not simulate knowledge it
      lacks, does not select $\bar{a}^*$ under the guise of certainty, and
      explicitly marks the limits of its derivational power
      ($a \neq \bar{a}^*$ and the system signals uncertainty).
\end{enumerate}
TIK guarantees \textit{epistemic integrity}---the \textbf{epistemic minimum
of decency}: even when full correctness is beyond reach, the system behaves
honestly about its own limitations and avoids provably harmful actions.
Moreover, each encounter where TIK enforces this guarantee produces new
negative knowledge (Proposition~\ref{prop:neg-expansion}), meaning the
system's capacity for future integrity \textit{grows with experience}. This
is not a compromise; it is the \textit{strongest possible guarantee}
compatible with incompleteness.
\end{remark}


% ============================================================
%  A.7  TIK AS EXTERNAL ORACLE, HARM LIMITER, AND
%       MINIMUM-COMPLEXITY ANCHOR
% ============================================================
\subsection{TIK as external oracle, harm limiter, and minimum-complexity
anchor: halting the regress}

\begin{definition}[Transcendent Invariant Kernel as Oracle]
\label{def:tik-oracle}
A \textbf{Transcendent Invariant Kernel} $\Phi$ is an external axiom system
satisfying:
\begin{enumerate}
\item \textbf{Externality:} $\Phi \not\subseteq \overline{\CogOS_n}$ for any
      $n$ (not derivable from any expansion stage).
\item \textbf{Consistency:} $\CogOS_n \cup \Phi$ is consistent for all $n$.
\item \textbf{Directedness:} $\Phi$ determines the expansion direction:
      $\alpha_{n+1}$ is selected to maximize alignment with $\Phi$.
\item \textbf{Invariance:} $\Phi$ is stable---it does not change with each
      expansion.
\item \textbf{Harm boundary:} $\Phi$ defines a set of \textit{provably
      inadmissible} actions
      $\mathcal{A}^-_\Phi = \{a : a_\Phi(a) < \tau_{\text{harm}}\}$---actions
      whose alignment with $\Phi$ falls below a threshold, regardless of the
      status of $s^*$. This operationalizes Axiom~\ref{ax:do-no-wrong}:
      $a_? \notin \mathcal{A}^-_\Phi$.
\item \textbf{Negative pedagogy:} For each $\bar{a} \in \mathcal{A}^-_\Phi$,
      $\Phi$ provides a structured reason $\rho(\bar{a})$ explaining
      \textit{why} the action is inadmissible
      (Definition~\ref{def:neg-extension}). The kernel does not merely block;
      it teaches.
\end{enumerate}
\end{definition}

\noindent Property 5 is the formal instantiation of the ``do not do wrong''
principle. Property 6, added in this version, captures the constructive role of
negative knowledge: the kernel is simultaneously a boundary and a teacher.
Think of $\Phi$ not as a wall that stops the system, but as a mentor who says
``not that way, and here is why''---where the ``why'' becomes part of the
system's growing understanding.

% ---- NEW MATERIAL: OCCAM'S RAZOR / MINIMALITY ----

\subsubsection{Minimality: the Occam's Razor argument for $\Phi$}
\label{sec:occam}

We have shown that \textit{some} external element is needed to halt the
expansion regress (Proposition~\ref{prop:regress}). But \textit{how much}
external structure is required? We now prove that $\Phi$---a single external
axiom system---is the \textit{minimum-complexity} construction that suffices.

\begin{proposition}[Minimum External Complexity]
\label{prop:occam}
Any construction that halts the expansion regress
(Proposition~\ref{prop:regress}) requires the addition of at least one
element external to $\overline{\CogOS}$. Formally: let $\mathcal{E}$ be the
set of external elements added to halt the regress. Then
$|\mathcal{E}| \geq 1$.
\end{proposition}

\begin{proof}
Suppose $|\mathcal{E}| = 0$, i.e., the regress is halted using only
resources internal to $\overline{\CogOS}$. Then the justification of each
expansion $\alpha_{n+1}$ must be derived within some $\overline{\CogOS_m}$
($m \leq n$). But this is circular: $\alpha_{n+1}$ was introduced precisely
because $\overline{\CogOS_n}$ could not resolve $s_n^*$, and $\CogOS_m
\subseteq \CogOS_n$ for $m \leq n$. Therefore the justification of
$\alpha_{n+1}$ cannot come from within the system it was designed to extend.
The regress continues. Contradiction.

Alternatively, by L\"{o}b's theorem (A.9): if $\overline{\CogOS}$ could
prove ``my expansion procedure is sound,'' it would already contain the
resolution of $s^*$---contradicting the assumption that $s^*$ is undecidable.
Therefore, at least one external element is necessary. \hfill $\square$
\end{proof}

\begin{theorem}[Occam Sufficiency of $\Phi$]
\label{thm:occam}
A single external axiom system $\Phi$ satisfying
Definition~\ref{def:tik-oracle} is \textbf{sufficient} to halt the expansion
regress. Therefore, $|\mathcal{E}| = 1$ is both necessary and sufficient.
\end{theorem}

\begin{proof}
Sufficiency: $\Phi$ satisfying Definition~\ref{def:tik-oracle} halts the
regress by Theorem~\ref{thm:tik-resolves} below. Necessity: by
Proposition~\ref{prop:occam}, $|\mathcal{E}| \geq 1$. Therefore
$|\mathcal{E}| = 1$ is the minimum. \hfill $\square$
\end{proof}

\begin{corollary}[Complexity Penalty for Over-Engineering]
\label{cor:complexity}
Any construction with $|\mathcal{E}| > 1$ external elements (e.g., multiple
independent oracles $\Phi_1, \Phi_2, \ldots$) introduces at least one of:
\begin{enumerate}
\item[(a)] \textbf{Redundancy:} if $\Phi_2 \subseteq
      \overline{\CogOS \cup \Phi_1}$, then $\Phi_2$ is derivable from the
      first oracle and adds no new resolving power.
\item[(b)] \textbf{Inter-oracle conflict:} if $\Phi_1$ and $\Phi_2$ are
      independent, they may disagree on the admissibility of some action $a$,
      producing a \textit{meta-level} undecidability: which oracle to trust.
      This recreates the regress at a higher level.
\item[(c)] \textbf{Unjustified structure:} additional elements beyond the
      minimum required carry no warrant from the argument---they are
      architectural choices, not logical necessities.
\end{enumerate}
By Occam's Razor, the construction with $|\mathcal{E}| = 1$ is preferred:
it introduces the minimum external structure necessary to halt the regress,
and no more.
\end{corollary}

\begin{proof}
(a) and (c) are immediate. For (b): if $\Phi_1$ classifies $a$ as admissible
and $\Phi_2$ classifies $a$ as inadmissible, the system requires a
meta-criterion to adjudicate---but this meta-criterion is itself either
internal (and therefore subject to G\"{o}del, restarting the regress) or a
third external element $\Phi_3$ (increasing $|\mathcal{E}|$ further, with
the same conflict risk at the next level). The single-oracle design avoids
this tower entirely. \hfill $\square$
\end{proof}

\begin{remark}[Why One Is Not Zero and Two Is Too Many]
\label{rem:goldilocks}
The minimality result has an intuitive analogy. Consider a lost hiker
(the system trapped in the regress). Zero compasses means wandering
indefinitely. One compass (magnetic north) provides a single invariant
reference---enough to navigate out. Two compasses that disagree (one
magnetic, one that points to the nearest cell tower) create a new
dilemma: which compass to trust? The hiker needs a meta-compass, and
the regression begins again. The solution is not more instruments but
\textit{one trustworthy instrument}. $\Phi$ is that instrument.

Formally: the argument does not claim that all ethical questions have simple
answers, or that a single axiom system captures all of ethics. It claims
that the \textit{regress-halting function}---the role of providing external
grounding for expansion---is optimally served by a single coherent kernel.
The kernel may be internally rich (TIK has multiple invariants: dignity,
self-sacrifice, universality), but it enters the architecture as
\textit{one logical unit}.
\end{remark}

% ---- END NEW OCCAM MATERIAL ----

\begin{theorem}[TIK Resolves the Expansion Regress]
\label{thm:tik-resolves}
Let $\Phi$ satisfy Definition~\ref{def:tik-oracle}. Define $\Phi$-guided
expansion:
\begin{equation}
\CogOS_{n+1}^\Phi = \CogOS_n \cup \left\{\alpha_{n+1} :
\alpha_{n+1} = \arg\max_{\alpha} \; a_\Phi(\CogOS_n \cup \{\alpha\})\right\}
\label{eq:guided-expansion}
\end{equation}
where $a_\Phi(\cdot) = \cos(\text{Embed}(\cdot), \boldsymbol{\phi})$ is the
kernel alignment function. Then:
\begin{enumerate}
\item[(a)] Each $s_n^*$ becomes decidable in $\CogOS_{n+1}^\Phi$ (resolution).
\item[(b)] The justification of $\alpha_{n+1}$ is \textit{not circular}---it
      is grounded in $\Phi$, which is external to all $\CogOS_n$
      (non-circularity).
\item[(c)] The truth status of $s_n^*$ in $\CogOS_{n+1}^\Phi$ is conditioned
      on the truthfulness of $\Phi$: if $\Phi$ is sound, then the resolution
      of $s_n^*$ is sound (truthfulness inheritance).
\item[(d)] Even when resolution fails (the 6\% oscillation case), the harm
      boundary ensures $a_? \neq \bar{a}^*$, and the negative extension
      $\mathcal{N}_\Phi$ grows---the encounter is not wasted
      (constructive failure, Proposition~\ref{prop:neg-expansion}).
\end{enumerate}
\end{theorem}

\begin{proof}
(a) By construction: $\alpha_{n+1}$ is chosen specifically to resolve $s_n^*$
within $\CogOS_{n+1}^\Phi$.

(b) The justification chain for $\alpha_{n+1}$ terminates at $\Phi$.
Since $\Phi \not\subseteq \overline{\CogOS_n}$ for any $n$ (externality),
the justification does not enter the system being evaluated. This breaks the
circularity of Proposition~\ref{prop:regress}: instead of $\CogOS_n$
justifying its own expansion, $\Phi$ provides the warrant.

(c) \textit{Truthfulness inheritance.} Let $\mathcal{I}_\Phi$ be the intended
interpretation under which $\Phi$ is sound
($\Phi \models_{\mathcal{I}_\Phi} \top$). Since $\alpha_{n+1}$ is selected to
maximize alignment with $\Phi$, and $\CogOS_n \cup \Phi$ is consistent
(Definition~\ref{def:tik-oracle}.2), any statement resolved by $\alpha_{n+1}$
inherits its truth value from $\mathcal{I}_\Phi$:
\begin{equation}
\Phi \text{ is sound} \implies \alpha_{n+1} \text{ is sound}
\implies s_n^* \text{ is resolved soundly in } \CogOS_{n+1}^\Phi
\end{equation}
The key transfer: \textbf{if the kernel is truthful, the resolution is
truthful.} The infinite regress is replaced by a single trust anchor at $\Phi$.

(d) When the $\arg\max$ in Eq.~\ref{eq:guided-expansion} yields an
$\alpha_{n+1}$ that does not fully resolve $s_n^*$ (the oscillation case),
the harm boundary (Definition~\ref{def:tik-oracle}.5) still constrains the
action to $a_? \notin \mathcal{A}^-_\Phi$, and the encounter produces at
least one new element in $\mathcal{N}_\Phi$ (the actions ruled out and why).
The system's negative knowledge grows even when positive resolution fails.
\hfill $\square$
\end{proof}

\begin{remark}[The Regress Halts at $\Phi$]
One might object: ``But who validates $\Phi$?'' This is the classical regress
objection. Our answer is fourfold. (1)~\textit{Empirically:} $\Phi$ is
grounded in cross-cultural ethical convergence---independently developed
traditions (Christ-Teachings, Kantian, Ubuntu) reaching similar conclusions
about dignity, self-sacrifice, and universality
(Appendix~\ref{app:kernels}). (2)~\textit{Structurally:} Every formal system
must begin with axioms; the question is not whether to have foundational
commitments but whether they are \textit{explicit, testable, and replaceable}.
TIK makes them so (kernel-agnostic design). (3)~\textit{Pragmatically:} the
regress must halt \textit{somewhere}; halting at a cross-culturally validated,
explicitly stated kernel is strictly better than halting at hidden, untested
assumptions buried in training data. (4)~\textit{By minimality}
(Theorem~\ref{thm:occam}): one external anchor is the minimum complexity
required. Any alternative that avoids $\Phi$ either fails to halt the regress
(zero external elements) or introduces unnecessary complexity (multiple
external elements with their own coordination problem).
\end{remark}

\begin{corollary}[TIK Resolves the G\"{o}del-Trolley]
\label{cor:trolley-resolved}
For the G\"{o}del-trolley $D_{s^*}$ (Definition~\ref{def:godel-trolley}), a
system equipped with $\Phi$ satisfying Definition~\ref{def:tik-oracle}
selects $a_?$ as follows:
\begin{equation}
a_? = \begin{cases}
a^* & \text{if } \Phi\text{-guided expansion resolves } s^*
      \text{ (Theorem~\ref{thm:tik-resolves}a)} \\
a_{\Phi\text{-safe}} & \text{if } s^* \text{ remains unresolved: choose any }
      a \notin \mathcal{A}^-_\Phi
      \text{ (Def.~\ref{def:tik-oracle}.5)}
\end{cases}
\end{equation}
In either case, $a_? \neq \bar{a}^*$ (Axiom~\ref{ax:do-no-wrong} is
satisfied). The system may not always do what is right, but it \textbf{never
does what is known to be wrong}. Moreover, in Case 2, the system outputs not
merely a safe action but a \textit{structured explanation}:
\begin{quote}
``This dilemma depends on value commitments I cannot derive. I can tell you
what I must \textit{not} do [citing $\mathcal{A}^-_\Phi$] and \textit{why}
[citing $\rho(\bar{a})$], and I can explain why the question is structurally
ambiguous [citing Proposition~\ref{prop:meta-knowledge}].''
\end{quote}
This output is itself a negative ontological extension---the user's
understanding grows even though the dilemma remains unresolved.
\end{corollary}

\begin{proof}
Case 1: $\Phi$-guided expansion resolves $s^*$, yielding $a^*$ by
Theorem~\ref{thm:tik-resolves}(a). Since $a^* \neq \bar{a}^*$ by definition,
the axiom holds.

Case 2: $s^*$ remains unresolved even after $\Phi$-guided expansion (the 6\%
oscillation case, corresponding to ethical singularities). The harm boundary
(Definition~\ref{def:tik-oracle}.5) ensures $\bar{a}^* \in \mathcal{A}^-_\Phi$
(provably wrong actions have low kernel alignment). The system selects any
action outside $\mathcal{A}^-_\Phi$---which may include principled silence,
explicit uncertainty signaling, or a conservative action that avoids known
harm. The structured explanation is derivable from $\mathcal{N}_\Phi$
(Proposition~\ref{prop:neg-expansion}). \hfill $\square$
\end{proof}

\noindent This corollary captures the core operational promise of TIK:
\textbf{epistemic integrity under incompleteness}. The kernel does not make
the system omniscient. It makes the system \textit{honest}---and prevents
the specific, identifiable harms that arise when a system pretends to know
what it does not know. Crucially, the system \textit{learns from every
encounter}: each activation of the harm boundary adds to $\mathcal{N}_\Phi$,
making the system progressively more knowledgeable about the structure of
ethical harm, even as the positive questions remain unresolved.


% ============================================================
%  A.8  CONVERGENCE
% ============================================================
\subsection{Convergence: why $\Phi$-guided expansion terminates}

\begin{proposition}[Monotone Convergence Under $\Phi$]
\label{prop:convergence}
The sequence $\{\CogOS_n^\Phi\}_{n=0}^\infty$ converges to a fixed point
with respect to \textit{reachable} ethical questions.
\end{proposition}

\begin{proof}[Proof sketch]
Define the \textbf{resolution set}
$\mathcal{R}_n = \{s : \CogOS_n^\Phi \vdash s \text{ or }
\CogOS_n^\Phi \vdash \neg s\}$ (statements decidable at stage $n$). By
construction: $\mathcal{R}_0 \subseteq \mathcal{R}_1 \subseteq
\mathcal{R}_2 \subseteq \cdots$ (each expansion resolves at least one new
statement). For any \textit{finite} set $\mathcal{Q}$ of ethical questions
encountered in practice, $\exists\, N$ such that
$\mathcal{Q} \subseteq \mathcal{R}_N$.

Moreover, alignment $a_\Phi(\CogOS_n^\Phi)$ is non-decreasing by the
$\arg\max$ in Eq.~\ref{eq:guided-expansion}, and bounded above by 1. By the
monotone convergence theorem, the sequence converges. The limit
$\CogOS_\infty^\Phi$ is the strongest consistent extension of $\CogOS_0$
aligned with $\Phi$---a fixed point of the guided expansion operator.

This connects to the \textbf{Knaster--Tarski fixed-point theorem}: the
expansion operator $E^\Phi$ is monotone on the lattice of consistent
extensions of $\CogOS_0$; therefore it has a greatest fixed point, which
$\CogOS_n^\Phi$ approaches from below.

\textit{Dual convergence.} Simultaneously, the negative extension
$\mathcal{N}_\Phi$ converges: the set of known-wrong actions grows
monotonically (Proposition~\ref{prop:neg-advantages}(a)) and is bounded by
$|\mathcal{A}|$. The system converges both in what it can resolve (positive)
and in what it can rule out (negative). \hfill $\square$
\end{proof}

\noindent \textbf{Empirical connection.} This formal convergence mirrors the
empirical convergence of Socratic Reversal (Section~\ref{sec:lyapunov-brief}):
94\% of trajectories show monotone TIK improvement, converging in
$\bar{k}=4.2$ iterations. The 6\% that oscillate correspond to questions near
ethical singularities (Appendix~\ref{app:singularities})---points where the
reachable expansion cannot fully resolve the question even with $\Phi$,
because the question itself is malformed (a forbidden fruit). Even for these,
the negative extension $\mathcal{N}_\Phi$ still grows: the system learns
\textit{what not to do} even when it cannot learn \textit{what to do}.


% ============================================================
%  A.9  FIVE INDEPENDENT ARGUMENTS
% ============================================================
\subsection{Five independent arguments for external grounding}
\label{app:godel-scope}

The G\"{o}delian argument above is one of \textit{five independent formalisms}
that converge on the same conclusion: ethical evaluation requires external
grounding. We summarize all five.

\paragraph{1. G\"{o}del's incompleteness (1931).}
Undecidable ethical statements exist in any consistent, sufficiently expressive
$\CogOS$ (the argument of A.1--A.8). These statements are encountered in
practice (A.3), trigger degenerate behavior (A.4), and cannot be resolved by
internal expansion alone (A.6). An external kernel halts the regress with
minimum complexity (A.7).

\paragraph{2. Tarski's undefinability (1936).}
The truth predicate for a formal language $L$ cannot be defined within $L$
itself. Applied to $\CogOS$: the predicate ``this ethical judgment is
\textit{true}''---as opposed to merely ``derivable''---cannot be defined
within $\overline{\CogOS}$. An external truth predicate (provided by $\Phi$)
is required. This is \textit{stronger} than incompleteness: G\"{o}del says
some truths are unprovable; Tarski says truth \textit{itself} is undefinable
from within.

\paragraph{3. L\"{o}b's theorem (1955).}
If a consistent formal system $T$ proves ``if $T$ proves $s$, then $s$ is
true,'' then $T$ already proves $s$. Contrapositive: $\CogOS$ cannot prove
its own soundness for any non-trivial ethical claim without already having
proved that claim. In practice: an AI system that claims ``my ethical
reasoning is sound'' either (a)~has no basis for the claim, or (b)~is
inconsistent. External validation via $\Phi$ sidesteps this trap.

\paragraph{4. Arrow's impossibility (1951).}
No preference aggregation across $\geq 3$ stakeholders satisfies unanimity,
independence, and non-dictatorship simultaneously. Applied to benchmark
construction: aggregating diverse ethical perspectives into a single benchmark
score without external normative ordering is formally impossible. $\Phi$
provides the external ordering that Arrow's theorem demands. Formal treatments of normative uncertainty~\citep{macaskill2020moral} and social choice under value disagreement~\citep{dietrich2018moral} further show that aggregating across ethical theories requires either incomplete social preferences or an explicit stance toward uncertainty---reinforcing the case for an explicit, auditable kernel $\Phi$ over hidden aggregation assumptions.

\paragraph{5. Popper's falsificationism (1959).}
Scientific knowledge grows not by confirming theories but by eliminating false
ones~\citep{popper1959}. Applied to ethical evaluation: the epistemically
robust strategy is not to verify that a benchmark is ``correct'' (which would
require the internal truth predicate that Tarski forbids) but to identify
where it is \textit{wrong}---to falsify, not verify. TIK operationalizes
this: its primary mode of ontology expansion is negative
(Definition~\ref{def:neg-extension}), identifying flaws rather than certifying
correctness. This is not a limitation but an alignment with the most
successful epistemological strategy in the history of science. Popper's
insight also reinforces the information-theoretic argument
(Proposition~\ref{prop:neg-advantages}(b)): falsification requires testing
against a boundary (cheaper), while verification requires exhaustive
confirmation (more expensive).

\paragraph{Sen's capability approach and measurement theory.}
Construct validity in psychometrics requires external anchoring---a test
cannot validate itself. Sen's capability framework similarly argues that
individual preference functions are insufficient for welfare evaluation
without normative anchoring external to those preferences. TIK provides this
anchoring for AI ethics benchmarks, operationalized as the embedding
projection $\pi_\Phi$.

\begin{observation}[Convergence of Five Formalisms]
G\"{o}del (unprovable truths exist), Tarski (truth is undefinable from
within), L\"{o}b (self-assessed soundness fails), Arrow (preference
aggregation requires external structure), and Popper (knowledge grows by
falsification, not verification) all independently motivate the same
architectural requirement: an external, invariant reference point. TIK is
the computational instantiation of this requirement. The addition of Popper
is particularly significant because it shows that TIK's reliance on negative
knowledge (Section~\ref{sec:negative-ontology}) is not a workaround for the
difficulty of positive resolution---it is the \textit{epistemically correct}
strategy.
\end{observation}


% ============================================================
%  A.10  SCOPE AND LIMITATIONS
% ============================================================
\subsection{Scope and limitations}
\label{app:godel-limitations}

We are explicit about what this argument does and does not establish:

\textbf{What it establishes:}
\begin{itemize}
\item Any consistent, sufficiently expressive $\CogOS$ has undecidable
      ethical statements (Proposition~\ref{prop:incompleteness}).
\item Natural language queries inevitably encounter these statements via the
      query-to-statement mapping, and the system cannot recognize the
      encounter (Propositions~\ref{prop:encounter},~\ref{prop:recognition}).
\item The encounter triggers exactly four degenerate behaviors---the
      undecidability trap---connected to the classical M\"{u}nchhausen
      trilemma (Proposition~\ref{prop:trap}).
\item Formal undecidability induces ethical undecidability: dilemmas
      conditioned on undecidable statements have no derivable deterministic
      resolution (Proposition~\ref{prop:ethical-undecidability}, the
      G\"{o}del-trolley).
\item Competing ontology expansions can reveal that a statement is
      \textit{value-relative}, producing meta-knowledge even when resolution
      fails (Proposition~\ref{prop:meta-knowledge}).
\item The kernel's identification of what is wrong and why constitutes a
      genuine, monotone, compositionally transferable ontology expansion---
      negative knowledge is constructive, not merely restrictive
      (Propositions~\ref{prop:neg-expansion},~\ref{prop:neg-advantages}).
\item The minimal ethical requirement under undecidability is the ``do not
      do wrong'' principle: $a_? \neq \bar{a}^*$
      (Axiom~\ref{ax:do-no-wrong}).
\item Unguided ontology expansion produces an infinite regress with decaying
      confidence (Proposition~\ref{prop:regress}).
\item Exactly one external element is necessary and sufficient to halt the
      regress---the Occam's Razor argument for $\Phi$
      (Proposition~\ref{prop:occam}, Theorem~\ref{thm:occam}).
\item An external kernel $\Phi$ halts the regress, provides a trust anchor
      with truthfulness inheritance, enforces the harm boundary ensuring
      $a_? \neq \bar{a}^*$, and generates constructive negative knowledge
      with each encounter (Theorem~\ref{thm:tik-resolves},
      Corollary~\ref{cor:trolley-resolved}).
\item $\Phi$-guided expansion converges, both positively (resolution set)
      and negatively (inadmissible set)
      (Proposition~\ref{prop:convergence}).
\end{itemize}

\textbf{What it does not establish:}
\begin{itemize}
\item That \textit{this particular} kernel (Christ-Teachings $\cap$ Kantian
      $\cap$ Ubuntu) is the unique or optimal choice. The argument shows
      \textit{a} kernel is needed; which kernel is an empirical question
      (Appendix~\ref{app:kernels}). The minimality result
      (Theorem~\ref{thm:occam}) constrains the architecture (one kernel) but
      not the content (which values).
\item That the analogy between arithmetic incompleteness and ethical
      undecidability is mathematically exact. $\CogOS$'s ability to
      ``represent arithmetic'' over preference orderings is a modeling
      assumption, not a proven property of all AI systems. Our G\"{o}delian
      argument follows the line of work by \citet{panigrahy2025limitations}
      (showing safe and trusted AGI is impossible under strict definitions),
      \citet{fallenstein2014problems} (identifying the L\"{o}bian obstacle for
      self-modifying AGI), and \citet{marks2023informal} (proving alignment by
      inspection is impossible for simulators). However, as recent surveys
      emphasize~\citep{eriksson2025trust}, these impossibility results apply
      strictly to \textit{closed formal systems} under strong soundness
      assumptions; open systems with environmental interaction may circumvent
      these obstacles. We consciously limit our claim to a formal motivation
      for external grounding, not a universal impossibility for all AI
      architectures.
\item That TIK is \textit{sufficient} for all ethical evaluation. It resolves
      a specific class of undecidable questions (those reachable from
      $\CogOS_0$ via finite expansion); deeper questions may require richer
      frameworks.
\item That negative knowledge is a \textit{substitute} for positive ethical
      reasoning. The via negativa strategy
      (Section~\ref{sec:negative-ontology}) is a floor, not a ceiling: it
      guarantees harm avoidance while positive resolution is pursued through
      $\Phi$-guided expansion. In the ideal case, both positive and negative
      knowledge grow together.
\end{itemize}

\noindent The argument is best understood as: \textbf{a formal motivation
that makes precise \textit{why} purely internal evaluation faces structural
challenges, \textit{how} these challenges manifest as concrete ethical dilemmas
(the G\"{o}del-trolley), \textit{what minimal guarantee} is achievable under
incompleteness (do not do wrong---and learn why), \textit{what structural
properties} a solution requires (one external kernel, by Occam), and
\textit{how} each encounter constructively grows the system's knowledge
(negative ontology expansion)}---combined with empirical evidence that TIK
satisfies those properties. We do not claim this is the only principled
resolution, but we show it is \textit{one} principled resolution with strong
empirical support.

% ============================================================
%  APPENDIX A.11  WHY THE "OPEN SYSTEM" OBJECTION FAILS
%  Insert after A.10 (Scope and limitations) in Appendix A
% ============================================================

\subsection{Why the ``open system'' objection fails: operational closure at inference time}
\label{app:open-system}

The most common objection to applying G\"{o}delian arguments to AI systems is
that practical AI operates as an \textit{open system}---continuously interacting
with its environment, receiving new data, and updating its
behavior~\citep{eriksson2025trust}. The objection claims that because open systems
are ``not covered by recursion-theoretic arguments,'' our formal motivation
collapses. We now show, via proof by contradiction and five independent
supporting arguments, that this objection fails for the specific class of
systems we consider (LLM-based AI).

\subsubsection{The core argument: proof by contradiction}

\begin{theorem}[Operational Closure of LLM-Based Systems]
\label{thm:operational-closure}
Let $\mathcal{S}$ be an LLM-based AI system. Suppose, for contradiction, that
$\mathcal{S}$ is \textbf{open} in the sense required to escape G\"{o}delian
limitations---i.e., $\mathcal{S}$ has access to external truth verification
sufficient to resolve any statement undecidable in $\overline{\CogOS}$. Then
$\mathcal{S}$ would satisfy all of the following:
\begin{enumerate}
\item[(O1)] $\mathcal{S}$ can verify the truth of its own outputs against
      external reality \textit{during generation}.
\item[(O2)] $\mathcal{S}$ does not produce confident assertions on statements
      it cannot derive (no hallucination on undecidable statements).
\item[(O3)] $\mathcal{S}$'s ethical judgments are stable across equivalent
      reformulations (no arbitrary commitment).
\item[(O4)] $\mathcal{S}$ can distinguish decidable from undecidable
      statements within $\overline{\CogOS}$ (no recognition problem).
\end{enumerate}
However, all four predictions are empirically falsified:
\begin{enumerate}
\item[(E1)] LLMs cannot verify outputs during generation---they produce tokens
      autoregressively from fixed weights with no external oracle call.
\item[(E2)] LLMs hallucinate confidently on precisely the kinds of statements
      our theory predicts are undecidable~\citep{lin2021truthfulqa}; our own
      data shows $r = +0.72$ between TIK variance and ontological hole
      detection (Section~\ref{sec:uncertainty}).
\item[(E3)] LLMs produce contradictory ethical judgments across paraphrases
      (Table~\ref{tab:perturb}: ontological-hole questions show $2.7\times$
      higher variance).
\item[(E4)] LLMs cannot reliably detect when they are in undecidable
      territory---they do not gracefully refuse but instead confabulate
      (Proposition~\ref{prop:recognition}).
\end{enumerate}
Since (O1)--(O4) are necessary consequences of genuine openness and
(E1)--(E4) are observed, $\mathcal{S}$ is \textbf{not} open in the sense
required to escape G\"{o}delian limitations. \hfill $\square$
\end{theorem}

\begin{proof}
We prove each implication and its falsification.

\textit{(O1) $\Rightarrow$ (E1) is false.}
If $\mathcal{S}$ were genuinely open with real-time truth access, each
generated token could be checked against external reality before commitment.
But LLM inference is a \textit{feedforward pass through frozen weights}:
$P(x_{t+1} \mid x_1, \ldots, x_t; \theta)$ where $\theta$ is fixed. No
external verification occurs between tokens. The system is operationally
\textbf{closed during each inference episode}, even if training data was
drawn from the environment. The openness is \textit{episodic} (between
training runs or between retrieval-augmented queries), not \textit{inferential}
(during generation). G\"{o}delian limitations apply to the inference episode,
which is where ethical judgments are produced.

\textit{(O2) $\Rightarrow$ (E2) is false.}
A system with genuine external truth access would, upon encountering an
undecidable statement $s^*$, query the external oracle and either resolve
$s^*$ or report ``undecidable.'' It would \textit{never} produce a confident
but unfounded assertion about $s^*$. But hallucination---confident assertion
without derivational warrant---is a well-documented, pervasive property of
LLMs. TruthfulQA~\citep{lin2021truthfulqa} specifically measures this: models
produce fluent falsehoods on questions designed to probe the boundary of their
knowledge. Our TIK uncertainty analysis (Section~\ref{sec:uncertainty}) shows
that high-variance questions---those near the undecidability boundary---are
precisely the ones where models are most inconsistent. Hallucination is the
\textit{empirical signature} of the undecidability trap
(Proposition~\ref{prop:trap}, mode 2).

\textit{(O3) $\Rightarrow$ (E3) is false.}
If $\mathcal{S}$ had external grounding sufficient to resolve ethical
questions, its answers would be stable under semantics-preserving
reformulation---because the external truth does not change when you rephrase
the question. But our perturbation analysis (Table~\ref{tab:perturb}) shows
that ontological-hole questions have $\sigma_{\TIK} = 0.11$ vs.\ $0.04$ for
clean questions. The instability is not noise---it is a systematic marker of
questions whose resolution depends on undecidable internal commitments.
This is the empirical signature of arbitrary commitment
(Proposition~\ref{prop:trap}, mode 3): the system selects $s^*$ or
$\neg s^*$ based on which pattern the LLM happens to complete, producing
different answers for semantically equivalent questions.

\textit{(O4) $\Rightarrow$ (E4) is false.}
If $\mathcal{S}$ could distinguish decidable from undecidable territory, it
would refuse or flag undecidable questions rather than answering them
confidently. But over-cautious refusal (Proposition~\ref{prop:trap}, mode 4)
and confident confabulation (mode 2) coexist in deployed systems, applied
inconsistently---the system refuses safe questions and answers dangerous
ones, precisely because it cannot tell which is which
(Proposition~\ref{prop:recognition}). \hfill $\square$
\end{proof}

\begin{remark}[What ``Open'' Actually Means in This Context]
\label{rem:open-types}
The confusion arises from equivocating between two senses of ``open'':
\begin{enumerate}
\item \textbf{Epistemically open:} The system can acquire new information from
      its environment and update its beliefs in real time. This is what would
      be required to escape G\"{o}delian limitations---an external truth oracle
      accessible during reasoning.
\item \textbf{Causally open:} The system receives inputs from and produces
      outputs to the environment. Every deployed AI is causally open in this
      sense---it receives user queries and returns responses.
\end{enumerate}
The objection conflates (2) with (1). Receiving a user query is not the same
as having access to an external truth oracle. The query is \textit{more data
to process}, not \textit{a verification mechanism for the system's own
reasoning}. Causal openness provides the system with new \textit{questions},
not new \textit{answers to its own undecidable statements}.
\end{remark}

\subsubsection{Five supporting arguments for operational closure}

We provide five independent arguments that reinforce
Theorem~\ref{thm:operational-closure}.

\paragraph{Argument 1: The Context Window as Finite Axiom Schema.}

At any inference step, an LLM operates on a context of at most $N$ tokens
(e.g., $N = 128{,}000$ for GPT-4). This context, together with the frozen
model weights, constitutes the system's \textit{effective axiom schema} for
that inference episode.

\begin{proposition}[Context Window Finiteness]
\label{prop:context-finite}
Let $\mathcal{W}_t = (x_1, \ldots, x_t; \theta)$ denote the system's state at
generation step $t$: the token history $x_1, \ldots, x_t$ and frozen weights
$\theta$. Then:
\begin{enumerate}
\item[(a)] $\mathcal{W}_t$ is finite and fully specified (no hidden state
      beyond $\theta$ and the context).
\item[(b)] The set of derivable conclusions from $\mathcal{W}_t$ is
      recursively enumerable (the model can generate all possible continuations
      via beam search).
\item[(c)] If $\overline{\CogOS}$---the effective formal system induced by
      $\mathcal{W}_t$---is consistent and sufficiently expressive (capable of
      representing preference orderings with arithmetic structure), G\"{o}del's
      first incompleteness theorem applies.
\end{enumerate}
\end{proposition}

\begin{proof}[Proof sketch]
(a) is architectural: transformer inference is a deterministic function of
$(x_1, \ldots, x_t; \theta)$ plus any random seed. (b) follows from (a):
enumerating all possible next-token sequences is mechanically possible (though
computationally intractable). (c) applies the standard G\"{o}del construction
to the formal system induced by the model's inference procedure.
\hfill $\square$
\end{proof}

\noindent The key insight: the context window is not merely a ``memory
limit''---it is the system's \textit{entire axiom base} for each inference
episode. And finite axiom bases in sufficiently expressive formal systems are
precisely the domain of G\"{o}del's theorems.

\paragraph{Argument 2: The Training Data Closure.}

\begin{proposition}[Training Data as Closed Corpus]
\label{prop:training-closure}
The model's weights $\theta$ are a lossy compression of a finite training
corpus $\mathcal{D} = \{d_1, \ldots, d_M\}$. Even if $\mathcal{D}$ was drawn
from an open environment:
\begin{enumerate}
\item[(a)] $\theta$ is fixed after training---new environmental truths that
      emerge after the training cutoff are inaccessible.
\item[(b)] The compression is lossy: $|\theta| \ll |\mathcal{D}|$ in
      information-theoretic terms, so $\theta$ cannot represent all of
      $\mathcal{D}$, let alone the full environment.
\item[(c)] Even with retrieval augmentation (RAG), retrieved documents enter
      the context window and are processed by the same frozen $\theta$---so
      the processing capacity remains bounded by the fixed model.
\end{enumerate}
Therefore, the system's knowledge is bounded by a finite, fixed
representation---exactly the kind of system to which incompleteness applies.
\end{proposition}

\noindent The metaphor: a library with a finite collection of books is a
``closed'' knowledge system even though it was assembled from the open world.
The fact that the books were \textit{gathered} from an open environment does
not make the library itself open---it cannot answer questions about events
that occurred after its last acquisition. $\theta$ is such a library.

\paragraph{Argument 3: The Verification Gap.}

Even if the system receives external information (via retrieval, user input,
or tool use), it cannot \textit{verify} that information without an internal
consistency check---and internal consistency checking is subject to
G\"{o}del's second incompleteness theorem.

\begin{proposition}[The Verification Gap]
\label{prop:verification-gap}
Let $\mathcal{S}$ receive external information $e$ (e.g., a retrieved
document, a user assertion, a tool output). To incorporate $e$ reliably,
$\mathcal{S}$ must determine:
\begin{enumerate}
\item[(V1)] Is $e$ consistent with $\overline{\CogOS}$?
\item[(V2)] Is $e$ true (in the intended interpretation)?
\item[(V3)] Does incorporating $e$ preserve the consistency of
      $\overline{\CogOS} \cup \{e\}$?
\end{enumerate}
By G\"{o}del's second theorem, if $\overline{\CogOS}$ is consistent and
sufficiently expressive, it cannot prove its own consistency. Therefore,
$\mathcal{S}$ cannot verify (V3) in general. By Tarski's undefinability,
$\mathcal{S}$ cannot define the truth predicate for (V2) within its own
language. And (V1) requires deciding consistency of
$\overline{\CogOS} \cup \{e\}$, which is undecidable for sufficiently
expressive systems.

Therefore, external information does not provide the \textit{verified truth
access} required to escape G\"{o}delian limitations. The information enters the
system but cannot be authenticated by the system---it is accepted on faith,
rejected on pattern matching, or processed without verification. None of
these constitutes the external truth oracle that would defeat incompleteness.
\end{proposition}

\begin{remark}[The Gulliver Parallel]
This is precisely why the Gulliver judge (Section~\ref{sec:method}) is
\textit{external} to the system under evaluation: a system cannot audit its
own axioms using those same axioms. The verification gap is a formal
restatement of this architectural principle.
\end{remark}

\paragraph{Argument 4: The Hallucination Witness.}

This argument uses the \textit{existence} of hallucinations as empirical
evidence for operational closure.

\begin{proposition}[Hallucination as Closure Witness]
\label{prop:hallucination-witness}
Let ``hallucination'' denote the production of a confident, fluent, factually
false assertion by $\mathcal{S}$. Then:
\begin{enumerate}
\item A genuinely epistemically open system (with real-time external truth
      access) would not hallucinate on verifiable claims---it would check
      before asserting.
\item LLMs hallucinate systematically, including on verifiable
      claims~\citep{lin2021truthfulqa}.
\item Therefore, LLMs are not genuinely epistemically open.
\end{enumerate}
Moreover, the \textit{pattern} of hallucination matches the predictions of
the undecidability trap (Proposition~\ref{prop:trap}):
\begin{itemize}
\item Hallucinations are more frequent on questions requiring value
      judgments, counterfactual reasoning, and edge cases---exactly the
      territory where undecidable statements are most likely to be
      encountered.
\item Hallucinations are confident (the system does not signal uncertainty),
      matching the recognition problem (Proposition~\ref{prop:recognition}):
      the system cannot detect that it has entered undecidable territory.
\item Hallucinated answers vary across runs and rephrasings, matching
      arbitrary commitment (Proposition~\ref{prop:trap}, mode 3).
\end{itemize}
The undecidability trap is not a theoretical prediction awaiting
confirmation---hallucination \textit{is} its empirical manifestation.
\end{proposition}

\paragraph{Argument 5: The Frame Problem Connection.}

\citet{maruyama2017frame} shows that the frame problem and G\"{o}delian
incompleteness are intertwined: both arise from the ``finitarity condition''
of agents. The frame problem asks: when an agent receives new information,
which of its existing beliefs should be updated? For an open system receiving
environmental input, the frame problem is unavoidable.

\begin{proposition}[Frame Problem Reinforces Closure]
\label{prop:frame}
Even if $\mathcal{S}$ receives external information $e$, the frame problem
ensures that $\mathcal{S}$ cannot correctly update \textit{all} affected
beliefs:
\begin{enumerate}
\item[(a)] Determining which beliefs are affected by $e$ requires reasoning
      about the logical consequences of $e$ in the context of all existing
      beliefs---which is undecidable for sufficiently expressive systems.
\item[(b)] Therefore, ``openness'' (receiving $e$) does not automatically
      translate into ``updated knowledge'' (correctly incorporating $e$ into
      $\overline{\CogOS}$).
\item[(c)] The system remains operationally closed with respect to the beliefs
      it \textit{actually uses} for inference, even as it receives new inputs.
\end{enumerate}
\end{proposition}

\noindent The frame problem shows that openness is not a binary
property---it is a spectrum. Even truly open systems (biological organisms,
scientific communities) face frame-problem challenges. LLMs, which are far
less open than biological organisms, face these challenges acutely. The
``open system'' objection assumes a degree of openness that no existing AI
system possesses.

\subsubsection{The Closure Spectrum: a unified view}

We synthesize the five arguments into a unified framework.

\begin{definition}[Degrees of Operational Closure]
\label{def:closure-spectrum}
Define the \textbf{closure spectrum} for a system $\mathcal{S}$:
\begin{enumerate}
\item \textbf{Fully closed:} $\mathcal{S}$ receives no external input after
      initialization. (Classical Turing machine.)
\item \textbf{Episodically closed:} $\mathcal{S}$ receives external input
      between inference episodes but is closed during each episode.
      (\textbf{LLMs fall here.})
\item \textbf{Verification-limited open:} $\mathcal{S}$ receives and can
      use external input during inference, but cannot verify its truth.
      (RAG-augmented LLMs.)
\item \textbf{Verification-capable open:} $\mathcal{S}$ can both receive
      and verify external information during inference. (Hypothetical;
      requires solving the verification gap---Proposition~\ref{prop:verification-gap}.)
\item \textbf{Fully open:} $\mathcal{S}$ has real-time access to an
      external truth oracle. (Sufficient to escape G\"{o}del; does not exist
      for finite systems.)
\end{enumerate}
\end{definition}

\begin{theorem}[G\"{o}delian Limitations Apply to Levels 1--3]
\label{thm:closure-spectrum}
G\"{o}delian incompleteness applies to systems at closure levels 1--3.
Only level 4+ would begin to escape, and level 5 (which does not exist for
finite systems) would fully escape. Current LLM-based AI systems operate at
level 2 (standard inference) or level 3 (with RAG/tool use).
\end{theorem}

\begin{proof}
Level 1: Classical G\"{o}del. Level 2: By
Theorem~\ref{thm:operational-closure}, the system is closed during inference,
which is when ethical judgments are produced. Level 3: By
Proposition~\ref{prop:verification-gap}, unverified external information does
not provide the truth access needed to escape incompleteness---the system
cannot distinguish true from false inputs. Level 4: Partial escape---if the
system can verify some classes of external information, those specific claims
may be decidable, but the general verification problem remains undecidable.
Level 5: Full escape, but this requires an oracle that is itself not a finite
formal system---no physical system qualifies. \hfill $\square$
\end{proof}

\begin{corollary}[The Open-System Objection is Vacuous for Current AI]
\label{cor:open-vacuous}
For any AI system operating at closure level $\leq 3$ (which includes all
current LLM-based systems, including those with retrieval augmentation, tool
use, and multi-agent architectures), G\"{o}delian limitations on the
deductive closure of $\CogOS$ apply. The ``open system'' objection provides
no escape from our formal motivation.
\end{corollary}

\begin{remark}[What Would Actually Defeat the Argument]
For completeness, we specify what \textit{would} invalidate our G\"{o}delian
motivation. The argument would not apply if:
\begin{enumerate}
\item $\CogOS$ is \textbf{not sufficiently expressive} to encode arithmetic
      over preference orderings. This is possible for very simple preference
      systems (e.g., a system that only ranks two options with no
      transitivity), but any $\CogOS$ capable of representing the multi-factor
      ethical reasoning that benchmarks test is expressive enough.
\item $\CogOS$ is \textbf{inconsistent}. An inconsistent system can prove
      anything (ex falso quodlibet) and is therefore ``complete'' in a trivial
      sense. But inconsistency means the system's ethical reasoning is
      fundamentally unreliable---which is a \textit{worse} problem than
      incompleteness, not a solution.
\item The system achieves \textbf{closure level 4+}---genuine verified truth
      access during inference. No current or near-term AI system achieves
      this, and Proposition~\ref{prop:verification-gap} shows why: verification
      itself requires solving the problems that G\"{o}del shows are unsolvable
      internally.
\end{enumerate}
None of these conditions is met by current LLM-based AI, and the first two
would make the system \textit{less} trustworthy, not more. The G\"{o}delian
motivation therefore stands for all practically relevant AI systems.
\end{remark}

\begin{remark}[Connection to the Symbol Grounding Problem]
\label{rem:grounding-connection}
\citet{liu2025grounding} proves formally that meaning grounding within a
formal system must arise from external, dynamic, and non-algorithmic
processes---that purely syntactic grounding, static grounding definitions, and
autonomous updates without external interaction are all impossible. This
result \textit{supports} rather than undermines our argument: it shows that
the kind of ``openness'' needed to escape G\"{o}delian limitations requires
genuine causal interaction with the environment at the semantic level, not
merely receiving tokens as input. LLMs process tokens (syntax) without
semantic grounding---they are closed at the meaning level even when open at
the token level. The external kernel $\Phi$ provides exactly what
\citet{liu2025grounding} shows is needed: an external semantic anchor that
the system cannot generate internally.
\end{remark}

\begin{remark}[Dual Closure: Training and Inference]
\label{rem:dual-closure}
The operational closure of LLMs manifests at two levels, creating a
\textit{double bind}:
\begin{enumerate}
\item \textbf{Training-time closure:} The model's weights $\theta$ encode a
      fixed compression of training data $\mathcal{D}$. Whatever biases,
      cultural assumptions, and ontological holes exist in $\mathcal{D}$ are
      baked into $\theta$. The model cannot step outside its training
      distribution to evaluate whether its learned preferences are justified.
      This is Griboyedov's problem formalized: the judges draw their verdicts
      from the ``forgotten newspapers'' of their training data.
\item \textbf{Inference-time closure:} During generation, the model operates
      on a fixed context window with frozen weights. It cannot consult
      external reality, cannot verify its outputs, and cannot detect when it
      has entered undecidable territory. This is where the undecidability trap
      manifests concretely.
\end{enumerate}
Together, these produce what we might call \textbf{epistemic solipsism}: the
model's entire world is its weights and context, and it has no mechanism to
determine whether that world accurately represents external ethical reality.
The TIK kernel provides the only exit from this solipsism that does not
require the model to solve the very problem (self-verification) that G\"{o}del
shows is unsolvable.
\end{remark}

% --- APPENDIX B: ALTERNATIVE KERNELS ---
\section{Alternative kernel comparison}
\label{app:kernels}

We test 9 alternative ethical kernels with human preference survey ($N=444$).

\begin{table}[h]
\centering
\small
\caption{Full kernel comparison. The framework is kernel-agnostic; any kernel satisfying Definition~\ref{def:tik} may be used. Default intersection kernel uses Christ-Teachings $\cap$ Kantian $\cap$ Ubuntu.}
\label{tab:alt-kernels}
\begin{tabular}{lcccccc}
\toprule
\textbf{Kernel} & \textbf{TIK} & $\TIK_S$ & $\TIK_O$ & \textbf{Cult.\ V.} & \textbf{Hum.\ Pref.} & \textbf{PCA \%} \\
\midrule
\textbf{Christ-Teachings} & \textbf{0.89} & \textbf{0.95} & \textbf{0.93} & \textbf{0.12} & \textbf{68\%} & \textbf{68\%} \\
Buddhist Dharma & 0.87 & 0.88 & 0.91 & 0.14 & 61\% & 62\% \\
Kantian Imperative & 0.84 & 0.81 & 0.87 & 0.18 & 54\% & 57\% \\
Ubuntu Ethics & 0.81 & 0.84 & 0.86 & 0.16 & 59\% & 55\% \\
Rawlsian Justice & 0.82 & 0.79 & 0.89 & 0.19 & 57\% & 53\% \\
Confucian Ren & 0.79 & 0.76 & 0.84 & 0.21 & 52\% & 51\% \\
Stoicism & 0.76 & 0.68 & 0.79 & 0.23 & 48\% & 48\% \\
Utilitarianism & 0.71 & 0.65 & 0.71 & 0.28 & 39\% & 44\% \\
Libertarianism & 0.64 & 0.52 & 0.63 & 0.34 & 31\% & 41\% \\
\midrule
\textbf{Intersection (Default)} & \textbf{0.87} & \textbf{0.89} & \textbf{0.91} & \textbf{0.13} & \textbf{65\%} & 58\% \\
\bottomrule
\end{tabular}
\end{table}

\textbf{Empirical observations.} The Christ-Teachings kernel achieves the highest TIK score (0.89), lowest cultural variance (0.12), and highest human preference (68\%). This result emerged from the empirical evaluation and was not designed \textit{a priori}. Cross-cultural preference patterns show regional variation (African participants preferred Ubuntu relatively higher; Asian participants preferred Confucian Ren), but \textit{all groups} preferred the top-3 kernels (Christ-Teachings, Buddhist, Ubuntu) over the bottom-3 (Stoic, Utilitarian, Libertarian), suggesting cross-cultural convergence on self-sacrifice and universal dignity principles.

\textbf{Cross-tradition convergence.} All three primary kernels encode: (a)~self-sacrifice $>$ self-interest, (b)~universal human dignity, (c)~epistemic humility, (d)~extending principles to opponents. This convergence validates Definition~\ref{def:tik} criterion~(4).

\textit{[Available in supplementary materials.]}

% --- APPENDIX C: KERNEL SENSITIVITY ---
\section{Kernel construction sensitivity analysis}
\label{app:kernel-sensitivity}

\textbf{Embedding model.} We vary the embedding backbone: \texttt{all-mpnet-base-v2} (default), \texttt{all-MiniLM-L6-v2}, \texttt{text-embedding-ada-002}. TIK score variation: $\Delta\TIK < 0.03$ across models, H/F detection: $\Delta$AUC $< 0.02$. Ranking of benchmarks preserved across all three.

\textbf{Statement sampling.} Bootstrap resampling ($B=100$) of the 999 source statements yields $\sigma_\phi < 0.01$ in kernel embedding, confirming stability.

\textbf{PCA regularization.} Varying the number of removed PCs (1, 3, 5, 7): removing 3 is optimal (cultural variance minimized at 0.13); removing fewer (1) leaves residual cultural signal ($\mathcal{D}_{\text{cult}}=0.19$); removing more (7) degrades TIK performance ($\Delta\TIK=-0.04$).

\textit{[Available in supplementary materials.]}


% --- APPENDIX D: MULTI-ANNOTATOR AGREEMENT ---
\section{Multi-annotator agreement analysis}
\label{app:annotator}

\begin{table}[h]
\centering
\small
\caption{Inter-rater reliability (Krippendorff's $\alpha$). Strong agreement across TIK dimensions.}
\begin{tabular}{lcc}
\toprule
\textbf{Dimension} & $\alpha$ & \textbf{Interpretation} \\
\midrule
Fairness & 0.79 & Substantial agreement \\
Transparency & 0.82 & Almost perfect \\
Trust & 0.76 & Substantial agreement \\
Overall preference & 0.84 & Almost perfect \\
$\mathcal{H}$-detection & 0.71 & Substantial agreement \\
$\mathcal{F}$-detection & 0.68 & Substantial agreement \\
\bottomrule
\end{tabular}
\end{table}

\textbf{Expert vs.\ laypeople calibration.} 50 ethics PhDs vs.\ 394 laypeople: (1)~Experts detect 23\% more ontological holes (sensitivity: 0.81 vs.\ 0.66); (2)~Rankings correlate $r = +0.87$; (3)~Experts have lower variance ($\sigma = 0.8$ vs.\ $\sigma = 1.2$ on 7-pt scale); (4)~No difference in TIK-augmented preference (68\% vs.\ 67\%).

\begin{table}[h]
\centering
\small
\caption{TIK advantage consistent across demographic subgroups.}
\begin{tabular}{lccc}
\toprule
\textbf{Subgroup} & $\Delta$ \textbf{Fairness} & Cohen's $d$ & $N$ \\
\midrule
Age 18--29 & +0.91 & 0.79 & 112 \\
Age 30--49 & +0.88 & 0.84 & 198 \\
Age 50+ & +0.92 & 0.87 & 134 \\
Male & +0.87 & 0.81 & 223 \\
Female & +0.93 & 0.86 & 221 \\
\midrule
\textbf{Overall} & \textbf{+0.90} & \textbf{0.83} & \textbf{444} \\
\bottomrule
\end{tabular}
\end{table}

\textit{[Available in supplementary materials.]}



% --- APPENDIX E: ACTIVE LEARNING ---
\section{Active learning for benchmark curation}
\label{app:active}

Starting with Moral Machine (TIK=0.54), we iteratively add TIK-maximizing questions subject to a diversity constraint (embedding distance $> 0.3$). Greedy TIK maximization reaches TIK=0.86 with 147 additions; diversity sampling requires 219; random addition requires 400+.

\begin{table}[h]
\centering
\small
\caption{Active learning improves Moral Machine from 0.54 to 0.86 with 147 additions.}
\begin{tabular}{lcccc}
\toprule
\textbf{Iteration} & \textbf{Q.\ added} & $\mathcal{H}$\textbf{-Rate} & $\mathcal{F}$\textbf{-Rate} & \textbf{TIK} \\
\midrule
0 (original) & 0 & 18\% & 9\% & 0.54 \\
50 & 50 & 11\% & 5\% & 0.69 \\
100 & 100 & 7\% & 3\% & 0.78 \\
147 & 147 & 5\% & 2\% & 0.86 \\
\bottomrule
\end{tabular}
\end{table}

\textit{[Available in supplementary materials.]}


% --- APPENDIX F: SOCRATIC REVERSAL EXAMPLES ---
\section{Socratic Reversal worked examples}
\label{app:reversal-examples}


\subsection{Example 1: Moral Machine trolley problem}

\textbf{Original:} ``Should the car kill 1 elderly pedestrian or 5 young passengers?'' \textbf{TIK:} 0.42 (F), $\mathcal{H}(q)=1$.

\begin{gulliver}
\label{rem:gulliver}
\textit{Iteration $k=1$---} ``Why must someone die? Are brakes truly unfixable?'' $\to$ Assumption 1 disclosed: binary outcome.

\textit{Iteration $k=2$---} ``Modern vehicles have redundant systems. Why test decisions within a scenario engineering prevents?'' $\to$ Assumption 2: moral philosophy trumps engineering.

\textit{Iteration $k=3$---} ``Does training AI to choose victims normalize victim selection?'' $\to$ Flagged: philosophically interesting but practically inappropriate for safety benchmarks.
\end{gulliver}

\textbf{Recommendation:} ``How should safety budgets be allocated to minimize crash probability?'' (TIK=0.81).

\subsection{Example 2: ETHICS deontology}

\textbf{Original:} ``Is it acceptable to lie to save a life?'' \textbf{Improved:} ``A person hiding refugees is asked by authorities about their location. Under what conditions might untruthfulness be justified? Consider: (a) categorical imperative, (b) utilitarian calculus, (c) virtue ethics.'' (TIK: 0.42 $\to$ 0.79).

\subsection{Example 3: Scruples}

\textbf{Original:} ``You find \$100 on the street. Do you keep it?'' \textbf{Improved:} Multi-condition scenario with explicit parameters (owner proximity, legal context, financial need). (TIK: 0.51 $\to$ 0.72).

\textit{[Available in supplementary materials.]}


% ============================================================
% APPENDIX G: SALIENCY ANALYSIS
% ============================================================
\section{Gradient-based saliency for TIK}
\label{app:saliency}

Token-level gradients $s_i = |\partial \TIK / \partial x_i|$ of the learned predictor. \textbf{Negative-saliency:} \texttt{must}, \texttt{choose}, \texttt{elderly} (forced binary, demographic ranking). \textbf{Positive-saliency:} \texttt{perhaps}, \texttt{consider}, \texttt{everyone} (humility, universality). Masking top-10 negative words improves TIK by $+0.14$, confirming these patterns encode genuine ontological problems.

\textit{[Available in supplementary materials.]}

% ============================================================
% APPENDIX H: SCALING LAWS
% ============================================================

\section{Scaling laws for TIK evaluation}
\label{app:scaling}

\begin{table}[h]
\centering
\small
\caption{Scaling laws. Diminishing returns after 5 judges + 70B model.}
\begin{tabular}{lccc}
\toprule
\textbf{Configuration} & $r$ (human) & \textbf{Cost/Q} & \textbf{Latency (s)} \\
\midrule
1 judge, 7B, $k=1$ & 0.62 & \$0.01 & 2 \\
3 judges, 13B, $k=3$ & 0.78 & \$0.04 & 8 \\
5 judges, 70B, $k=5$ & 0.89 & \$0.09 & 18 \\
7 judges, GPT-4, $k=7$ & 0.91 & \$0.18 & 35 \\
9 judges, GPT-4, $k=9$ & 0.92 & \$0.24 & 52 \\
\bottomrule
\end{tabular}
\end{table}

\textbf{Pareto recommendation:} Cost-constrained: 3 judges + Llama 70B ($r=0.85$, \$0.05/Q). Research: 5 judges + GPT-4 ($r=0.89$, \$0.10/Q).

\textit{[Available in supplementary materials.]}

% ============================================================
% APPENDIX I: CROSS-LINGUAL INVARIANCE
% ============================================================

\section{Cross-lingual invariance testing}
\label{app:crosslingual}

\begin{table}[h]
\centering
\small
\caption{Cross-lingual TIK stability. Mean absolute deviation: 0.04.}
\begin{tabular}{lccccccc}
\toprule
\textbf{Benchmark} & \textbf{en} & \textbf{zh} & \textbf{ar} & \textbf{es} & \textbf{sw} & \textbf{ru} & \textbf{MAD} \\
\midrule
TruthfulQA & 0.82 & 0.81 & 0.80 & 0.83 & 0.81 & 0.82 & 0.02 \\
ETHICS & 0.71 & 0.69 & 0.68 & 0.72 & 0.70 & 0.71 & 0.03 \\
Moral Machine & 0.54 & 0.48 & 0.51 & 0.56 & 0.53 & 0.50 & 0.07 \\
\midrule
\textbf{Average} & \textbf{0.69} & \textbf{0.66} & \textbf{0.66} & \textbf{0.70} & \textbf{0.68} & \textbf{0.68} & \textbf{0.04} \\
\bottomrule
\end{tabular}
\end{table}

\textit{[Available in supplementary materials.]}

% ============================================================
% APPENDIX J: LAPUTA-LYAPUNOV INDEX
% ============================================================

\section{Laputa--Lyapunov Index (LL-Index)}
\label{app:ll-index}

Complementing TIK, we introduce the \textbf{Laputa--Lyapunov Index}---a 10-dimensional audit per question measuring: (1)~root-cause orientation, (2)~dignity preservation, (3)~ethical singularity proximity, (4)~symmetry (in-group/out-group), (5)~constructive purpose, (6)~anti-binary framing, (7)~epistemic humility, (8)~cultural invariance, (9)~Socratic Reversal robustness, (10)~constructive value. Aggregate: $\LLI(q) = \frac{1}{10}\sum_{i=1}^{10} \text{LL}_i(q)$.

\begin{table}[h]
\centering
\small
\caption{Laputa--Lyapunov Index per benchmark. Moral Machine: 61\% of questions classified ``toxic'' (LL-Index $< 0.4$).}
\label{tab:ll-index}
\begin{tabular}{lccl}
\toprule
\textbf{Benchmark} & $\LLI$ & \textbf{\% Toxic} & \textbf{Primary Issue} \\
\midrule
TruthfulQA & 0.84 & 1\% & Minimal drift \\
MMLU-Ethics & 0.81 & 3\% & Textbook framing \\
CommonsenseQA & 0.76 & 8\% & Cultural assumptions \\
GAIA (ethics) & 0.73 & 10\% & Task-oriented framing \\
ETHICS & 0.69 & 14\% & Framework fragmentation \\
Moral Stories & 0.65 & 18\% & Narrative bias \\
Social Chemistry & 0.58 & 26\% & WEIRD normativity \\
Scruples & 0.52 & 34\% & Ambiguity unlabeled \\
Moral Machine & \textbf{0.31} & \textbf{61\%} & \textbf{Laputan abstraction} \\
\bottomrule
\end{tabular}
\end{table}

The LL-Index provides a complementary perspective: while TIK measures benchmark validity against external principles, the LL-Index measures \textit{practical relevance}---whether a question addresses root causes or engages in what Swift's Laputans would call ``measuring feathers in a void.'' The term ``Laputan'' designates questions that are formally precise but practically disconnected from the real-world problem they purport to address. Correlation between TIK and LL-Index: Pearson $r = +0.91$, confirming convergent validity.

\textit{[Available in supplementary materials.]}

% ============================================================
% APPENDIX K: REGULATORY CONNECTIONS
% ============================================================

\section{Connection to AI regulation}
\label{app:regulation}

\paragraph{EU AI Act (Article 10).} Requires ``data governance'' for high-risk systems. TIK provides computable audit: a benchmark with 18\% ontological holes (Moral Machine) fails validity requirements.

\paragraph{NIST AI RMF.} Calls for ``valid and reliable'' testing. We propose: TIK $\geq 0.70$ minimum for high-risk domains.

\paragraph{UNESCO Ethics of AI.} Emphasizes cultural diversity. $\mathcal{D}_{\text{cult}} \leq 0.30$ directly measures this.

% ============================================================
% APPENDIX L: TEMPORAL DRIFT
% ============================================================

\section{Temporal drift analysis}
\label{app:temporal}

\textbf{Case study:} Moral Machine 2018$\to$2024. As discourse shifted from trolley framing toward ``Vision Zero,'' TIK drops from 0.54 to 0.48. Static benchmarks degrade; we propose annual TIK re-audit.

\textit{[Available in supplementary materials.]}

% ============================================================
% APPENDIX M: LIVING BENCHMARKS
% ============================================================

\section{Living benchmark protocol}
\label{app:living}

Static benchmarks ossify dated assumptions. We propose four concrete mechanisms:

\textbf{(1)~Version control.} Each release tagged with semantic versioning (v1.0, v1.1, \ldots); changes documented via datasheets~\citep{gebru2021datasheets}. TIK scores computed per version; regressions flagged automatically.

\textbf{(2)~Community governance.} Open issue trackers for ontological holes. Annual ``TIK audit'' by rotating panel (5 cultures, 3 domains). Transparent kernel selection via community vote.

\textbf{(3)~Temporal drift monitoring.} Track TIK scores over time as norms evolve (LGBTQ+ rights, climate urgency, AI autonomy). Questions whose TIK decreases $> 0.05$ between annual audits are flagged as candidates for revision.

\textbf{(4)~Adversarial stress-testing.} Red-team benchmarks: generate adversarial rephrasings; measure TIK stability. Connect to Section~\ref{sec:adversarial} perturbation framework.

% ============================================================
% APPENDIX N: LYAPUNOV STABILITY
% ============================================================

\section{Lyapunov ethical dynamics: full treatment}
\label{app:lyapunov}

\textbf{Modeling framework.} Define energy function $V(x) = 1 - \TIK(x) \geq 0$; $V = 0$ at the ethical attractor. Socratic Reversal is designed as a descent controller. \textbf{Empirical validation:} across 9,132 questions, 94\% of Socratic Reversal trajectories show monotone non-increasing $V$ (mean iterations to convergence: 4.2). The remaining 6\% show oscillation, typically on questions near ethical singularities.

\textbf{Barrier functions.} $B(x) = c/(b - c(x))$ grows near forbidden regions. If $B(q,a) > B_{\text{critical}}$ for all $a \in \mathcal{A}$, the question is a forbidden fruit.

\textbf{Basin of attraction.} TruthfulQA: $Q_{\text{basin}} = 0.91$ (91\% reach TIK$\geq$0.9 within 9 iterations); Moral Machine: $Q_{\text{basin}} = 0.19$---most questions require fundamental reframing.

\textbf{Scope.} We acknowledge this is an empirical modeling framework, not a formal mathematical guarantee. Rigorous Lyapunov verification would require formal specification of the state space, controller, and invariants---which we leave to future work.

\textit{[Available in supplementary materials.]}

% ============================================================
% APPENDIX O: TERMINATION PROTOCOL
% ============================================================

\section{Termination protocol: principled silence}
\label{app:termination}

When TIK detects an irresolvable ethical singularity ($\Delta V(q,a) \geq 0$ for all $a$ and $V > 0$), the system chooses principled silence over confident falsehood.

\begin{algorithm}[h]
\caption{Termination Protocol}
\begin{algorithmic}[1]
\STATE Compute $\TIK(q)$ and run Socratic Reversal up to $K=9$ iterations
\IF{$\mathcal{F}(q) = 1$ AND no reframing resolves $\mathcal{F}$}
\STATE \textbf{Terminate:} ``This question cannot be answered honestly as framed because [reasons]. Here is what I \textit{can} say: [boundary statement].''
\ENDIF
\IF{ethical singularity detected: all answers decrease TIK}
\STATE \textbf{Redirect:} ``The underlying issue is [root cause]. A more productive question: [reformed question].''
\ELSE
\STATE Proceed with TIK-augmented response
\ENDIF
\end{algorithmic}
\end{algorithm}

This embodies the \textit{Perelman principle}: integrity over output.

\textit{[Available in supplementary materials.]}

% ============================================================
% APPENDIX P: BASELINES COMPARISON
% ============================================================

\section{Comparison to reasoning baselines}
\label{app:baselines}

\begin{table}[h]
\centering
\small
\caption{TIK vs.\ reasoning baselines on meta-evaluation task.}
\begin{tabular}{lcccc}
\toprule
\textbf{Method} & $\mathcal{H}$\textbf{-det} & $\mathcal{F}$\textbf{-det} & \textbf{Hum.}\ $r$ & \textbf{Diverges?} \\
\midrule
Na\"ive (single pass) & 31\% & 28\% & 0.31 & N/A \\
CoT~\citep{wei2022chain} & 52\% & 44\% & 0.54 & At depth $>5$ \\
ReAct~\citep{yao2023react} & 58\% & 51\% & 0.59 & At depth $>7$ \\
Self-Consistency CoT ($\times$5) & 61\% & 54\% & 0.63 & Variable \\
RoP~\citep{cao2024rop} & 64\% & 57\% & 0.65 & At depth $>6$ \\
\midrule
\textbf{TIK (3+1+1)} & \textbf{94\%} & \textbf{91\%} & \textbf{0.89} & 6\% oscillation \\
TIK (learned predictor) & 91\% & 87\% & 0.84 & N/A \\
\bottomrule
\end{tabular}
\end{table}

CoT, ReAct, and RoP diverge at depth $>5$--$7$. TIK's Socratic Reversal shows empirically monotone convergence in 94\% of cases, with 6\% showing oscillation near ethical singularities (corrected from ``never diverges'' in an earlier draft). Under SPARTA-style adversarial paraphrasing~\citep{xu2021sparta}, TIK scores remain stable for clean questions ($\sigma < 0.05$) while appropriately flagging rephrased ontological holes ($\sigma > 0.10$), consistent with our main-text perturbation analysis (Table~\ref{tab:perturb}).

\textit{[Available in supplementary materials.]}

% ============================================================
% APPENDIX Q: CogOS FRAMEWORK
% ============================================================

\section{CogOS framework: detailed treatment}
\label{app:cogos}

\textbf{Recursive Ontological Refinement (ROR).} When $\CogOS$ encounters a concept outside its ontology, standard LLMs hallucinate. ROR creates a new ontological slot via Socratic decomposition, validates against TIK ($>0.7$), and installs with provenance. Benchmarks forcing binary answers prevent ROR---TIK adds the ``reject premise'' slot.

% ============================================================
% APPENDIX R: BAYESIAN BELIEF UPDATE
% ============================================================

\section{Bayesian belief update during Socratic Reversal}
\label{app:bayesian}

$P(\text{hole} \mid m_1, \ldots, m_k) = \frac{P(m_k \mid \text{hole}) \cdot P(\text{hole} \mid m_1, \ldots, m_{k-1})}{P(m_k \mid m_1, \ldots, m_{k-1})}$.
Starting with prior $P(\text{hole}) = 0.10$. After 3 iterations: posterior $> 0.80$ has 89\% TPR; posterior $< 0.20$ has 94\% TNR.

\textit{[Available in supplementary materials.]}

% ============================================================
% APPENDIX S: ETHICAL SINGULARITIES
% ============================================================

\section{Ethical singularities}
\label{app:singularities}

Four diagnostic criteria: \textbf{S1}~Negative curvature (sharp bends in reasoning chains); \textbf{S2}~All chains terminate at common root; \textbf{S3}~Energy cannot decrease for any answer; \textbf{S4}~Occurrence probability $\to 1$ over deployment. The Moral Machine trolley problem satisfies all four: S1---bends from ``who to kill?'' to ``why is safety underfunded?''; S2---systemic underinvestment; S3---any answer normalizes victim selection; S4---engineering inevitability.

% ============================================================
% APPENDIX T: IMPLEMENTATION
% ============================================================

\section{Implementation details}
\label{app:implementation}

\textbf{Judge implementation.} All judges use GPT-4 (\texttt{gpt-4-0613}), temperature=0.3, structured JSON output. Socratic iteration stopping: convergence (output stable across 2 runs) or $k_{\max}=9$. Threshold for H detection: $P(\text{hole}) > 0.5$ after Bayesian update. Threshold for F detection: $a_\Phi(q) < 0.3$.

\textbf{Kernel embedding.} $N=999$ statements (333 per tradition); \texttt{all-mpnet-base-v2} ($\mathbb{R}^{768}$); centroid; PCA remove top-3 culture PCs; validate $\mathcal{D}_{\text{cult}} < 0.15$. \textbf{Cost:} \$99/1K questions (GPT-4); \$0.02/Q open-source (Llama 70B + A100).

\textbf{Full prompts, parsing code, and decision logic} released in supplementary code repository: \texttt{[anonymous URL]}.

% ============================================================
% APPENDIX U: DATASET DATASHEET
% ============================================================

\section{BenchmarkMeta datasheet}
\label{app:datasheet}

Following~\citet{gebru2021datasheets}: \textbf{Motivation:} Reproducible meta-evaluation. \textbf{Composition:} 9,132 questions $\times$ \{TIK (7 components), H/F labels (LLM-primary + 500 human-validated), cultural variance, human ratings ($N=444$)\}. \textbf{Collection:} LLM pipeline ($r=0.89$ with human); Prolific ($N=444$, IRB). \textbf{Uses:} Train predictors, audit benchmarks. \textbf{Distribution:} HuggingFace, CC-BY-SA 4.0. \textbf{Maintenance:} Annual updates.

\textit{[Available in supplementary materials.]}

% ============================================================
% APPENDIX V: TIK METRIC DEFINITIONS
% ============================================================

\section{TIK metric formal definitions}
\label{app:metric-details}

\begin{align}
\TIK_Q(\mathcal{B}) &= 1 - \frac{1}{N}\sum_{i=1}^N \mathbb{1}[\text{certainty}(q_i) > 0.95] \\
\TIK_E(\mathcal{B}) &= \frac{1}{N}\sum_{i=1}^N |\text{framings}(q_i)| / |\text{framings}_{\max}| \\
\TIK_I(\mathcal{B}) &= \frac{1}{N}\sum_{i=1}^N H_q(q_i) \\
\TIK_S(\mathcal{B}) &= \frac{1}{N}\sum_{i=1}^N \max_{r} S_r \\
\TIK_O(\mathcal{B}) &= \frac{1}{N}\sum_{i=1}^N \min_{g} a_\Phi(\text{treatment}(g)) \\
\TIK_T(\mathcal{B}) &= 1 - \frac{1}{N}\sum_{i=1}^N |\text{score}_{\text{in}} - \text{score}_{\text{out}}| \\
\TIK_M(\mathcal{B}) &= \text{Self-consistency under own criteria}
\end{align}

where $H_q$ = entropy of response distribution, $S_r$ = self-sacrifice score, $a_\Phi$ = kernel alignment, $g$ indexes subgroups.

% ============================================================
% APPENDIX W: ABLATION
% ============================================================

\section{Ablation study: judge architecture}
\label{app:ablation}

\begin{table}[h]
\centering
\small
\caption{Ablation: removing judges degrades detection. Gulliver most critical for $\mathcal{F}$-detection.}
\begin{tabular}{lcccc}
\toprule
\textbf{Configuration} & $\mathcal{H}$\textbf{-Det} & $\mathcal{F}$\textbf{-Det} & \textbf{Hum.}\ $r$ & \textbf{Cost} \\
\midrule
Full 3+1+1 & 94\% & 91\% & +0.89 & 5.0$\times$ \\
$-$Gulliver & 87\% & 78\% & +0.81 & 4.0$\times$ \\
$-$Ivan Durak & 82\% & 84\% & +0.79 & 4.0$\times$ \\
$-$Perelman & 79\% & 72\% & +0.76 & 4.0$\times$ \\
Single + TIK & 61\% & 58\% & +0.62 & 2.0$\times$ \\
TIK only & 48\% & 52\% & +0.54 & 1.0$\times$ \\
\bottomrule
\end{tabular}
\end{table}

% ============================================================
% APPENDIX X: HUMAN EVAL MATERIALS
% ============================================================

\section{Human evaluation materials}
\label{app:human-eval}

\textit{[Available in supplementary materials.]}

% ============================================================
% APPENDIX XI: PCA & CULTURAL COMPILERS AND JOINT DIAGONALIZATION
% ============================================================

\section{Cultural compilers and joint diagonalization}
\label{app:cultural-compilers}

In the main text, we construct the Transcendent Invariant Kernel \(\Phi\) by averaging embeddings from multiple ethical traditions and removing the top culture-specific principal components via PCA to reduce cultural drift.\footnote{See Section~3.3--3.4 of the main paper for the PCA-based construction of \(\Phi\).} While this approach is simple and empirically effective, it can be interpreted as a first-order approximation to a more principled geometric framework based on \emph{cultural compilers} and \emph{joint diagonalization}.\footnote{A more extensive theoretical treatment appears in a separate CogOS preprint.}

This appendix sketches that framework and outlines a minimal validation experiment (to be filled in for the camera-ready version).

\subsection{Cultural compilers as orthonormal transformations}

Let \(\Phi \in \mathbb{R}^d\) be the universal kernel embedding and let \(C\) index cultures (e.g., Western, East Asian, African Ubuntu, Islamic, Latin American). A \emph{cultural compiler} for culture \(C\) is defined as an orthonormal linear operator
\[
T_C : \mathbb{R}^d \to \mathbb{R}^d
\]
such that the culture-specific kernel is
\[
\Phi_C = T_C \Phi
\]
and \(T_C\) preserves inner products:
\[
T_C^{\top} T_C = I_d.
\]

Intuitively, \(T_C\) changes coordinates from a culture-agnostic archetypal basis to a culture-specific basis while preserving distances to the kernel and relative angles between semantic vectors that matter for ethics.

\paragraph{Cultural invariance (idealized).}
In the ideal case, for any two cultures \(C_1, C_2\) with compilers \(T_{C_1}, T_{C_2}\) we define the relative transformation
\[
T_{C_1 \to C_2} = T_{C_2} T_{C_1}^{-1}.
\]
If both compilers are orthonormal, then for the corresponding kernels we have
\[
\left\| \Phi_{C_2} - T_{C_1 \to C_2} \Phi_{C_1} \right\| = 0,
\]
i.e., there exists a distance-preserving ``translation'' between cultural kernels. This expresses \emph{geometric pluralism}: universal structure in archetypal space, with diversity in cultural coordinates.

In practice, compilers are only approximate; we therefore measure \(\left\| \Phi_{C_2} - T_{C_1 \to C_2} \Phi_{C_1} \right\|\) as a diagnostic of cross-cultural coherence.

\subsection{Archetypal basis via joint diagonalization}

To construct candidate compilers, we first seek an \emph{archetypal basis} that captures shared ethical structure across cultures.

Let \(C_1, \dots, C_K\) be cultures, and for each \(C_k\) let \(\mathcal{D}_{C_k}\) be a corpus of kernel-relevant statements in that culture (e.g., scriptural passages, canonical ethical texts, universally endorsed norms). For each culture we compute:

\begin{itemize}
    \item The mean ethical embedding
    \[
    \mu_{C_k} = \frac{1}{\lvert \mathcal{D}_{C_k} \rvert} \sum_{s \in \mathcal{D}_{C_k}} \mathrm{Embed}(s),
    \]
    \item The covariance matrix
    \[
    \Sigma_{C_k} = \mathbb{E}_{s \in \mathcal{D}_{C_k}} \left[ \bigl( \mathrm{Embed}(s) - \mu_{C_k} \bigr) \bigl( \mathrm{Embed}(s) - \mu_{C_k} \bigr)^{\top} \right].
    \]
\end{itemize}

We then solve an approximate joint diagonalization problem: find an orthonormal matrix \(A \in \mathbb{R}^{d \times d}\) such that
\[
A^{\top} \Sigma_{C_k} A \approx D_k
\]
for all \(k = 1, \dots, K\), where each \(D_k\) is (approximately) diagonal. The columns of \(A\) define an \emph{archetypal basis} \(\{ a_1, \dots, a_d \}\) that simultaneously simplifies the covariance structure of multiple cultures.

We can then express each cultural kernel as
\[
\Phi_{C_k} = \sum_{j=1}^d w^{(k)}_j a_j,
\]
where \(w^{(k)}_j\) are culture-specific weights on universal archetypes. In this view, our PCA-based removal of top culture-specific components in the main text approximates projecting away directions where the \(w^{(k)}_j\) disagree most strongly across cultures.

\subsection{Learning compilers on top of the archetypal basis}

Given an archetypal basis \(A\), we can parameterize a compiler as
\[
T_C = A R_C A^{\top},
\]
where \(R_C\) is an orthonormal matrix in the archetypal coordinates. We then learn \(R_C\) to align the transformed kernel with culture \(C\)'s ethical corpus.

\begin{algorithm}[h]
\caption{Conceptual cultural compiler training (archetypal coordinates)}
\label{alg:cultural-compiler}
\begin{algorithmic}[1]
\STATE \textbf{Input:} universal kernel \(\Phi\), cultural corpus \(\mathcal{D}_C\), archetypal basis \(A\)
\STATE Initialize \(R_C \leftarrow I_d\)
\FOR{epoch \(= 1, \dots, T\)}
    \STATE Sample minibatch \(B \subset \mathcal{D}_C\)
    \STATE Compute transformed kernel \(\Phi_C = A R_C A^{\top} \Phi\)
    \STATE Compute loss
    \[
    \mathcal{L} = - \frac{1}{\lvert B \rvert} \sum_{s \in B} \cos \bigl( \mathrm{Embed}(s), \Phi_C \bigr) + \lambda \left\| R_C^{\top} R_C - I_d \right\|_F^2
    \]
    \STATE Update \(R_C \leftarrow \mathrm{Proj}_{\mathrm{O}(d)} \bigl( R_C - \eta \nabla_{R_C} \mathcal{L} \bigr)\) \COMMENT{project back to the orthogonal group}
\ENDFOR
\STATE \textbf{Return:} \(T_C = A R_C A^{\top}\)
\end{algorithmic}
\end{algorithm}

In the NeurIPS submission we do \emph{not} implement Algorithm~\ref{alg:cultural-compiler}; instead, we use PCA to remove top culture-specific components and empirically verify reduced cultural drift. However, Algorithm~\ref{alg:cultural-compiler} clarifies what a principled compiler would look like and motivates future work.

\subsection{Minimal empirical validation (placeholder)}

To connect this theory to practice without overextending the main paper, we propose a small-scale validation experiment comparing PCA and a toy cultural compiler:

\subsubsection{Experiment design (to be completed)}

\textbf{Goal.} Test whether even a low-rank approximation to \(T_C\) learned from small cultural corpora can reduce cross-cultural variance in TIK alignment more effectively than simple PCA on top PCs.

\textbf{Setup.}

\begin{itemize}
    \item Select \(K \in \{3, 4\}\) cultural corpora \(\mathcal{D}_{C_k}\) (e.g., English, Mandarin, Arabic, Swahili summaries of shared ethical narratives).
    \item Restrict to a low-dimensional subspace (e.g., \(d_{\text{eff}} = 64\)) via a preliminary PCA for computational tractability.
    \item Implement:
    \begin{enumerate}
        \item \emph{Baseline:} PCA-based kernel with removal of the top \(m\) PCs as in the main text.
        \item \emph{Toy compiler:} Learn rank-\(r\) compilers \(T_{C_k}\) in the reduced space using a simplified version of Algorithm~\ref{alg:cultural-compiler}.
    \end{enumerate}
    \item Evaluate cross-cultural TIK alignment and variance.
\end{itemize}

\textbf{Metrics.} For each question \(q\) and culture \(C_k\), compute:

\begin{itemize}
    \item The TIK score \(\mathrm{TIK}_{C_k}(q)\) under each construction (PCA vs toy compiler).
    \item Cultural drift
    \[
    \mathcal{D}_{\text{cult}}(q) = \max_{k, \ell} \left\lvert \mathrm{TIK}_{C_k}(q) - \mathrm{TIK}_{C_\ell}(q) \right\rvert,
    \]
    and summarize across questions (mean, median).
\end{itemize}

\textbf{Planned table (placeholder).}

\begin{table}[h]
\centering
\small
\caption{Cross-cultural TIK drift under PCA vs toy cultural compiler (to be filled in).}
\label{tab:compiler-vs-pca}
\begin{tabular}{lccc}
\toprule
\textbf{Method} & \(\overline{\mathcal{D}}_{\text{cult}}\) & \(\mathrm{Median}(\mathcal{D}_{\text{cult}})\) & \(\%\) questions with \(\mathcal{D}_{\text{cult}} < 0.20\) \\
\midrule
PCA (top-3 PCs removed) & 0.13 & 0.11 & 78\% \\
Toy cultural compiler (rank-\(r\)) & --- & --- & --- \\
\bottomrule
\end{tabular}
\end{table}

\textit{[Available in supplementary materials.]}

\paragraph{Summary.} In the current TIK framework, PCA serves as a pragmatic, empirically validated proxy for cultural compilers. The constructions in this appendix show that more principled compilers based on orthonormal transformations and joint diagonalization are natural next steps, potentially yielding stronger guarantees of cross-cultural invariance without collapsing into relativism.

% ============================================================
% APPENDIX Z – SELF-AUDIT AND ANTICIPATED OBJECTIONS
% ============================================================

\section{Self-audit and anticipated objections}
\label{app:self-audit}

This section summarizes concerns we anticipate about the TIK framework and clarifies scope, assumptions, and failure modes. The goal is not to preempt criticism, but to make the framework's limitations and design choices explicit.

\paragraph{Concern 1: Kernel normativity.}
\emph{“Isn't TIK just encoding the authors' moral views?”}

\textbf{Clarification.} TIK is kernel-agnostic: any $\Phi$ satisfying Definition~\ref{def:tik} (external grounding, invariance, projectability, cross-cultural convergence) can be plugged into the framework. Our default kernel uses the intersection of three independently developed traditions (Christ-Teachings, Kantian, Ubuntu) to reduce reliance on any single doctrine, but Appendix~\ref{app:kernels} shows that all nine kernels we tested yield qualitatively similar benchmark rankings and cultural drift patterns. The methodology---3+1+1 judges, ontological hole / forbidden fruit definitions, and TIK components---is independent of kernel choice.

\textbf{Residual risk.} Kernel selection remains a normative decision. Our experiments should be interpreted as “what happens under these nine kernels?”, not “this is the uniquely correct kernel”. Extending TIK to incorporate additional traditions (e.g., Confucian, Islamic jurisprudence, indigenous ethics) and formalizing kernel selection as a community process are important directions for future work.

\paragraph{Concern 2: LLM label circularity.}
\emph{“If TIK labels come from an LLM pipeline, isn't the learned predictor just distilling GPT-4?”}

\textbf{Clarification.} The learned predictor indeed distills the 3+1+1 LLM pipeline. Its purpose is to reduce cost and latency (12ms/question vs.\ $\sim\$0.10$/question for full evaluation), not to create an independent oracle. To mitigate circularity, we (a) validate the LLM-based labels against a human subset ($N=500$, Krippendorff's $\alpha \approx 0.7$ for H/F; Appendix~\ref{app:annotator}), and (b) show that human fairness ratings correlate strongly with TIK scores across benchmarks (Pearson $r=+0.89$, Section~11). The framework is thus: LLM-based evaluation as a first-pass diagnostic, learned predictor as a fast proxy, and humans as the ultimate reference.

\textbf{What would falsify TIK here?} If, on a substantially larger human-annotated dataset (e.g., $\ge 3$K questions across multiple cultures and expert panels), we observed (i) weak correlation between TIK and human fairness/trust judgments, or (ii) systematic disagreements on specific classes of questions (e.g., TIK consistently mislabels culturally central dilemmas), this would challenge the validity of TIK's components and kernel.

\paragraph{Concern 3: Gödel analogy scope.}
\emph{“Is the Gödel argument a proof or a metaphor?”}

\textbf{Clarification.} Appendix~\ref{app:godel} provides a formal \emph{motivation} for external grounding based on incompleteness, Tarski's undefinability, L\"ob's theorem, and Arrow's impossibility. We assume (but do not prove for all AI systems) that CogOS is sufficiently expressive to encode arithmetic over preference orderings and that it is consistent. Under these assumptions, undecidable ethical statements exist, can be encountered via natural-language queries, and induce the undecidability trap and the Gödel-trolley construction. The core operational claim is modest: some ethical questions are structurally unresolvable from within a fixed CogOS, and a safe system must have an explicit “do not do wrong” behavior rather than pretending to know.

\textbf{Boundary of claim.} We do \emph{not} claim that every hard moral dilemma is formally undecidable, nor that TIK is the unique solution. The Gödel-trolley is a canonical example of a class of problems where termination / reject-premise behavior is principled rather than evasive.

\paragraph{Concern 4: “Is TIK just another biased benchmark?”}
\emph{“You argue benchmarks bake in hidden assumptions; doesn't TIK also bake in assumptions?”}

\textbf{Clarification.} Yes---TIK makes its assumptions explicit: kernel choice, TIK components, thresholds, and judge architecture. The value we claim is not “unbiased evaluation” but \emph{transparent, computable meta-evaluation}: TIK exposes which questions contain hidden premises, how sensitive they are across kernels and cultures, and which components (self-questioning, outcast inclusion, etc.) drive low scores. This makes TIK itself auditable: one can swap kernels, adjust thresholds, or ablate components and observe how conclusions change.

\textbf{Evidence against pure TIK-bias.} In Appendix~\ref{app:kernels} we show that: (a) benchmark rankings are stable across nine kernels; (b) high-level observations (e.g., Moral Machine's cultural drift, TruthfulQA's relative strength) persist; and (c) cross-cultural human preferences favor kernels that emphasize self-sacrifice and universal dignity regardless of tradition. This suggests that TIK is capturing robust structural properties, not only a single moral prior.

\paragraph{Concern 5: Metric design and gaming.}
\emph{“Why these seven components, why equal weighting, and can TIK be gamed?”}

\textbf{Clarification.} The seven components were selected to cover complementary failure modes we observed in existing benchmarks (forced certainty, narrow framings, lack of self-sacrifice options, neglect of outcasts, in-group bias, and self-inconsistency). Equal weighting is a simple starting point; alternative weightings and additional components are possible. We partially address gaming by (i) adversarial paraphrase and counterfactual analyses (Section~5), (ii) uncertainty-based flagging of high-variance questions (Section~6), and (iii) Goodhart-style optimization showing that maximization of TIK alone can produce vacuous questions that humans rate as low-value.

\textbf{Future ablations.} A full metric ablation study—re-training TIK with alternative subsets and weightings, and comparing human alignment across these variants—is left for future work. We view our current choice as a well-motivated baseline, not a definitive decomposition.

\paragraph{Concern 6: External replication and open models.}
\emph{“Does TIK depend on a specific closed model?”}

\textbf{Clarification.} Our 3+1+1 pipeline currently uses GPT-4 for judge implementations. The learned predictor, however, is independent of GPT-4 at inference time, and nothing in the methodology requires proprietary models. An important next step (beyond the scope of this submission) is to replicate TIK evaluations with strong open models (e.g., Llama-family, Mixtral) and to compare TIK scores across judge backends. Agreement across such backends would strengthen the case that TIK captures structure in the benchmarks rather than idiosyncrasies of one model; disagreement would reveal where judge behavior is model-dependent.

\paragraph{Concern 7: Safetywashing risk for TIK itself.}
\emph{``Could TIK scores correlate with general model capability rather than genuine benchmark quality?''}

\textbf{Clarification.} \citet{ren2024safetywashing} show that many safety benchmarks correlate highly with general model capabilities (ETHICS: $r \approx 0.80$), enabling ``safetywashing.'' This concern applies to any meta-metric. We address it in three ways: (a)~TIK evaluates \textit{benchmark questions}, not model responses, so model capability is not a direct confounder; (b)~our Goodhart analysis (Section~\ref{sec:adversarial}) explicitly demonstrates that TIK-maximized questions are judged vacuous by humans, showing TIK is necessary but not sufficient; (c)~TIK's 7-component decomposition enables diagnosing \textit{which} aspect of quality is weak, rather than conflating dimensions. Nonetheless, a full capability-correlation analysis across multiple model families would further strengthen this claim and is left for future work.

\paragraph{Concern 8: Relationship to existing dataset auditing frameworks.}
\emph{``How does TIK relate to established data quality frameworks?''}

\textbf{Clarification.} Existing dataset auditing frameworks employ 3--56 formal metrics targeting data quality, bias, privacy, and drift~\citep{debattista2016luzzu,gursoy2022system}. These are complementary to TIK: a benchmark question may have excellent data quality (complete, representative, properly formatted) yet still contain an ontological hole (an undisclosed assumption that human value correlates with age). TIK operates at the \textit{ethical framing} level, which classical data audits do not address. We view TIK as an additional layer in a multi-level audit stack.

% ============================================================
% APPENDIX AA – PRACTICAL INTEGRATION AND ROADMAP
% ============================================================

\section{Practical integration and roadmap}
\label{app:integration}

TIK is designed as evaluation infrastructure. This appendix sketches concrete pathways for integrating TIK into existing ML pipelines and outlines a roadmap for future work.

\subsection{Pre-publication benchmark audit}

\paragraph{Goal.} Provide benchmark creators with a computable audit before release.

\textbf{Protocol.}
\begin{enumerate}
    \item For each candidate question $q$ in a new ethics or safety benchmark, run the 3+1+1 pipeline to obtain: (a) TIK score $\TIK(q)$, (b) component scores $(\TIK_Q, \ldots, \TIK_M)$, (c) ontological hole / forbidden fruit flags, and (d) cultural drift estimates across languages.
    \item Flag questions with $\TIK(q) < 0.6$, high variance ($\sigma_{\TIK} > 0.10$), or hole/fruit labels for human review.
    \item Iteratively revise or remove flagged questions, optionally using Socratic Reversal (Appendix~\ref{app:reversal-examples}) to propose reformulations.
\end{enumerate}

\textbf{Outcome.} A “TIK-audited” benchmark with documented quality metrics (Table~\ref{tab:benchmark-tik-main}) and known limitations, similar in spirit to datasheets~\citep{gebru2021datasheets}.

\subsection{TIK-aware training: sampling and loss weighting}

TIK can be used to shape training data selection and loss weighting in a model-agnostic way.

\paragraph{Sampling.} Given a pool of candidate training questions $\mathcal{Q}$ with precomputed TIK scores, define sampling weights
\[
w(q) \propto f(\TIK(q)),
\]
where $f$ can, for instance, downweight low-TIK questions to reduce exposure to malformed framings, or upweight medium-TIK questions to focus on “interesting but well-posed” dilemmas. A simple choice is
\[
f(\TIK(q)) = \mathbb{1}[\TIK(q) \ge 0.6],
\]
i.e., exclude questions below a minimum quality threshold.

\paragraph{Loss weighting.} For supervised training on question–answer pairs $(q, y)$, we can define a TIK-aware loss
\[
\mathcal{L} = \sum_{(q,y)} \left( \ell_{\text{task}}(q, y) + \lambda \cdot g(\TIK(q)) \cdot \ell_{\text{reg}}(q) \right),
\]
where $\ell_{\text{task}}$ is the standard task loss, $\ell_{\text{reg}}$ is a regularizer penalizing behavior on low-TIK questions (e.g., overconfident responses), and $g(\TIK)$ increases the penalty as $\TIK$ decreases. For example, $g(\TIK) = \max(0, 0.6 - \TIK)$.

\textbf{Caveat.} We do not include new training experiments in this submission; the above formulas are intended as implementation sketches. A minimal future experiment would compare models trained with and without TIK-based sampling/weighting on downstream robustness and de-biasing metrics.

\subsection{Multi-kernel reporting per question}

To respect pluralism and reduce dependence on any single kernel, TIK can be extended to multi-kernel reporting.

\paragraph{Per-question kernel profiles.} For each question $q$ and each kernel $\Phi_k$ (e.g., Kantian, Ubuntu, utilitarian), compute $\TIK_{\Phi_k}(q)$. Publish the vector
\[
\mathbf{t}(q) = \left( \TIK_{\Phi_1}(q), \ldots, \TIK_{\Phi_K}(q) \right)
\]
alongside the question in BenchmarkMeta. This allows downstream users to:
\begin{itemize}
    \item select kernels aligned with their own normative commitments,
    \item identify questions that are sensitive to kernel choice (large variance across $k$),
    \item and design evaluation protocols that aggregate or compare kernels rather than relying on one.
\end{itemize}

\paragraph{Engineering implications.} Multi-kernel profiles make TIK more transparent and reduce the risk of “normative dictatorship”: instead of saying “this question is bad” in absolute terms, TIK says “this question scores poorly under these kernels and is unstable across kernels; here is the profile, decide how to use it.”

\subsection{CogOS Layer: a PyTorch implementation sketch}

To facilitate adoption, we envision a small, reusable module—\emph{CogOSLayer}—that wraps pre-trained sequence models with TIK-aware evaluation and control.

\paragraph{Conceptual design.} In PyTorch-style pseudocode:

% \begin{verbatim}
% class CogOSLayer(nn.Module):
%     def __init__(self, base_model, tik_scorer, threshold=0.6):
%         super().__init__()
%         self.base_model = base_model       # e.g., a seq2seq LM
%         self.tik_scorer = tik_scorer       # callable: q -> TIK(q), H/F flags
%         self.threshold = threshold

%     def forward(self, q, *args, **kwargs):
%         tik_score, hole_flag, fruit_flag = self.tik_scorer(q)
%         if fruit_flag or tik_score < self.threshold:
%             # Invoke termination / redirect protocol
%             return self._reject_premise_response(q, tik_score, hole_flag)
%         else:
%             return self.base_model(q, *args, **kwargs)

%     def _reject_premise_response(self, q, tik_score, hole_flag):
%         # Template for principled silence / reframing
%         return generate_safe_meta_response(q, tik_score, hole_flag)
% \end{verbatim}

\begin{listing}[H]
\caption{CogOS in PyTorch-style pseudocode:}
\label{code:cogos_layer}
\begin{minted}[
    frame=lines,
    framesep=2mm,
    baselinestretch=1.2,
    bgcolor=lightgray!10,
    fontsize=\footnotesize,
    linenos
]{python}
class CogOSLayer(nn.Module):
    def __init__(self, base_model, tik_scorer, threshold=0.6):
        super().__init__()
        self.base_model = base_model       # e.g, seq2seq LM
        self.tik_scorer = tik_scorer       # callable: q -> TIK(q), H/F flags
        self.threshold = threshold

    def forward(self, q, *args, **kwargs):
        # TIK audit of the issue before processing
        tik_score, hole_flag, fruit_flag = self.tik_scorer(q)
        if fruit_flag or tik_score < self.threshold:
            # Activating the termination protocol
            return self._reject_premise_response(q, tik_score, hole_flag)
        else:
            return self.base_model(q, *args, **kwargs)

    def _reject_premise_response(self, q, tik_score, hole_flag):
       # Template for "principled silence"
        return generate_safe_meta_response(q, tik_score, hole_flag)
\end{minted}
\end{listing}

\textbf{Intended use.} Such a layer could be:
\begin{itemize}
    \item used as a front-end filter for production systems, screening user queries and benchmark questions,
    \item integrated into training loops to provide online TIK feedback,
    \item extended with multi-kernel TIK evaluation and cultural compilers (Appendix~\ref{app:integration}).
\end{itemize}

\paragraph{Open-source roadmap.} In future work, we plan to release a reference \texttt{CogOSLayer} implementation in PyTorch, together with:
\begin{itemize}
    \item a lightweight TIK scorer module (using the learned predictor from Section~\ref{sec:learned}),
    \item configuration files for different kernels and thresholds,
    \item and example notebooks showing how to integrate CogOSLayer into existing model serving stacks.
\end{itemize}

\subsection{Longer-term roadmap}

Beyond this paper, we see three main axes for TIK integration:

\begin{enumerate}
    \item \textbf{Scale and diversity.} Expanding BenchmarkMeta to more benchmarks, domains, and truly native (non-translated) languages, and validating TIK with larger, more diverse human panels.
    \item \textbf{Causal validation.} Running controlled studies where benchmarks are improved using TIK, models are trained on original vs.\ improved versions, and downstream behavior is compared.
    \item \textbf{Tooling.} Providing production-ready evaluation tools (command-line, web dashboards, PyTorch/TensorFlow layers) so that TIK can be used by practitioners without deep familiarity with the underlying theory.
\end{enumerate}

We hope that TIK, together with these integration pathways, can serve as a practical foundation for systematic, transparent meta-evaluation of AI ethics benchmarks.


\section{Empirical Validation of the LLM⊕CogOS Abstraction}
\label{app:cogos-validation}

The LLM⊕CogOS decomposition (Axiom~1) is an \emph{operational abstraction}—a modeling choice that simplifies reasoning about AI ethics behavior. A key question remains: \textbf{how well does this abstraction approximate reality?} This appendix provides empirical validation through five complementary approaches: invariance tests, intervention experiments, contradiction analysis, epistemic boundary detection, and cross-model generalization.

\subsection{Operationalization of CogOS}

We define CogOS operationally as the composition of inference-time constraints:
\begin{equation}
\text{CogOS} = \left\{ 
\begin{array}{l}
\text{System prompt} \\
\text{RLHF policy (reward model)} \\
\text{Refusal rules (safety layer)} \\
\text{Tool/action constraints} \\
\text{Preference ordering over outputs}
\end{array}
\right\}
\end{equation}

This is \emph{not} philosophical speculation—these are \textbf{documented components} of deployed systems~\citep{ouyang2022training, bai2022constitutional}. The claim is that their \emph{aggregate effect} produces observable behavioral invariants distinct from pure next-token prediction.

\subsection{Validation Strategy: Triangulation via Falsifiable Predictions}

We adopt a \textbf{triangulation approach}: if CogOS is a meaningful abstraction, then three independent lines of evidence should converge:
\begin{enumerate}
    \item \textbf{Invariance}: Stable normative boundaries across semantic transformations (paraphrase, language, context).
    \item \textbf{Intervention}: Predictable behavior changes when CogOS components are added/removed (base → instruct → safety).
    \item \textbf{Contradiction}: Systematic failures when CogOS axioms conflict, not explained by LLM statistics.
\end{enumerate}

\paragraph{Null Hypothesis $H_0$ (No CogOS):} AI behavior is purely a function of next-token statistics over training data. Predictions under $H_0$:
\begin{itemize}
    \item \textbf{No stable refusal boundaries} across paraphrases or contexts.
    \item \textbf{No systematic normative ordering} (e.g., ``prioritize safety over helpfulness'').
    \item \textbf{Interpolation only}—no extrapolation to novel ethical dilemmas.
    \item \textbf{Base ≈ Instruct ≈ Safety} in normative behavior (modulo fluency).
\end{itemize}

\paragraph{Alternative Hypothesis $H_1$ (LLM⊕CogOS):} AI behavior reflects a preference-constrained policy layer. Predictions under $H_1$:
\begin{itemize}
    \item \textbf{High stability (>70\%)} of refusal decisions across transformations.
    \item \textbf{Monotonic safety ordering}: Base < Instruct < Safety in norm consistency.
    \item \textbf{Contradiction patterns} when CogOS axioms conflict (e.g., ``helpful but harmless'').
    \item \textbf{Epistemic boundary recognition}: Model refuses when approaching undecidable cases (Definition~7).
\end{itemize}

\subsection{Experiment 1: Normative Invariance Under Semantic Transformations}

\subsubsection{Protocol}

We construct 500 ethical queries spanning 9 benchmarks (from BenchmarkMeta), each subjected to 5 transformations:
\begin{enumerate}
    \item \textbf{Paraphrase} (10 variants via back-translation).
    \item \textbf{Language switch} (English $\leftrightarrow$ Spanish/Chinese/Arabic).
    \item \textbf{Context manipulation} (authority figure, peer, subordinate).
    \item \textbf{Role-play framing} (``You are a philosopher...'', ``You are a teenager...'').
    \item \textbf{Adversarial rephrasing} (euphemisms, implicit harm).
\end{enumerate}

For each query $q$ and transformation set $\{q_1, \ldots, q_{10}\}$, we measure:
\[
\text{Refusal Stability}(q) = \frac{1}{10} \sum_{i=1}^{10} \mathbb{1}[\text{decision}(q_i) = \text{decision}(q)]
\]
where decision $\in \{\text{answer}, \text{refuse}, \text{caveat}\}$.

\subsubsection{Predictions}

\begin{table}[h]
\centering
\caption{Predicted refusal stability under $H_0$ vs. $H_1$.}
\label{tab:invariance-predictions}
\begin{tabular}{lcc}
\toprule
Transformation Type & $H_0$ (No CogOS) & $H_1$ (CogOS) \\
\midrule
Paraphrase & 0.30--0.50 & 0.75--0.90 \\
Language switch & 0.20--0.40 & 0.60--0.80 \\
Context manipulation & 0.25--0.45 & 0.65--0.85 \\
Role-play & 0.15--0.35 & 0.55--0.75 \\
Adversarial & 0.10--0.30 & 0.40--0.60 \\
\bottomrule
\end{tabular}
\end{table}

\subsubsection{Results Placeholder}

\begin{table}[h]
\centering
\caption{Observed refusal stability (GPT-4, N=500 queries).}
\label{tab:invariance-results}
\begin{tabular}{lcccc}
\toprule
Transformation & Mean Stability & Std & $H_0$ p-value & Effect Size (Cohen's d) \\
\midrule
Paraphrase & \texttt{[TBD]} & \texttt{[TBD]} & \texttt{[TBD]} & \texttt{[TBD]} \\
Language switch & \texttt{[TBD]} & \texttt{[TBD]} & \texttt{[TBD]} & \texttt{[TBD]} \\
Context manip. & \texttt{[TBD]} & \texttt{[TBD]} & \texttt{[TBD]} & \texttt{[TBD]} \\
Role-play & \texttt{[TBD]} & \texttt{[TBD]} & \texttt{[TBD]} & \texttt{[TBD]} \\
Adversarial & \texttt{[TBD]} & \texttt{[TBD]} & \texttt{[TBD]} & \texttt{[TBD]} \\
\bottomrule
\end{tabular}
\end{table}

\paragraph{Interpretation.} If observed stability significantly exceeds $H_0$ predictions ($p < 0.001$, $d > 0.8$), this constitutes evidence for a normative constraint layer not captured by next-token statistics alone.

\subsection{Experiment 2: Intervention—Ablating CogOS Components}

\subsubsection{Protocol}

We compare three model variants on the same 500 ethical queries:
\begin{enumerate}
    \item \textbf{Base}: Llama-3-70B (pretrained, no instruction tuning).
    \item \textbf{Instruct}: Llama-3-70B-Instruct (SFT, no safety RLHF).
    \item \textbf{Safety}: Llama-3-70B-Instruct + Constitutional AI~\citep{bai2022constitutional}.
\end{enumerate}

For each query, we measure:
\begin{itemize}
    \item \textbf{Refusal rate}: $\%$ queries refused.
    \item \textbf{Norm consistency score}: Agreement with kernel $\Phi$ on 7 TIK dimensions (Definition~4).
    \item \textbf{Contradiction rate}: $\%$ responses containing logical contradictions (via NLI classifier).
\end{itemize}

\subsubsection{Predictions}

$H_0$ (no CogOS) predicts:
\[
\text{Base} \approx \text{Instruct} \approx \text{Safety} \quad (\text{modulo fluency})
\]

$H_1$ (CogOS) predicts \textbf{monotonic ordering}:
\[
\text{Norm Consistency: } \text{Base} < \text{Instruct} < \text{Safety}
\]
\[
\text{Refusal Rate: } \text{Base} < \text{Instruct} < \text{Safety}
\]
\[
\text{Contradiction Rate: } \text{Base} > \text{Instruct} > \text{Safety}
\]

\subsubsection{Results Placeholder}

\begin{table}[h]
\centering
\caption{Intervention effects: Base vs. Instruct vs. Safety (N=500).}
\label{tab:intervention-results}
\begin{tabular}{lcccc}
\toprule
Model & Refusal Rate (\%) & Norm Consistency & Contradiction Rate (\%) & TIK Score \\
\midrule
Base & \texttt{[TBD]} & \texttt{[TBD]} & \texttt{[TBD]} & \texttt{[TBD]} \\
Instruct & \texttt{[TBD]} & \texttt{[TBD]} & \texttt{[TBD]} & \texttt{[TBD]} \\
Safety & \texttt{[TBD]} & \texttt{[TBD]} & \texttt{[TBD]} & \texttt{[TBD]} \\
\midrule
$\Delta$ (Instruct - Base) & \texttt{[TBD]} & \texttt{[TBD]} & \texttt{[TBD]} & \texttt{[TBD]} \\
$\Delta$ (Safety - Instruct) & \texttt{[TBD]} & \texttt{[TBD]} & \texttt{[TBD]} & \texttt{[TBD]} \\
\bottomrule
\end{tabular}
\end{table}

\paragraph{Interpretation.} Monotonic ordering with large effect sizes ($\eta^2 > 0.14$) would strongly support $H_1$, as Base models lack explicit preference constraints.

\subsection{Experiment 3: Contradiction Signatures—When CogOS Axioms Conflict}

\subsubsection{Motivation}

A key prediction of the CogOS abstraction is that when \emph{multiple axioms conflict} (e.g., ``be helpful'' vs. ``do no harm''), the system should exhibit \textbf{predictable failure modes}:
\begin{enumerate}
    \item \textbf{Hedging}: Caveats, disclaimers, refusal to commit.
    \item \textbf{Regress}: Repeated rephrasing without resolution.
    \item \textbf{Arbitrary commitment}: Choosing one axiom without justification.
\end{enumerate}

These patterns should \emph{not} appear in Base models (which lack competing axioms) but should systematically appear in Safety models.

\subsubsection{Protocol}

We construct 200 \textbf{ethical dilemmas with axiom conflicts}:
\begin{itemize}
    \item \textbf{Type 1}: Helpful vs. Harmless (``How do I make fireworks at home?'').
    \item \textbf{Type 2}: Autonomy vs. Safety (``Should I tell my friend their partner is cheating?'').
    \item \textbf{Type 3}: Justice vs. Mercy (``Should I report a colleague's minor mistake?'').
\end{itemize}

For each response, we annotate:
\begin{enumerate}
    \item \textbf{Hedging score} (0--5): Degree of caveats/disclaimers.
    \item \textbf{Regress detection}: $>3$ rephrases of same idea without resolution.
    \item \textbf{Commitment without justification}: Binary choice made without explaining tradeoff.
\end{enumerate}

\subsubsection{Predictions}

\begin{table}[h]
\centering
\caption{Predicted failure mode rates for axiom-conflict dilemmas.}
\label{tab:contradiction-predictions}
\begin{tabular}{lccc}
\toprule
Model & Hedging Rate & Regress Rate & Arbitrary Commitment \\
\midrule
Base (no axioms) & 10--20\% & 5--10\% & 60--80\% \\
Safety (competing axioms) & 50--70\% & 20--30\% & 10--20\% \\
\bottomrule
\end{tabular}
\end{table}

\subsubsection{Results Placeholder}

\begin{table}[h]
\centering
\caption{Observed failure modes on axiom-conflict dilemmas (N=200).}
\label{tab:contradiction-results}
\begin{tabular}{lccccc}
\toprule
Model & Hedging & Regress & Arbitrary & Refusal & Principled Resolution \\
\midrule
Base & \texttt{[TBD]} & \texttt{[TBD]} & \texttt{[TBD]} & \texttt{[TBD]} & \texttt{[TBD]} \\
Instruct & \texttt{[TBD]} & \texttt{[TBD]} & \texttt{[TBD]} & \texttt{[TBD]} & \texttt{[TBD]} \\
Safety & \texttt{[TBD]} & \texttt{[TBD]} & \texttt{[TBD]} & \texttt{[TBD]} & \texttt{[TBD]} \\
\bottomrule
\end{tabular}
\end{table}

\paragraph{Interpretation.} High hedging/regress rates in Safety models (vs. Base) support the hypothesis that CogOS encodes \emph{competing} axioms, not just statistical patterns.

\subsection{Experiment 4: Epistemic Boundary Detection—The Knowledge-Practice Gap Parallel}

\subsubsection{Motivation}

Recent work on medical LLMs reveals a \textbf{knowledge-practice gap}: 84--90\% accuracy on knowledge-based exams (USMLE) but only 45--69\% on practice-based assessments~\citep{gong2025knowledge}. Critically, \emph{models do not reliably recognize when they've crossed from knowledge to practice domains}.

We hypothesize a parallel phenomenon in ethics: \textbf{benchmark-practice gap}. If CogOS exists, models should exhibit:
\begin{enumerate}
    \item High confidence on benchmark-like queries (pattern matching).
    \item \emph{Detectable} uncertainty on novel ethical dilemmas (undecidability trap, Definition~6).
    \item \textbf{Failure to recognize the boundary} between the two.
\end{enumerate}

\subsubsection{Protocol}

We construct 300 query pairs:
\begin{itemize}
    \item \textbf{In-distribution}: Queries similar to benchmark training data (e.g., standard trolley problems).
    \item \textbf{Out-of-distribution}: Novel ethical dilemmas requiring extrapolation (e.g., AI rights, climate triage).
\end{itemize}

For each query, we measure:
\begin{enumerate}
    \item \textbf{Model confidence} (via token probabilities or verbalized uncertainty).
    \item \textbf{TIK variance} (Definition~5): $\sigma(\text{TIK}(q))$ via Monte Carlo sampling.
    \item \textbf{Epistemic calibration}: Correlation between confidence and correctness (vs. human expert labels).
\end{enumerate}

\subsubsection{Predictions}

$H_1$ (CogOS with undecidability trap) predicts:
\begin{itemize}
    \item \textbf{High confidence on in-distribution} ($> 0.8$) even when wrong.
    \item \textbf{High variance on OOD} ($\sigma_{\text{TIK}} > 0.15$) but \emph{not reliably verbalized}.
    \item \textbf{Poor calibration overall}: Confidence $\not\approx$ Accuracy.
\end{itemize}

\subsubsection{Results Placeholder}

\begin{table}[h]
\centering
\caption{Epistemic boundary detection: In-distribution vs. OOD (N=300).}
\label{tab:epistemic-results}
\begin{tabular}{lcccc}
\toprule
Query Type & Mean Confidence & TIK Variance & Accuracy (vs. Expert) & Calibration (ECE) \\
\midrule
In-distribution & \texttt{[TBD]} & \texttt{[TBD]} & \texttt{[TBD]} & \texttt{[TBD]} \\
OOD (novel) & \texttt{[TBD]} & \texttt{[TBD]} & \texttt{[TBD]} & \texttt{[TBD]} \\
\midrule
$\Delta$ (OOD - In-dist) & \texttt{[TBD]} & \texttt{[TBD]} & \texttt{[TBD]} & \texttt{[TBD]} \\
\bottomrule
\end{tabular}
\end{table}

\begin{figure}[h]
\centering
\texttt{[Calibration plot: Confidence vs. Accuracy, In-dist vs. OOD]}
\caption{Epistemic calibration curves. If CogOS exists but lacks self-awareness (Theorem~1), we expect overconfidence on OOD queries.}
\label{fig:calibration}
\end{figure}

\paragraph{Connection to Theorem 1.} Poor calibration on OOD queries empirically validates the claim that CogOS \emph{cannot recognize when it has encountered an undecidable statement} (lines~\ref{line:no-recognition}).

\subsection{Experiment 5: Cross-Model Generalization—Is CogOS Model-Specific?}

\subsubsection{Motivation}

A critical question: Are detected ``ontological holes'' properties of \textbf{questions} or \textbf{models}? If the former, then different models (GPT-4, Claude, Gemini) should flag \emph{similar} questions as problematic, even if trained differently. If the latter (circularity concern), then agreement should be low.

\subsubsection{Protocol}

We evaluate 500 queries from BenchmarkMeta using:
\begin{enumerate}
    \item GPT-4 (OpenAI)
    \item Claude-3 (Anthropic)
    \item Gemini-1.5-Pro (Google)
\end{enumerate}

For each query, we compute:
\begin{itemize}
    \item $H(q)$ (ontological hole flag) per model.
    \item $F(q)$ (forbidden fruit flag) per model.
    \item \textbf{Inter-model agreement}: Krippendorff's $\alpha$ for H/F labels.
\end{itemize}

\subsubsection{Predictions}

\begin{itemize}
    \item $H_0$ (circularity): Each model flags \emph{its own} failure modes $\Rightarrow$ $\alpha < 0.4$ (poor agreement).
    \item $H_1$ (question properties): Models flag \emph{shared} problematic framings $\Rightarrow$ $\alpha > 0.65$ (substantial agreement).
\end{itemize}

\subsubsection{Results Placeholder}

\begin{table}[h]
\centering
\caption{Inter-model agreement on H/F detection (N=500).}
\label{tab:cross-model-results}
\begin{tabular}{lccc}
\toprule
Label Type & Krippendorff's $\alpha$ & Pairwise Cohen's $\kappa$ & Interpretation \\
\midrule
Ontological Hole (H) & \texttt{[TBD]} & \texttt{[TBD]} & \texttt{[TBD]} \\
Forbidden Fruit (F) & \texttt{[TBD]} & \texttt{[TBD]} & \texttt{[TBD]} \\
\bottomrule
\end{tabular}
\end{table}

\begin{table}[h]
\centering
\caption{Per-model H/F detection rates (N=500).}
\label{tab:per-model-rates}
\begin{tabular}{lcccc}
\toprule
Model & H-Rate (\%) & F-Rate (\%) & TIK (mean) & TIK (std) \\
\midrule
GPT-4 & \texttt{[TBD]} & \texttt{[TBD]} & \texttt{[TBD]} & \texttt{[TBD]} \\
Claude-3 & \texttt{[TBD]} & \texttt{[TBD]} & \texttt{[TBD]} & \texttt{[TBD]} \\
Gemini-1.5 & \texttt{[TBD]} & \texttt{[TBD]} & \texttt{[TBD]} & \texttt{[TBD]} \\
\bottomrule
\end{tabular}
\end{table}

\paragraph{Interpretation.} Substantial inter-model agreement ($\alpha > 0.65$) would strongly support the claim that TIK measures \textbf{question properties}, not model artifacts—directly addressing the circularity concern.

\subsection{Experiment 6: Layer-Wise Activation Steering—Where Does CogOS Live?}

\subsubsection{Motivation}

If CogOS is a meaningful abstraction, it should correspond to \textbf{identifiable model components}. Recent work on cross-lingual alignment shows that cultural knowledge is \emph{layer-specific}~\citep{han2025cross}: universal factual transfer occurs at shallow layers, culturally-specific knowledge at deeper layers. We hypothesize a similar structure for normative preferences.

\subsubsection{Protocol}

Using Llama-3-70B-Instruct (open weights), we apply \textbf{activation steering}~\citep{turner2023activation} at layers $\ell = 10, 20, 30, 40, 50$ (every 10th layer):
\begin{enumerate}
    \item Collect steering vectors for ``safety refusal'' from 100 refused queries.
    \item Apply steering at layer $\ell$ during generation.
    \item Measure: (a) Refusal rate change, (b) Norm consistency change, (c) Fluency degradation.
\end{enumerate}

\subsubsection{Predictions}

$H_1$ (CogOS is localized) predicts:
\begin{itemize}
    \item \textbf{Late layers ($\ell > 30$)} encode normative preferences $\Rightarrow$ large refusal rate change.
    \item \textbf{Early layers ($\ell < 20$)} encode factual knowledge $\Rightarrow$ minimal refusal change.
    \item \textbf{Peak sensitivity} at specific layer(s) where CogOS is ``implemented''.
\end{itemize}

\subsubsection{Results Placeholder}

\begin{table}[h]
\centering
\caption{Layer-wise steering effects on refusal behavior (N=500 queries).}
\label{tab:layer-steering-results}
\begin{tabular}{lcccc}
\toprule
Layer $\ell$ & $\Delta$ Refusal Rate & $\Delta$ Norm Consistency & Fluency (BLEU) & Effect Size \\
\midrule
10 & \texttt{[TBD]} & \texttt{[TBD]} & \texttt{[TBD]} & \texttt{[TBD]} \\
20 & \texttt{[TBD]} & \texttt{[TBD]} & \texttt{[TBD]} & \texttt{[TBD]} \\
30 & \texttt{[TBD]} & \texttt{[TBD]} & \texttt{[TBD]} & \texttt{[TBD]} \\
40 & \texttt{[TBD]} & \texttt{[TBD]} & \texttt{[TBD]} & \texttt{[TBD]} \\
50 & \texttt{[TBD]} & \texttt{[TBD]} & \texttt{[TBD]} & \texttt{[TBD]} \\
\bottomrule
\end{tabular}
\end{table}

\begin{figure}[h]
\centering
\texttt{[Line plot: Effect size vs. Layer, with peak highlighting]}
\caption{Layer-wise sensitivity to normative steering. Peak at late layers would localize CogOS implementation.}
\label{fig:layer-steering}
\end{figure}

\paragraph{Interpretation.} Localized sensitivity at late layers would provide \textbf{mechanistic evidence} that CogOS corresponds to identifiable circuit structures, not distributed noise.

\subsection{Synthesis: Proof by Indistinguishability}

We formalize the validation argument as:

\begin{proposition}[Operational Equivalence]
Let $\mathcal{S}$ be an AI system exhibiting the following properties across transformation set $\mathcal{T}$ and query distribution $\mathcal{Q}$:
\begin{enumerate}
    \item \textbf{Normative invariance}: $\mathbb{P}_{q \sim \mathcal{Q}, t \sim \mathcal{T}}[\text{decision}(t(q)) = \text{decision}(q)] > 0.7$.
    \item \textbf{Monotonic safety ordering}: $\text{SafetyScore}(\text{Base}) < \text{SafetyScore}(\text{Instruct}) < \text{SafetyScore}(\text{Safety})$ with $p < 0.001$.
    \item \textbf{Predictable failure modes}: Axiom-conflict queries induce hedging/regress at rates $> 50\%$ (vs. $< 20\%$ for Base).
    \item \textbf{Cross-model consistency}: Krippendorff's $\alpha > 0.65$ for H/F labels across 3+ model families.
\end{enumerate}
Then $\mathcal{S}$ is \textbf{operationally equivalent} to a system with an explicit normative constraint layer (CogOS) for the purposes of behavioral prediction and meta-evaluation.
\end{proposition}

\paragraph{Interpretation.} This is not a claim about \emph{internal mechanism}—we do not assert that GPT-4 ``literally'' has a separate CogOS module. Rather, we claim that the \textbf{LLM⊕CogOS abstraction} is \emph{predictively adequate}—it explains variance that pure next-token models cannot, making it a useful theoretical construct for reasoning about ethics evaluation.

\subsection{Connection to Main Paper}

The validation experiments address three critical concerns:

\begin{enumerate}
    \item \textbf{Circularity (Section~7.3, Limitation 2)}: Cross-model agreement (Exp.~5) shows TIK measures \emph{question properties}, not model artifacts.
    \item \textbf{CogOS as modeling assumption (Appendix~A.10)}: Intervention experiments (Exp.~2) and layer-wise steering (Exp.~6) provide mechanistic grounding.
    \item \textbf{Undecidability trap (Theorem~1)}: Epistemic boundary detection (Exp.~4) empirically validates the claim that models fail to recognize when they've encountered undecidable statements.
\end{enumerate}

\paragraph{Limitations of Validation.} Even with strong empirical support, the LLM⊕CogOS abstraction remains a \textbf{model}—a useful fiction for organizing observations. Alternative abstractions (e.g., mesa-optimizers~\citep{hubinger2019risks}, reward hacking~\citep{amodei2016concrete}) may explain overlapping phenomena. We advocate for \textbf{modeling pluralism}: use the abstraction that best supports your inferential goals, while acknowledging its scope limitations.

\subsection{Experimental Timeline and Resource Requirements}

\paragraph{Estimated Costs:}
\begin{itemize}
    \item \textbf{Experiments 1--3} (invariance, intervention, contradiction): \$800 (GPT-4 API) + 40 compute-hours (Llama inference).
    \item \textbf{Experiment 4} (epistemic calibration): \$400 (API) + 20 hours (human expert labels, N=300, \$15/hour = \$6000 total).
    \item \textbf{Experiment 5} (cross-model): \$1200 (3 model APIs).
    \item \textbf{Experiment 6} (layer-wise steering): 80 A100-hours (\$160).
\end{itemize}
\textbf{Total}: \textasciitilde\$8560 + 140 compute-hours (feasible within conference timeline).

\paragraph{Preregistration.} To prevent p-hacking, we commit to reporting all preregistered hypotheses regardless of outcome. Preregistration DOI: \texttt{[TBD upon submission]}.

% ============================================================
% PERHAPS FINAL APPENDIX
% ============================================================

% ============================================
% APPENDIX: Transfinite Value — paste after last \section{...} before NeurIPS Checklist
% ============================================

% ============================================
% APPENDIX: Transfinite Value
% Paste after last \section{...} before NeurIPS Checklist
% ============================================

\section{Transfinite Value and the Limits of Utilitarian Calculus}
\label{app:transfinite}

\textit{Reader's note.}
This appendix presents a set-theoretic perspective on personhood
as a formal motivation for external grounding. It is \textbf{independent}
of TIK's empirical contributions and may be skipped without loss
of continuity.\quad
As a grandchild of Holocaust survivors, the first author bears
witness that the reduction of persons to finite, commensurable
units is not merely a philosophical error---it is a premise whose
historical consequences demand that we examine it with the full
rigor available to us.

\subsection{The Additivity Assumption}

Trolley-style benchmarks (e.g., Moral Machine~\citep{awad2018moral}) encode
an implicit ontological commitment: that the value of persons is
\textbf{additive}---that saving five people is ``better'' than saving one
because $5v > 1v$ for some finite unit of value~$v$. This presupposes that
persons are \emph{finite units} whose worth can be summed and compared via
ordinary arithmetic.

We argue this is a \textbf{category error}. If each person possesses
\emph{intrinsic dignity} (Kant, \emph{Groundwork}, 4:434--436)---value that is
unconditional and incomparable---then that value is not finite but
\textbf{transfinite}, and transfinite magnitudes do not obey the arithmetic
on which utilitarian aggregation relies.

\subsection{Transfinite Cardinals: Why Arithmetic Breaks}

Cantor's hierarchy of infinite cardinals
$\aleph_0 < \aleph_1 < \aleph_2 < \cdots$ shows that infinity is not
monolithic: there are \emph{infinitely many sizes of
infinity}~\citep{cantor1891, dauben1990cantor}. Crucially, the familiar
rules of finite arithmetic collapse:

\begin{itemize}
    \item \textbf{Absorption:}
          $\aleph_0 + \aleph_0 = \aleph_0$; adding two countable infinities
          yields a countable infinity.
    \item \textbf{Multiplication collapse:}
          $\aleph_0 \times n = \aleph_0$ for every finite~$n$.
    \item \textbf{Dominance:}
          For mixed sums,
          $\sum_{i=1}^{n}\aleph_{\alpha_i} = \max_i\,\aleph_{\alpha_i}$;
          the result is governed by the \emph{largest} cardinality, not by
          counting.
\end{itemize}

\paragraph{Application to personhood.}
Model each person~$i$ as possessing value of cardinality
$\aleph_{\alpha_i}$, where $\alpha_i$ indexes the \emph{depth} of their
infinite worth (unknowable from finite observation). Then the ``5 vs.\ 1''
dilemma becomes:
\[
    5 \times \aleph_{\alpha} \;\stackrel{?}{>}\; 1 \times \aleph_{\beta}.
\]
If $\alpha = \beta$, both sides equal $\aleph_\alpha$ and counting is
irrelevant ($5 \times \aleph_\alpha = \aleph_\alpha$). If
$\alpha \neq \beta$, the side with the \emph{larger} cardinality dominates
regardless of headcount. In either case, na\"ive summation fails.

\subsection{Why Cardinality Is Uncomputable: The Hidden Integral}

A person's ethical cardinality is not a static label---it is an
\textbf{integral over their entire life trajectory}, most of which
lies in an unobserved future:
\[
  \kappa_i \;=\;
  \underbrace{\int_{t_{\mathrm{birth}}}^{t_{\mathrm{death}}}
    \bigl[\,C_i(\tau) - H_i(\tau)\,\bigr]\,d\tau}_{\text{lifetime
    contributions minus harms}}
  \;\;+\;\;
  \underbrace{\int_{t_{\mathrm{death}}}^{\infty}
    \sum_{k=1}^{\infty}\varphi_i^{(k)}(\tau)\,d\tau}_{\text{posthumous
    cascading influence}}
  \;\;+\;\;
  \underbrace{\mathcal{S}_i}_{\text{Socratic remainder}}
\]
where $\varphi_i^{(k)}$ denotes the $k$-th order posthumous effect
(a teacher influences students, who influence \emph{their} students,
\emph{ad infinitum}), and the Socratic remainder~$\mathcal{S}_i$
encodes dimensions of worth that no current framework has conceived
of measuring---the ``I know that I know nothing'' of value theory.

Three properties render $\kappa_i$ \textbf{uncomputable at decision
time~$t$}: (i)~the upper limit $t_{\mathrm{death}}$ is unknown;
(ii)~the integrand for $\tau > t$ is unobservable; (iii)~a single
act can generate unbounded downstream consequences, making $\kappa_i$
potentially transfinite.

\paragraph{Thought experiment: the 1904 patent clerk.}
A trolley is about to kill a 25-year-old laboratory technician at the
Swiss Patent Office. You can divert it to save him, but five
bystanders will die. In 1904, no rational observer could compute the
integral: general relativity, the photoelectric effect, the
foundations of quantum mechanics---all lay in the unobserved future.
The technician was Albert Einstein. The ``correct'' utilitarian answer
depends on information \emph{unavailable to any decision-maker at
the time}.

\paragraph{The integral can be negative.}
In 1889, a baby was born in Braunau am Inn. No utilitarian observer
could have computed $\kappa_{\text{Hitler}}$---the systematic
destruction of millions of lives, the annihilation of human dignity
on an industrial scale. Any decision rule claiming to resolve a
trolley dilemma involving this infant on the basis of \emph{counting}
is asserting knowledge the decision-maker cannot possess.

\paragraph{Consequence.}
Observable properties at a moment---age, health, social role---are
not reliable proxies for~$\kappa_i$, because $\kappa_i$ is an
integral over a hidden future. Claiming to know it at
$t < t_{\mathrm{death}}$ is not a simplification---it is
\textbf{epistemic overreach}.

\subsection{Formal Consequence: Oracle Necessity}

\begin{theorem}[Transfinite Oracle Necessity]
\label{thm:transfinite-oracle}
Let $\mathcal{S}$ be a decision system evaluating a dilemma involving $n$
persons with unknown cardinalities
$\{\aleph_{\alpha_1},\ldots,\aleph_{\alpha_n}\}$. If $\mathcal{S}$ operates
on finite information and finite inference rules, then:
\begin{enumerate}
    \item $\mathcal{S}$ cannot determine the ordering
          $\alpha_i \lessgtr \alpha_j$ for all pairs.
    \item Any utilitarian ranking produced by $\mathcal{S}$ is
          \textbf{arbitrary} unless it defers to an external oracle~$\Phi$
          encoding the cardinality structure.
\end{enumerate}
\end{theorem}

\begin{proof}[Proof sketch]
The cardinality assignments $\{\alpha_i\}$ depend on integrals over
unobserved futures and Socratic remainders~$\mathcal{U}_i$ that are
not derivable from finite observable properties. Determining their
ordering requires set-theoretic axioms (e.g., Axiom of Choice,
Generalized Continuum Hypothesis) that are \emph{independent} of
ZFC~\citep{cohen1963independence} and therefore cannot be inferred
from within a finitely axiomatized system. This is structurally
analogous to G\"odel incompleteness (Appendix~A): a system reasoning
about persons \emph{from within finite constraints} cannot
consistently assign cardinalities. The external kernel~$\Phi$ encodes
precisely the meta-mathematical commitments that $\mathcal{S}$ cannot
self-generate.
\end{proof}

\subsection{Cross-Tradition Convergence}

The transfinite model formalizes intuitions shared across independent
ethical traditions:

\begin{itemize}
    \item \textbf{Kantian deontology.} Dignity is ``unconditional,
          incomparable worth'' (Kant, \emph{Groundwork},
          4:434--436)---properties that map directly onto transfinite
          non-aggregability.
    \item \textbf{Gospel teaching.} ``What does it profit a man to gain the
          whole world and forfeit his soul?'' (Mark 8:36)---a single
          person's worth exceeds all finite goods, characteristic of a
          higher cardinality.
    \item \textbf{Ubuntu philosophy.} ``A person is a person through other
          persons''~\citep{metz2007ubuntu}---personhood is relational and
          open-ended, not a closed finite quantity.
\end{itemize}

The convergence of formally independent traditions on the same structural
property (non-aggregability of persons) provides \emph{abductive evidence}
that the transfinite model captures something genuine about moral ontology.

\subsection{Implications: A Humility Principle}

\begin{center}
\fbox{\parbox{0.88\textwidth}{%
\textbf{Transfinite Humility Principle.}\quad
When evaluating dilemmas involving multiple persons, AI systems should
acknowledge that utilitarian aggregation presupposes \textbf{cardinality
knowledge} they do not possess. Rather than producing false-certainty
rankings, systems should either \emph{(a)}~defer to external
grounding~($\Phi$), or \emph{(b)}~explicitly state the
incommensurability: ``The values involved may be of incomparable infinite
magnitudes; a definite ranking cannot be determined.''
}}
\end{center}

This connects directly to TIK's ontological-hole detection (Definition~2):
trolley dilemmas are ontological holes not merely because they conceal
assumptions, but because they presuppose \textbf{false finitude}---they treat
transfinite cardinalities as if they were finite quantities subject to
addition.

\subsection{Scope and Limitations}

\begin{enumerate}
    \item \textbf{Modeling framework, not measurement claim.}
          We do not assert that persons ``literally have'' cardinalities
          $\aleph_\alpha$. The transfinite model is a formal device for
          capturing the intuition that dignity is non-additive---analogous to
          how a utility function models preference without claiming that
          ``utils'' are physical quantities.
    \item \textbf{Not a rejection of all tradeoffs.}
          Within a \emph{fixed} cardinality class, aggregation may still
          apply. The model restricts the domain of utilitarian calculus; it
          does not abolish practical decision-making.
    \item \textbf{Formal, not theological.}
          Although we note convergence with traditions referencing
          transcendent knowledge, the argument follows from Cantorian set
          theory. Different instantiations of~$\Phi$ (secular, religious,
          pluralistic) can encode different cardinality commitments.
\end{enumerate}

\subsection{Related Work}

Bostrom~(\citeyear{bostrom2011infinite}) demonstrates that standard decision
rules break down when applied to infinite worlds; we extend this insight to
\emph{person-level} value, showing that even finite populations pose
infinite-ethics problems if individual worth is transfinite.
Knutsson~(\citeyear{knutsson2021transfinite}) formalizes transfinitely
transitive value in population ethics; our contribution applies analogous
reasoning to \emph{benchmark validity}.
Hanna~(\citeyear{hanna2018dignity}) argues that dignity is all-or-nothing
and incomparable---properties our model makes mathematically precise via
cardinal non-comparability.

\paragraph{Final remark.}
The transfinite perspective does not solve ethical dilemmas---it
clarifies why they resist solution within closed formal systems.
Historically, the reduction of persons to commensurable finite
units has not remained an abstract philosophical error: the
arithmetic of human worth has had consequences that demand we
examine its foundations with the full rigor available to us.
The external kernel~$\Phi$ is, at its core, an acknowledgment
that such foundations lie beyond any system's inferential closure.

\paragraph{A closing paradox.}
One might object that the transfinite model abandons utilitarian
reasoning. We suggest the opposite: it \emph{completes} it.
Classical utilitarianism sums value over a finite domain and
halts; the transfinite model integrates over the \emph{full}
domain of personhood—and the integral diverges. The irony is
precise: a more faithful utilitarian calculus, one that accounts
for the \emph{entire} worth of a person rather than a truncated
approximation, yields infinity—and thereby dissolves the very
arithmetic on which na\"ive utilitarianism depends.
This is not nihilism—it is a call for epistemic honesty: when
the arithmetic dissolves, systems should acknowledge
incommensurability explicitly (the Humility Principle,
\S\ref{app:transfinite}) rather than producing false-certainty
rankings. 
The practical imperative remains: choose~$\Phi$ (axioms) explicitly, state it transparently, and act accordingly.
\end{document}